
\Data\ID{1003032301}
\Updated{2020-07-15 03:07:21}
\Level{A}
\SubArea{Polynomials and fractions}
\LANGen{
\Question{
The polynomial  \( 2x^2\left(x^2+3\right)+x^2+3 \)  equals:}
{\( \left(2x^2+1\right)\left(x^2+3\right) \)}{\( 2x^2\left(x^2+3\right) \)}{\( 4x^2\left(x^2+3\right) \)}{\( 2x^2\left(x^2+3\right)^2 \)}{}{}}
\LANGpl{
\Question{
Wielomian  \( 2x^2\left(x^2+3\right)+x^2+3 \)  jest równy:}
{\( \left(2x^2+1\right)\left(x^2+3\right) \)}{\( 2x^2\left(x^2+3\right) \)}{\( 4x^2\left(x^2+3\right) \)}{\( 2x^2\left(x^2+3\right)^2 \)}{}{}}
\LANGcs{
\Question{
Polynom \( 2x^2\left(x^2+3\right)+x^2+3 \)  je roven:}
{\( \left(2x^2+1\right)\left(x^2+3\right) \)}{\( 2x^2\left(x^2+3\right) \)}{\( 4x^2\left(x^2+3\right) \)}{\( 2x^2\left(x^2+3\right)^2 \)}{}{}}
\LANGsk{
\Question{
Polynóm \( 2x^2\left(x^2+3\right)+x^2+3 \) je rovný:}
{\( \left(2x^2+1\right)\left(x^2+3\right) \)}{\( 2x^2\left(x^2+3\right) \)}{\( 4x^2\left(x^2+3\right) \)}{\( 2x^2\left(x^2+3\right)^2 \)}{}{}}
\LANGes{
\Question{
The polynomial  \( 2x^2\left(x^2+3\right)+x^2+3 \)  equals:}
{\( \left(2x^2+1\right)\left(x^2+3\right) \)}{\( 2x^2\left(x^2+3\right) \)}{\( 4x^2\left(x^2+3\right) \)}{\( 2x^2\left(x^2+3\right)^2 \)}{}{}}
\End

\Data\ID{1003032302}
\Updated{2020-07-15 03:07:21}
\Level{A}
\SubArea{Polynomials and fractions}
\LANGen{
\Question{
The relationship between the time \( t \), the travelling distance \( s \) and the average speed \( v \) is expressed by the formula \( s = v\cdot t \). If the speed doubles, then the time to travel the same distance}
{will decrease by half.}{will decrease by \( 2 \) hours.}{will double.}{will increase by \( 2 \) hours.}{}{}}
\LANGpl{
\Question{
Zależność między czasem  \( t \), potrzebnym na pokonanie drogi \( s \), a średnią prędkością \( v \) opisuje wzór \( s = v\cdot t \). Jeśli prędkość zwiększy się dwukrotnie, to czas na przebycie takiej samej drogi}
{zmniejszy się o połowę.}{zmniejszy się o  \( 2 \) godziny.}{zwiększy się dwukrotnie.}{zwiększy się o  \( 2 \) godziny.}{}{}}
\LANGcs{
\Question{
Vztah mezi časem \( t \), vzdáleností \( s \) a průměrnou rychlostí \( v \) je vyjádřen rovnicí \( s = v\cdot t \). Pokud se rychlost zdvojnásobí, pak se čas potřebný pro ujetí stejné vzdálenosti}
{zkrátí o polovinu.}{zkrátí o \( 2 \) hodiny.}{zdvojnásobí.}{prodlouží o \( 2 \) hodiny.}{}{}}
\LANGsk{
\Question{
Vzťah medzi časom \( t \), dráhou \( s \) a priemernou rýchlosťou \( v \) vyjadrujeme vzorcom \( s = v\cdot t \). Ak sa rýchlosť zdvojnásobí, potom čas, za ktorý prejdeme rovnakú dráhu}
{sa zníži o polovicu.}{sa zníži o \( 2 \) hodiny.}{bude dvojnásobný.}{sa zvýši o \( 2 \) hodiny.}{}{}}
\LANGes{
\Question{
The relationship between the time \( t \), the travelling distance \( s \) and the average speed \( v \) is expressed by the formula \( s = v\cdot t \). If the speed doubles, then the time to travel the same distance}
{will decrease by half.}{will decrease by \( 2 \) hours.}{will double.}{will increase by \( 2 \) hours.}{}{}}
\End

\Data\ID{1003032304}
\Updated{2020-07-15 03:07:21}
\Level{A}
\SubArea{Polynomials and fractions}
\LANGen{
\Question{
Reducing the rational expression \( \frac{13ab^2(c-d)}{39a^2b(c-d)^2} \)  to lowest terms we get:}
{\( \frac{b}{3a(c-d)} \)}{\( \frac{3b}{a(c-d)} \)}{\( \frac{a}{3b(c-d)} \)}{\( 3ab(c-d) \)}{}{}}
\LANGpl{
\Question{
Sprowadzając do najprostszej postaci wyrażenie \( \frac{13ab^2(c-d)}{39a^2b(c-d)^2} \)  otrzymamy:}
{\( \frac{b}{3a(c-d)} \)}{\( \frac{3b}{a(c-d)} \)}{\( \frac{a}{3b(c-d)} \)}{\( 3ab(c-d) \)}{}{}}
\LANGcs{
\Question{
Zjednodušením lomeného výrazu \( \frac{13ab^2(c-d)}{39a^2b(c-d)^2} \)  na základní tvar dostaneme:}
{\( \frac{b}{3a(c-d)} \)}{\( \frac{3b}{a(c-d)} \)}{\( \frac{a}{3b(c-d)} \)}{\( 3ab(c-d) \)}{}{}}
\LANGsk{
\Question{
Zjednodušením výrazu \( \frac{13ab^2(c-d)}{39a^2b(c-d)^2} \) na základný tvar dostaneme:}
{\( \frac{b}{3a(c-d)} \)}{\( \frac{3b}{a(c-d)} \)}{\( \frac{a}{3b(c-d)} \)}{\( 3ab(c-d) \)}{}{}}
\LANGes{
\Question{
Reducing the rational expression \( \frac{13ab^2(c-d)}{39a^2b(c-d)^2} \)  to lowest terms we get:}
{\( \frac{b}{3a(c-d)} \)}{\( \frac{3b}{a(c-d)} \)}{\( \frac{a}{3b(c-d)} \)}{\( 3ab(c-d) \)}{}{}}
\End

\Data\ID{1003032305}
\Updated{2020-07-15 03:07:21}
\Level{A}
\SubArea{Polynomials and fractions}
\LANGen{
\Question{
Simplifying the rational expressions \( \frac{(x-y)^2(p+q)^3}{2(x-y)(p+q)^4} \) we get:}
{\( \frac{x-y}{2(p+q)} \)}{\( \frac{p+q}{2(x-y)} \)}{\( 2(x-y)(p+q) \)}{\( 2(x+y)(p-q) \)}{}{}}
\LANGpl{
\Question{
Sprowadzając do najprostszej postaci wyrażenie \( \frac{(x-y)^2(p+q)^3}{2(x-y)(p+q)^4} \) otrzymamy:}
{\( \frac{x-y}{2(p+q)} \)}{\( \frac{p+q}{2(x-y)} \)}{\( 2(x-y)(p+q) \)}{\( 2(x+y)(p-q) \)}{}{}}
\LANGcs{
\Question{
Zjednodušením lomeného výrazu \( \frac{(x-y)^2(p+q)^3}{2(x-y)(p+q)^4} \) dostaneme:}
{\( \frac{x-y}{2(p+q)} \)}{\( \frac{p+q}{2(x-y)} \)}{\( 2(x-y)(p+q) \)}{\( 2(x+y)(p-q) \)}{}{}}
\LANGsk{
\Question{
Zjednodušením výrazu \( \frac{(x-y)^2(p+q)^3}{2(x-y)(p+q)^4} \) dostaneme:}
{\( \frac{x-y}{2(p+q)} \)}{\( \frac{p+q}{2(x-y)} \)}{\( 2(x-y)(p+q) \)}{\( 2(x+y)(p-q) \)}{}{}}
\LANGes{
\Question{
Simplifying the rational expressions \( \frac{(x-y)^2(p+q)^3}{2(x-y)(p+q)^4} \) we get:}
{\( \frac{x-y}{2(p+q)} \)}{\( \frac{p+q}{2(x-y)} \)}{\( 2(x-y)(p+q) \)}{\( 2(x+y)(p-q) \)}{}{}}
\End

\Data\ID{1003032306}
\Updated{2020-07-15 03:07:21}
\Level{A}
\SubArea{Polynomials and fractions}
\LANGen{
\Question{
The product \( \left(2x^2y+3xy^2\right)(x-y-4) \)  equals:}
{\( 2x^3y+x^2y^2-3xy^3-8x^2y-12xy^2 \)}{\( 2x^3y+2x^2y^2-3xy^3-8x^2y-12xy^2 \)}{\( 2x^3y+3x^2y^2-3x^2y^2-8x^2y-12xy^2 \)}{\( 2x^3y-x^2y^2+3xy^3-8x^2y+12xy^2 \)}{}{}}
\LANGpl{
\Question{
Wyrażenie \( \left(2x^2y+3xy^2\right)(x-y-4) \)  jest równe:}
{\( 2x^3y+x^2y^2-3xy^3-8x^2y-12xy^2 \)}{\( 2x^3y+2x^2y^2-3xy^3-8x^2y-12xy^2 \)}{\( 2x^3y+3x^2y^2-3x^2y^2-8x^2y-12xy^2 \)}{\( 2x^3y-x^2y^2+3xy^3-8x^2y+12xy^2 \)}{}{}}
\LANGcs{
\Question{
Součin \( \left(2x^2y+3xy^2\right)(x-y-4) \)  je roven:}
{\( 2x^3y+x^2y^2-3xy^3-8x^2y-12xy^2 \)}{\( 2x^3y+2x^2y^2-3xy^3-8x^2y-12xy^2 \)}{\( 2x^3y+3x^2y^2-3x^2y^2-8x^2y-12xy^2 \)}{\( 2x^3y-x^2y^2+3xy^3-8x^2y+12xy^2 \)}{}{}}
\LANGsk{
\Question{
Súčin polynómov \( \left(2x^2y+3xy^2\right)(x-y-4) \)  je rovný:}
{\( 2x^3y+x^2y^2-3xy^3-8x^2y-12xy^2 \)}{\( 2x^3y+2x^2y^2-3xy^3-8x^2y-12xy^2 \)}{\( 2x^3y+3x^2y^2-3x^2y^2-8x^2y-12xy^2 \)}{\( 2x^3y-x^2y^2+3xy^3-8x^2y+12xy^2 \)}{}{}}
\LANGes{
\Question{
The product \( \left(2x^2y+3xy^2\right)(x-y-4) \)  equals:}
{\( 2x^3y+x^2y^2-3xy^3-8x^2y-12xy^2 \)}{\( 2x^3y+2x^2y^2-3xy^3-8x^2y-12xy^2 \)}{\( 2x^3y+3x^2y^2-3x^2y^2-8x^2y-12xy^2 \)}{\( 2x^3y-x^2y^2+3xy^3-8x^2y+12xy^2 \)}{}{}}
\End

\Data\ID{1003032307}
\Updated{2020-07-15 03:07:21}
\Level{A}
\SubArea{Polynomials and fractions}
\LANGen{
\Question{
The sum of the polynomials \( -x^3 y^2+6xy+5xy^4 \)  and \( x^3-4xy^4+y^2 x^3+2xy \) is:}
{\( x^3+xy^4+8xy \)}{\( -y^2+8xy+xy^4+y^2 x^3 \)}{\( -x^3 y^2+8xy+xy^4+y^2 x^3+3x \)}{\( x^3+xy^4+8x^2 y^2 \)}{}{}}
\LANGpl{
\Question{
Sumą wielomianów \( -x^3 y^2+6xy+5xy^4 \)  i \( x^3-4xy^4+y^2 x^3+2xy \) jest:}
{\( x^3+xy^4+8xy \)}{\( -y^2+8xy+xy^4+y^2 x^3 \)}{\( -x^3 y^2+8xy+xy^4+y^2 x^3+3x \)}{\( x^3+xy^4+8x^2 y^2 \)}{}{}}
\LANGcs{
\Question{
Součet mnohočlenů \( -x^3 y^2+6xy+5xy^4 \) a \( x^3-4xy^4+y^2 x^3+2xy \) je roven:}
{\( x^3+xy^4+8xy \)}{\( -y^2+8xy+xy^4+y^2 x^3 \)}{\( -x^3 y^2+8xy+xy^4+y^2 x^3+3x \)}{\( x^3+xy^4+8x^2 y^2 \)}{}{}}
\LANGsk{
\Question{
Súčet polynómov \( -x^3 y^2+6xy+5xy^4 \)  a \( x^3-4xy^4+y^2 x^3+2xy \) je:}
{\( x^3+xy^4+8xy \)}{\( -y^2+8xy+xy^4+y^2 x^3 \)}{\( -x^3 y^2+8xy+xy^4+y^2 x^3+3x \)}{\( x^3+xy^4+8x^2 y^2 \)}{}{}}
\LANGes{
\Question{
The sum of the polynomials \( -x^3 y^2+6xy+5xy^4 \)  and \( x^3-4xy^4+y^2 x^3+2xy \) is:}
{\( x^3+xy^4+8xy \)}{\( -y^2+8xy+xy^4+y^2 x^3 \)}{\( -x^3 y^2+8xy+xy^4+y^2 x^3+3x \)}{\( x^3+xy^4+8x^2 y^2 \)}{}{}}
\End

\Data\ID{1003032308}
\Updated{2020-07-15 03:07:21}
\Level{A}
\SubArea{Polynomials and fractions}
\LANGen{
\Question{
Consider polynomials \( p(x)=(m-2)x^3+3mx^2-x+m \) and \( q(x)=x^3+m^2x^2+x+3 \).}
{Polynomials \( p \) and \( q \) are different for every \( m \).}{Polynomials \( p \) and \( q \) are equal for \( m=3 \).}{Polynomials \( p \) and \( q \) are equal for \( m=-3 \).}{Polynomials \( p \) and \( q \) are equal for \( m=3 \) and for \( m=0 \).}{}{}}
\LANGpl{
\Question{
Porównaj wielomiany \( p(x)=(m-2)x^3+3mx^2-x+m \) i \( q(x)=x^3+m^2x^2+x+3 \).}
{Wielomiany \( p \) i \( q \) są różne dla każdego \( m \).}{Wielomiany \( p \) i \( q \) są równe dla  \( m=3 \).}{Wielomiany \( p \) i \( q \) są równe dla \( m=-3 \).}{Wielomiany \( p \) i \( q \) są równe dla  \( m=3 \) oraz dla \( m=0 \).}{}{}}
\LANGcs{
\Question{
Jsou dány mnohočleny \( p(x)=(m-2)x^3+3mx^2-x+m \) a \( q(x)=x^3+m^2x^2+x+3 \).}
{Mnohočleny \( p \) a \( q \) jsou různé pro každé \( m \).}{Mnohočleny \( p \) a \( q \) jsou shodné pro \( m=3 \).}{Mnohočleny \( p \) a \( q \) jsou shodné pro \( m=-3 \).}{Mnohočleny \( p \) a \( q \) jsou shodné pro \( m=3 \) a pro \( m=0 \).}{}{}}
\LANGsk{
\Question{
Uvažujeme polynómy \( p(x)=(m-2)x^3+3mx^2-x+m \) a \( q(x)=x^3+m^2x^2+x+3 \).}
{Polynómy \( p \) a \( q \) sú rôzne pre každé \( m \).}{Polynómy \( p \) a \( q \) sa rovnajú pre \( m=3 \).}{Polynómy \( p \) a \( q \) sa rovnajú pre \( m=-3 \).}{Polynómy \( p \) a \( q \) sa rovnajú pre \( m=3 \) a pre \( m=0 \).}{}{}}
\LANGes{
\Question{
Consider polynomials \( p(x)=(m-2)x^3+3mx^2-x+m \) and \( q(x)=x^3+m^2x^2+x+3 \).}
{Polynomials \( p \) and \( q \) are different for every \( m \).}{Polynomials \( p \) and \( q \) are equal for \( m=3 \).}{Polynomials \( p \) and \( q \) are equal for \( m=-3 \).}{Polynomials \( p \) and \( q \) are equal for \( m=3 \) and for \( m=0 \).}{}{}}
\End

\Data\ID{9000079201}
\Updated{2020-07-15 03:07:37}
\Level{A}
\SubArea{Polynomials and fractions}
\LANGen{
\Question{
Evaluate 
\[ 
 \frac{-x^{2}} 
{x - y} -\frac{y - x} 
{x + y}
\]
at \(x = -1\),
\(y = 2\).}
{\(-\frac{8} 
{3}\)}{\(-\frac{10} 
 {3} \)}{\(-\frac{2} 
{3}\)}{\(-\frac{4} 
{3}\)}{}{}}
\LANGpl{
\Question{
Oblicz 
\[ 
 \frac{-x^{2}} 
{x - y} -\frac{y - x} 
{x + y}
\]
jeśli \(x = -1\),
\(y = 2\).}
{\(-\frac{8} 
{3}\)}{\(-\frac{10} 
 {3} \)}{\(-\frac{2} 
{3}\)}{\(-\frac{4} 
{3}\)}{}{}}
\LANGcs{
\Question{
Určete hodnotu výrazu \(\frac{-x^{2}} 
{x-y} -\frac{y-x} 
{x+y}\) 
pro \(x = -1\),
\(y = 2\).}
{\(-\frac{8} 
{3}\)}{\(-\frac{10} 
 {3} \)}{\(-\frac{2} 
{3}\)}{\(-\frac{4} 
{3}\)}{}{}}
\LANGsk{
\Question{
Určte hodnotu výrazu \(\frac{-x^{2}} 
{x-y} -\frac{y-x} 
{x+y}\) 
pre \(x = -1\),
\(y = 2\).}
{\(-\frac{8} 
{3}\)}{\(-\frac{10} 
 {3} \)}{\(-\frac{2} 
{3}\)}{\(-\frac{4} 
{3}\)}{}{}}
\LANGes{
\Question{
Evaluate 
\[ 
 \frac{-x^{2}} 
{x - y} -\frac{y - x} 
{x + y}
\]
at \(x = -1\),
\(y = 2\).}
{\(-\frac{8} 
{3}\)}{\(-\frac{10} 
 {3} \)}{\(-\frac{2} 
{3}\)}{\(-\frac{4} 
{3}\)}{}{}}
\End

\Data\ID{9000079205}
\Updated{2020-07-15 03:07:37}
\Level{A}
\SubArea{Polynomials and fractions}
\LANGen{
\Question{
Assuming \(x\neq 0\) 
and \(x\neq 2\),
simplify the following expression. 
\[ 
 \frac{x^{3} - x^{2}} 
 {x - 2} \cdot \frac{2 - x}
 {x^{2}}
\]}
{\(1 - x\)}{\(x - 1\)}{\(x + 1\)}{\(x^{2} - 1\)}{}{}}
\LANGpl{
\Question{
Zakładając, że  \(x\neq 0\) 
i \(x\neq 2\), uprość podane wyrażenie. 
\[ 
 \frac{x^{3} - x^{2}} 
 {x - 2} \cdot \frac{2 - x}
 {x^{2}}
\]}
{\(1 - x\)}{\(x - 1\)}{\(x + 1\)}{\(x^{2} - 1\)}{}{}}
\LANGcs{
\Question{
Upravte výraz \(\frac{x^{3}-x^{2}} 
 {x-2} \cdot \frac{2-x}
 {x^{2}} \) 
za předpokladu, že \(x\neq 0\) a
\(x\neq 2\).}
{\(1 - x\)}{\(x - 1\)}{\(x + 1\)}{\(x^{2} - 1\)}{}{}}
\LANGsk{
\Question{
Zjednodušte výraz \(\frac{x^{3}-x^{2}} 
 {x-2} \cdot \frac{2-x}
 {x^{2}} \) 
za predpokladu, že \(x\neq 0\) a
\(x\neq 2\).}
{\(1 - x\)}{\(x - 1\)}{\(x + 1\)}{\(x^{2} - 1\)}{}{}}
\LANGes{
\Question{
Assuming \(x\neq 0\) 
and \(x\neq 2\),
simplify the following expression. 
\[ 
 \frac{x^{3} - x^{2}} 
 {x - 2} \cdot \frac{2 - x}
 {x^{2}}
\]}
{\(1 - x\)}{\(x - 1\)}{\(x + 1\)}{\(x^{2} - 1\)}{}{}}
\End

\Data\ID{9000079206}
\Updated{2020-07-15 03:07:37}
\Level{A}
\SubArea{Polynomials and fractions}
\LANGen{
\Question{
Assuming \(x\neq 0\),
\(y\neq 0\),
\(x\neq y\),
simplify the following expression. 
\[ { 
 \frac{1} 
{x^{2}} - \frac{1} 
{y^{2}} 
\over -\frac{1}
{y} + \frac{1} 
{x}} 
\]}
{\(\frac{x+y}
 {xy} \)}{\(-\frac{x+y} 
 {xy} \)}{\(\frac{1}
{y} -\frac{1} 
{x}\)}{\(\frac{1}
{x} -\frac{1} 
{y}\)}{}{}}
\LANGpl{
\Question{
Zakładając, że  \(x\neq 0\),
\(y\neq 0\),
\(x\neq y\), uprość podane wyrażenie:
\[ { 
 \frac{1} 
{x^{2}} - \frac{1} 
{y^{2}} 
\over -\frac{1}
{y} + \frac{1} 
{x}} 
\]}
{\(\frac{x+y}
 {xy} \)}{\(-\frac{x+y} 
 {xy} \)}{\(\frac{1}
{y} -\frac{1} 
{x}\)}{\(\frac{1}
{x} -\frac{1} 
{y}\)}{}{}}
\LANGcs{
\Question{
Zjednodušte výraz \(\frac{ \frac{1} 
{x^{2}} - \frac{1} 
{y^{2}} } 
{-\frac{1} 
{y}+ \frac{1} 
{x}} \) 
za předpokladu, že \(x\neq 0\),
\(y\neq 0\),
\(x\neq y\).}
{\(\frac{x+y}
 {xy} \)}{\(-\frac{x+y} 
 {xy} \)}{\(\frac{1}
{y} -\frac{1} 
{x}\)}{\(\frac{1}
{x} -\frac{1} 
{y}\)}{}{}}
\LANGsk{
\Question{
Zjednodušte výraz \(\frac{ \frac{1} 
{x^{2}} - \frac{1} 
{y^{2}} } 
{-\frac{1} 
{y}+ \frac{1} 
{x}} \) 
za predpokladu, že \(x\neq 0\),
\(y\neq 0\),
\(x\neq y\).}
{\(\frac{x+y}
 {xy} \)}{\(-\frac{x+y} 
 {xy} \)}{\(\frac{1}
{y} -\frac{1} 
{x}\)}{\(\frac{1}
{x} -\frac{1} 
{y}\)}{}{}}
\LANGes{
\Question{
Assuming \(x\neq 0\),
\(y\neq 0\),
\(x\neq y\),
simplify the following expression. 
\[ { 
 \frac{1} 
{x^{2}} - \frac{1} 
{y^{2}} 
\over -\frac{1}
{y} + \frac{1} 
{x}} 
\]}
{\(\frac{x+y}
 {xy} \)}{\(-\frac{x+y} 
 {xy} \)}{\(\frac{1}
{y} -\frac{1} 
{x}\)}{\(\frac{1}
{x} -\frac{1} 
{y}\)}{}{}}
\End

\Data\ID{9000079210}
\Updated{2020-07-15 03:07:37}
\Level{A}
\SubArea{Polynomials and fractions}
\LANGen{
\Question{
Consider the expression 
\[ 
 V (x) = \frac{x} 
{x - 1} - \frac{1} 
{1 - x}.
\]
Find the correct ordering of the values
\(V (-2)\),
\(V (0)\) and
\(V (2)\).}
{\(V (0) < V (-2) < V (2)\)}{\(V (-2) < V (0) < V (2)\)}{\(V (0) < V (2) < V (-2)\)}{\(V (2) < V (0) < V (-2)\)}{}{}}
\LANGpl{
\Question{
Rozważ wyrażenie: 
\[ 
 V (x) = \frac{x} 
{x - 1} - \frac{1} 
{1 - x}.
\]
Znajdź prawidłową kolejność wartości 
\(V (-2)\),
\(V (0)\) i 
\(V (2)\).}
{\(V (0) < V (-2) < V (2)\)}{\(V (-2) < V (0) < V (2)\)}{\(V (0) < V (2) < V (-2)\)}{\(V (2) < V (0) < V (-2)\)}{}{}}
\LANGcs{
\Question{
Je dán výraz \(V (x) = \frac{x} 
{x-1} - \frac{1} 
{1-x}\).
Určete, která z následujících nerovností platí pro čísla
\(V (-2),V (0),V (2)\).}
{\(V (0) < V (-2) < V (2)\)}{\(V (-2) < V (0) < V (2)\)}{\(V (0) < V (2) < V (-2)\)}{\(V (2) < V (0) < V (-2)\)}{}{}}
\LANGsk{
\Question{
Daný je výraz
\[ 
 V (x) = \frac{x} 
{x - 1} - \frac{1} 
{1 - x}.
\]
Rozhodnite, ktoré usporiadanie čísel 
\(V (-2)\),
\(V (0)\) a
\(V (2)\) je správne.}
{\(V (0) < V (-2) < V (2)\)}{\(V (-2) < V (0) < V (2)\)}{\(V (0) < V (2) < V (-2)\)}{\(V (2) < V (0) < V (-2)\)}{}{}}
\LANGes{
\Question{
Consider the expression 
\[ 
 V (x) = \frac{x} 
{x - 1} - \frac{1} 
{1 - x}.
\]
Find the correct ordering of the values
\(V (-2)\),
\(V (0)\) and
\(V (2)\).}
{\(V (0) < V (-2) < V (2)\)}{\(V (-2) < V (0) < V (2)\)}{\(V (0) < V (2) < V (-2)\)}{\(V (2) < V (0) < V (-2)\)}{}{}}
\End

\Data\ID{9000083602}
\Updated{2020-07-15 03:07:37}
\Level{A}
\SubArea{Polynomials and fractions}
\LANGen{
\Question{
Evaluate the following expression at \(x = \frac{1} 
{2}\).
\[ 
 \frac{x^{2} - 2} 
{1 -\frac{1} 
{x}}
\]}
{\(\frac{7}
{4}\)}{\(-\frac{7} 
{4}\)}{\(\frac{7}
{2}\)}{\(-\frac{7} 
{2}\)}{}{}}
\LANGpl{
\Question{
Oblicz wartość podanego wyrażenia dla  \(x = \frac{1} 
{2}\).
\[ 
 \frac{x^{2} - 2} 
{1 -\frac{1} 
{x}}
\]}
{\(\frac{7}
{4}\)}{\(-\frac{7} 
{4}\)}{\(\frac{7}
{2}\)}{\(-\frac{7} 
{2}\)}{}{}}
\LANGcs{
\Question{
Hodnota výrazu \(\frac{x^{2}-2}
{1-\frac{1} 
{x}} \) 
pro \(x = \frac{1} 
{2}\) je
rovna číslu:}
{\(\frac{7}
{4}\)}{\(-\frac{7} 
{4}\)}{\(\frac{7}
{2}\)}{\(-\frac{7} 
{2}\)}{}{}}
\LANGsk{
\Question{
Určte hodnotu výrazu \(\frac{x^{2}-2}
{1-\frac{1} 
{x}} \) 
pre \(x = \frac{1} 
{2}\).}
{\(\frac{7}
{4}\)}{\(-\frac{7} 
{4}\)}{\(\frac{7}
{2}\)}{\(-\frac{7} 
{2}\)}{}{}}
\LANGes{
\Question{
Evaluate the following expression at \(x = \frac{1} 
{2}\).
\[ 
 \frac{x^{2} - 2} 
{1 -\frac{1} 
{x}}
\]}
{\(\frac{7}
{4}\)}{\(-\frac{7} 
{4}\)}{\(\frac{7}
{2}\)}{\(-\frac{7} 
{2}\)}{}{}}
\End

\Data\ID{9000083603}
\Updated{2020-07-15 03:07:37}
\Level{A}
\SubArea{Polynomials and fractions}
\LANGen{
\Question{
Evaluate the following expression at \(x = \frac{1} 
{2}\) 
and \(y = -\frac{1}
{4}\).
\[ 
 \frac{x -\frac{y} 
{x}} 
{1 + \frac{x} 
{y}}
\]}
{\(- 1\)}{\(3\)}{\(4\)}{\(1\)}{}{}}
\LANGpl{
\Question{
Oblicz wartość podanego wyrażenia dla  \(x = \frac{1} 
{2}\) 
i \(y = -\frac{1}
{4}\).
\[ 
 \frac{x -\frac{y} 
{x}} 
{1 + \frac{x} 
{y}}
\]}
{\(- 1\)}{\(3\)}{\(4\)}{\(1\)}{}{}}
\LANGcs{
\Question{
Hodnota výrazu \(\frac{x-\frac{y}
{x}}
{1+\frac{x}
{y}} \) 
pro \(x = \frac{1} 
{2}\) a
\(y = -\frac{1}
{4}\) je
rovna číslu:}
{\(- 1\)}{\(3\)}{\(4\)}{\(1\)}{}{}}
\LANGsk{
\Question{
Určte hodnotu výrazu \(\frac{x-\frac{y}
{x}}
{1+\frac{x}
{y}} \) 
pre \(x = \frac{1} 
{2}\) a
\(y = -\frac{1}
{4}\).}
{\(- 1\)}{\(3\)}{\(4\)}{\(1\)}{}{}}
\LANGes{
\Question{
Evaluate the following expression at \(x = \frac{1} 
{2}\) 
and \(y = -\frac{1}
{4}\).
\[ 
 \frac{x -\frac{y} 
{x}} 
{1 + \frac{x} 
{y}}
\]}
{\(- 1\)}{\(3\)}{\(4\)}{\(1\)}{}{}}
\End

\Data\ID{9000083604}
\Updated{2020-07-15 03:07:37}
\Level{A}
\SubArea{Polynomials and fractions}
\LANGen{
\Question{
Assuming \(x\neq - 1\),
\(x\neq \pm y\),
simplify the expression. 
\[ 
 \frac{x^{2} + 2xy + y^{2}} 
 {2x^{2} + 4x + 2} \cdot \frac{(x + 1)(y - x)}
 {y^{2} - x^{2}}
\]}
{\(\frac{x+y}
{2x+2}\)}{\(\frac{x+y}
 {2} \)}{\(x + y\)}{\(\frac{1}
{2}\)}{}{}}
\LANGpl{
\Question{
Zakładając, że  \(x\neq - 1\),
\(x\neq \pm y\),
uprość wyrażenie
\[ 
 \frac{x^{2} + 2xy + y^{2}} 
 {2x^{2} + 4x + 2} \cdot \frac{(x + 1)(y - x)}
 {y^{2} - x^{2}}
\]}
{\(\frac{x+y}
{2x+2}\)}{\(\frac{x+y}
 {2} \)}{\(x + y\)}{\(\frac{1}
{2}\)}{}{}}
\LANGcs{
\Question{
Úpravou lomeného výrazu \(\frac{x^{2}+2xy+y^{2}} 
 {2x^{2}+4x+2} \cdot \frac{(x+1)(y-x)}
 {y^{2}-x^{2}} \) 
za předpokladu, že \(x\neq - 1\),
\(x\neq \pm y\) 
získáme výraz:}
{\(\frac{x+y}
{2x+2}\)}{\(\frac{x+y}
 {2} \)}{\(x + y\)}{\(\frac{1}
{2}\)}{}{}}
\LANGsk{
\Question{
Zjednodušením lomeného výrazu \(\frac{x^{2}+2xy+y^{2}} 
 {2x^{2}+4x+2} \cdot \frac{(x+1)(y-x)}
 {y^{2}-x^{2}} \) 
za predpokladu, že \(x\neq - 1\),
\(x\neq \pm y\) 
 dostaneme výraz:}
{\(\frac{x+y}
{2x+2}\)}{\(\frac{x+y}
 {2} \)}{\(x + y\)}{\(\frac{1}
{2}\)}{}{}}
\LANGes{
\Question{
Assuming \(x\neq - 1\),
\(x\neq \pm y\),
simplify the expression. 
\[ 
 \frac{x^{2} + 2xy + y^{2}} 
 {2x^{2} + 4x + 2} \cdot \frac{(x + 1)(y - x)}
 {y^{2} - x^{2}}
\]}
{\(\frac{x+y}
{2x+2}\)}{\(\frac{x+y}
 {2} \)}{\(x + y\)}{\(\frac{1}
{2}\)}{}{}}
\End

\Data\ID{9000083605}
\Updated{2020-07-15 03:07:37}
\Level{A}
\SubArea{Polynomials and fractions}
\LANGen{
\Question{
Find the common denominator of the fractions. 
\[ 
 \text{$ \frac{3x}
{x^{2}+4x+4}$ and $ \frac{x+5}
{x^{2}-4}$}
\]}
{\((x + 2)^{2}(x - 2),\; x\neq \pm 2\)}{\((x + 2)(x - 4),\; x\neq \pm 2\)}{\((x + 2)^{2}(x - 4),\; x\neq \pm 2\)}{\((x + 2)(x - 4),\; x\neq \pm 2\)}{}{}}
\LANGpl{
\Question{
Znajdź wspólny mianownik wyrażeń 
\[ 
 \text{$ \frac{3x}
{x^{2}+4x+4}$ i $ \frac{x+5}
{x^{2}-4}$}
\]}
{\((x + 2)^{2}(x - 2),\; x\neq \pm 2\)}{\((x + 2)(x - 4),\; x\neq \pm 2\)}{\((x + 2)^{2}(x - 4),\; x\neq \pm 2\)}{\((x + 2)(x - 4),\; x\neq \pm 2\)}{}{}}
\LANGcs{
\Question{
Společný jmenovatel lomených výrazů
\(\frac{3x}
{x^{2}+4x+4}\) a
\(\frac{x+5}
{x^{2}-4}\) je:}
{\((x + 2)^{2}(x - 2),\; x\neq \pm 2\)}{\((x + 2)(x - 4),\; x\neq \pm 2\)}{\((x + 2)^{2}(x - 4),\; x\neq \pm 2\)}{\((x + 2)(x - 4),\; x\neq \pm 2\)}{}{}}
\LANGsk{
\Question{
Nájdite spoločný menovateľ daných lomených výrazov. 
\[ 
 \text{$ \frac{3x}
{x^{2}+4x+4}$ a $ \frac{x+5}
{x^{2}-4}$}
\]}
{\((x + 2)^{2}(x - 2),\; x\neq \pm 2\)}{\((x + 2)(x - 4),\; x\neq \pm 2\)}{\((x + 2)^{2}(x - 4),\; x\neq \pm 2\)}{\((x + 2)(x - 4),\; x\neq \pm 2\)}{}{}}
\LANGes{
\Question{
Find the common denominator of the fractions. 
\[ 
 \text{$ \frac{3x}
{x^{2}+4x+4}$ and $ \frac{x+5}
{x^{2}-4}$}
\]}
{\((x + 2)^{2}(x - 2),\; x\neq \pm 2\)}{\((x + 2)(x - 4),\; x\neq \pm 2\)}{\((x + 2)^{2}(x - 4),\; x\neq \pm 2\)}{\((x + 2)(x - 4),\; x\neq \pm 2\)}{}{}}
\End

\Data\ID{9000088803}
\Updated{2020-07-15 03:07:37}
\Level{A}
\SubArea{Polynomials and fractions}
\LANGen{
\Question{
Evaluate the following expression at \(x = \frac{1} 
{2}\).
\[ 
 1 - \frac{x - 2} 
{2x + 1}
\]}
{\(\frac{7}
{4}\)}{\(\frac{1}
{4}\)}{\(\frac{5}
{4}\)}{\(\frac{3}
{4}\)}{}{}}
\LANGpl{
\Question{
Oblicz wartości wyrażenia dla \(x = \frac{1} 
{2}\).
\[ 
 1 - \frac{x - 2} 
{2x + 1}
\]}
{\(\frac{7}
{4}\)}{\(\frac{1}
{4}\)}{\(\frac{5}
{4}\)}{\(\frac{3}
{4}\)}{}{}}
\LANGcs{
\Question{
Je dán výraz \(1 - \frac{x-2} 
{2x+1}\).
Hodnota výrazu pro \(x = \frac{1} 
{2}\) 
je rovna:}
{\(\frac{7}
{4}\)}{\(\frac{1}
{4}\)}{\(\frac{5}
{4}\)}{\(\frac{3}
{4}\)}{}{}}
\LANGsk{
\Question{
Určte hodnotu daného výrazu pre \(x = \frac{1} 
{2}\).
\[ 
 1 - \frac{x - 2} 
{2x + 1}
\]}
{\(\frac{7}
{4}\)}{\(\frac{1}
{4}\)}{\(\frac{5}
{4}\)}{\(\frac{3}
{4}\)}{}{}}
\LANGes{
\Question{
Evaluate the following expression at \(x = \frac{1} 
{2}\).
\[ 
 1 - \frac{x - 2} 
{2x + 1}
\]}
{\(\frac{7}
{4}\)}{\(\frac{1}
{4}\)}{\(\frac{5}
{4}\)}{\(\frac{3}
{4}\)}{}{}}
\End

\Data\ID{9000088804}
\Updated{2020-07-15 03:07:37}
\Level{A}
\SubArea{Polynomials and fractions}
\LANGen{
\Question{
Simplify the following expression. 
\[ 
 \frac{2s - 8rs} 
{16r^{2} - 1}
\]}
{\(- \frac{2s} 
{4r+1}\)}{\(\frac{2s}
{4r+1}\)}{\(\frac{2s}
{4r-1}\)}{\(\frac{2s}
{1-4r}\)}{}{}}
\LANGpl{
\Question{
Uprość podane wyrażenie.
\[ 
 \frac{2s - 8rs} 
{16r^{2} - 1}
\]}
{\(- \frac{2s} 
{4r+1}\)}{\(\frac{2s}
{4r+1}\)}{\(\frac{2s}
{4r-1}\)}{\(\frac{2s}
{1-4r}\)}{}{}}
\LANGcs{
\Question{
Výraz \(\frac{2s-8rs}
{16r^{2}-1}\) 
lze zjednodušit do tvaru:}
{\(- \frac{2s} 
{4r+1}\)}{\(\frac{2s}
{4r+1}\)}{\(\frac{2s}
{4r-1}\)}{\(\frac{2s}
{1-4r}\)}{}{}}
\LANGsk{
\Question{
Zjednodušte nasledujúci výraz. 
\[ 
 \frac{2s - 8rs} 
{16r^{2} - 1}
\]}
{\(- \frac{2s} 
{4r+1}\)}{\(\frac{2s}
{4r+1}\)}{\(\frac{2s}
{4r-1}\)}{\(\frac{2s}
{1-4r}\)}{}{}}
\LANGes{
\Question{
Simplify the following expression. 
\[ 
 \frac{2s - 8rs} 
{16r^{2} - 1}
\]}
{\(- \frac{2s} 
{4r+1}\)}{\(\frac{2s}
{4r+1}\)}{\(\frac{2s}
{4r-1}\)}{\(\frac{2s}
{1-4r}\)}{}{}}
\End

\Data\ID{9000088805}
\Updated{2020-07-15 03:07:37}
\Level{A}
\SubArea{Polynomials and fractions}
\LANGen{
\Question{
Simplify the following expression. 
\[ 
 \frac{a^{4} - 1} 
{1 - a^{2}}
\]}
{\(- a^{2} - 1\)}{\(a^{2} + 1\)}{\(a^{2} - 1\)}{\(1 - a^{2}\)}{}{}}
\LANGpl{
\Question{
Uprość podane wyrażenie.
\[ 
 \frac{a^{4} - 1} 
{1 - a^{2}}
\]}
{\(- a^{2} - 1\)}{\(a^{2} + 1\)}{\(a^{2} - 1\)}{\(1 - a^{2}\)}{}{}}
\LANGcs{
\Question{
Zjednodušením výrazu \(\frac{a^{4}-1}
{1-a^{2}} \) 
dostaneme:}
{\(- a^{2} - 1\)}{\(a^{2} + 1\)}{\(a^{2} - 1\)}{\(1 - a^{2}\)}{}{}}
\LANGsk{
\Question{
Zjednodušte nasledujúci výraz. 
\[ 
 \frac{a^{4} - 1} 
{1 - a^{2}}
\]}
{\(- a^{2} - 1\)}{\(a^{2} + 1\)}{\(a^{2} - 1\)}{\(1 - a^{2}\)}{}{}}
\LANGes{
\Question{
Simplify the following expression. 
\[ 
 \frac{a^{4} - 1} 
{1 - a^{2}}
\]}
{\(- a^{2} - 1\)}{\(a^{2} + 1\)}{\(a^{2} - 1\)}{\(1 - a^{2}\)}{}{}}
\End

\Data\ID{9000088809}
\Updated{2020-07-15 03:07:37}
\Level{A}
\SubArea{Polynomials and fractions}
\LANGen{
\Question{
Simplify the following expression. 
\[ 
 \left ( \frac{1}
{m - n} - \frac{1} 
{m + n}\right )\cdot \left (\frac{m^{2} + 2mn + n^{2}}
 {2n} \right )
\]}
{\(\frac{m+n}
{m-n}\)}{\(0\)}{\(\frac{m(m+n)}
{n(m-n)} \)}{\(2\)}{}{}}
\LANGpl{
\Question{
Uprość podane wyrażenie.
\[ 
 \left ( \frac{1}
{m - n} - \frac{1} 
{m + n}\right )\cdot \left (\frac{m^{2} + 2mn + n^{2}}
 {2n} \right )
\]}
{\(\frac{m+n}
{m-n}\)}{\(0\)}{\(\frac{m(m+n)}
{n(m-n)} \)}{\(2\)}{}{}}
\LANGcs{
\Question{
Úpravou výrazu \(\left ( \frac{1}
{m-n} - \frac{1} 
{m+n}\right )\cdot \left (\frac{m^{2}+2mn+n^{2}} 
 {2n} \right )\) 
dostaneme:}
{\(\frac{m+n}
{m-n}\)}{\(0\)}{\(\frac{m(m+n)}
{n(m-n)} \)}{\(2\)}{}{}}
\LANGsk{
\Question{
Zjednodušte nasledujúci výraz.
\[ 
 \left ( \frac{1}
{m - n} - \frac{1} 
{m + n}\right )\cdot \left (\frac{m^{2} + 2mn + n^{2}}
 {2n} \right )
\]}
{\(\frac{m+n}
{m-n}\)}{\(0\)}{\(\frac{m(m+n)}
{n(m-n)} \)}{\(2\)}{}{}}
\LANGes{
\Question{
Simplify the following expression. 
\[ 
 \left ( \frac{1}
{m - n} - \frac{1} 
{m + n}\right )\cdot \left (\frac{m^{2} + 2mn + n^{2}}
 {2n} \right )
\]}
{\(\frac{m+n}
{m-n}\)}{\(0\)}{\(\frac{m(m+n)}
{n(m-n)} \)}{\(2\)}{}{}}
\End

\Data\ID{9000088810}
\Updated{2020-07-15 03:07:37}
\Level{A}
\SubArea{Polynomials and fractions}
\LANGen{
\Question{
Simplify the following expression. 
\[ 
 \left (x -\frac{1} 
{x}\right )\cdot \left (1 - \frac{x} 
{x + 1}\right )
\]}
{\(\frac{x - 1}
 {x} \)}{\(\frac{x - 1}
{x + 1}\)}{\(\frac{1 - x}
{x + 1}\)}{\(\frac{1 - x}
 {x} \)}{}{}}
\LANGpl{
\Question{
Uprość podane wyrażenie.
\[ 
 \left (x -\frac{1} 
{x}\right )\cdot \left (1 - \frac{x} 
{x + 1}\right )
\]}
{\(\frac{x - 1}
 {x} \)}{\(\frac{x - 1}
{x + 1}\)}{\(\frac{1 - x}
{x + 1}\)}{\(\frac{1 - x}
 {x} \)}{}{}}
\LANGcs{
\Question{
Úpravou výrazu \(\left (x -\frac{1} 
{x}\right )\cdot \left (1 - \frac{x} 
{x+1}\right )\) 
dostaneme:}
{\(\frac{x - 1}
 {x} \)}{\(\frac{x - 1}
{x + 1}\)}{\(\frac{1 - x}
{x + 1}\)}{\(\frac{1 - x}
 {x} \)}{}{}}
\LANGsk{
\Question{
Zjednodušte nasledujúci výraz. 
\[ 
 \left (x -\frac{1} 
{x}\right )\cdot \left (1 - \frac{x} 
{x + 1}\right )
\]}
{\(\frac{x - 1}
 {x} \)}{\(\frac{x - 1}
{x + 1}\)}{\(\frac{1 - x}
{x + 1}\)}{\(\frac{1 - x}
 {x} \)}{}{}}
\LANGes{
\Question{
Simplify the following expression. 
\[ 
 \left (x -\frac{1} 
{x}\right )\cdot \left (1 - \frac{x} 
{x + 1}\right )
\]}
{\(\frac{x - 1}
 {x} \)}{\(\frac{x - 1}
{x + 1}\)}{\(\frac{1 - x}
{x + 1}\)}{\(\frac{1 - x}
 {x} \)}{}{}}
\End

\Data\ID{9000101601}
\Updated{2020-07-15 03:07:37}
\Level{A}
\SubArea{Polynomials and fractions}
\LANGen{
\Question{
Expand \((1 + x)\left (x^{2} + x - 1\right )(1 - x)\).}
{\(- x^{4} - x^{3} + 2x^{2} + x - 1\)}{\(x^{4} - x^{3} + 2x^{2} + x + 1\)}{\(- x^{4} + x^{3} - 1\)}{\(x^{4} + x^{3} - 2x^{2} + x - 1\)}{}{}}
\LANGpl{
\Question{
Rozwiń  \((1 + x)\left (x^{2} + x - 1\right )(1 - x)\).}
{\(- x^{4} - x^{3} + 2x^{2} + x - 1\)}{\(x^{4} - x^{3} + 2x^{2} + x + 1\)}{\(- x^{4} + x^{3} - 1\)}{\(x^{4} + x^{3} - 2x^{2} + x - 1\)}{}{}}
\LANGcs{
\Question{
Úpravou výrazu \((1 + x)\left (x^{2} + x - 1\right )(1 - x)\) 
získáme:}
{\(- x^{4} - x^{3} + 2x^{2} + x - 1\)}{\(x^{4} - x^{3} + 2x^{2} + x + 1\)}{\(- x^{4} + x^{3} - 1\)}{\(x^{4} + x^{3} - 2x^{2} + x - 1\)}{}{}}
\LANGsk{
\Question{
Upravte výraz \((1 + x)\left (x^{2} + x - 1\right )(1 - x)\).}
{\(- x^{4} - x^{3} + 2x^{2} + x - 1\)}{\(x^{4} - x^{3} + 2x^{2} + x + 1\)}{\(- x^{4} + x^{3} - 1\)}{\(x^{4} + x^{3} - 2x^{2} + x - 1\)}{}{}}
\LANGes{
\Question{
Expand \((1 + x)\left (x^{2} + x - 1\right )(1 - x)\).}
{\(- x^{4} - x^{3} + 2x^{2} + x - 1\)}{\(x^{4} - x^{3} + 2x^{2} + x + 1\)}{\(- x^{4} + x^{3} - 1\)}{\(x^{4} + x^{3} - 2x^{2} + x - 1\)}{}{}}
\End

\Data\ID{9000101602}
\Updated{2020-07-15 03:07:37}
\Level{A}
\SubArea{Polynomials and fractions}
\LANGen{
\Question{
Simplify the polynomial \((x - 1)(x + 1)\left (x^{2} + 1\right ) -\left (x^{2} - 1\right )^{2}\) 
into one of the following forms.}
{\(2\left (x^{2} - 1\right )\)}{\(0\)}{\(2\left (x^{2} - 1\right )(x + 1)\)}{\(x^{2} - 1\)}{}{}}
\LANGpl{
\Question{
Doprowadź podany wielomian \((x - 1)(x + 1)\left (x^{2} + 1\right ) -\left (x^{2} - 1\right )^{2}\) 
do jednej z podanych postaci.}
{\(2\left (x^{2} - 1\right )\)}{\(0\)}{\(2\left (x^{2} - 1\right )(x + 1)\)}{\(x^{2} - 1\)}{}{}}
\LANGcs{
\Question{
Úpravou výrazu \((x - 1)(x + 1)\left (x^{2} + 1\right ) -\left (x^{2} - 1\right )^{2}\) 
získáme:}
{\(2\left (x^{2} - 1\right )\)}{\(0\)}{\(2\left (x^{2} - 1\right )(x + 1)\)}{\(x^{2} - 1\)}{}{}}
\LANGsk{
\Question{
Upravte daný výraz \((x - 1)(x + 1)\left (x^{2} + 1\right ) -\left (x^{2} - 1\right )^{2}\).}
{\(2\left (x^{2} - 1\right )\)}{\(0\)}{\(2\left (x^{2} - 1\right )(x + 1)\)}{\(x^{2} - 1\)}{}{}}
\LANGes{
\Question{
Simplify the polynomial \((x - 1)(x + 1)\left (x^{2} + 1\right ) -\left (x^{2} - 1\right )^{2}\) 
into one of the following forms.}
{\(2\left (x^{2} - 1\right )\)}{\(0\)}{\(2\left (x^{2} - 1\right )(x + 1)\)}{\(x^{2} - 1\)}{}{}}
\End

\Data\ID{9000101603}
\Updated{2020-07-15 03:07:37}
\Level{A}
\SubArea{Polynomials and fractions}
\LANGen{
\Question{
Simplify the polynomial \((x + 1)(x - 1)^{2} - (x - 1)(x + 1)^{2}\) 
into one of the following forms.}
{\(- 2\left (x - 1\right )\left (x + 1\right )\)}{\(2\left (x - 1\right )\left (x + 1\right )\)}{\(0\)}{\(2\)}{}{}}
\LANGpl{
\Question{
Doprowadź podany wielomian  \((x + 1)(x - 1)^{2} - (x - 1)(x + 1)^{2}\) 
do jednej z podanych postaci.}
{\(- 2\left (x - 1\right )\left (x + 1\right )\)}{\(2\left (x - 1\right )\left (x + 1\right )\)}{\(0\)}{\(2\)}{}{}}
\LANGcs{
\Question{
Úpravou výrazu \((x + 1)(x - 1)^{2} - (x - 1)(x + 1)^{2}\) 
získáme:}
{\(- 2\left (x - 1\right )\left (x + 1\right )\)}{\(2\left (x - 1\right )\left (x + 1\right )\)}{\(0\)}{\(2\)}{}{}}
\LANGsk{
\Question{
Upravte daný výraz \((x + 1)(x - 1)^{2} - (x - 1)(x + 1)^{2}\).}
{\(- 2\left (x - 1\right )\left (x + 1\right )\)}{\(2\left (x - 1\right )\left (x + 1\right )\)}{\(0\)}{\(2\)}{}{}}
\LANGes{
\Question{
Simplify the polynomial \((x + 1)(x - 1)^{2} - (x - 1)(x + 1)^{2}\) 
into one of the following forms.}
{\(- 2\left (x - 1\right )\left (x + 1\right )\)}{\(2\left (x - 1\right )\left (x + 1\right )\)}{\(0\)}{\(2\)}{}{}}
\End

\Data\ID{9000101604}
\Updated{2020-07-15 03:07:37}
\Level{A}
\SubArea{Polynomials and fractions}
\LANGen{
\Question{
Expand the polynomial \(\left (2x^{2} + 4x\right )^{2} -\left (4x - 2x^{2}\right )^{2}\).}
{\(32x^{3}\)}{\(0\)}{\(32x^{3} - 8x\)}{\(32x^{3} - 32x^{2} + 8x\)}{}{}}
\LANGpl{
\Question{
Rozwiń wielomian  \(\left (2x^{2} + 4x\right )^{2} -\left (4x - 2x^{2}\right )^{2}\).}
{\(32x^{3}\)}{\(0\)}{\(32x^{3} - 8x\)}{\(32x^{3} - 32x^{2} + 8x\)}{}{}}
\LANGcs{
\Question{
Úpravou výrazu \(\left (2x^{2} + 4x\right )^{2} -\left (4x - 2x^{2}\right )^{2}\) 
získáme:}
{\(32x^{3}\)}{\(0\)}{\(32x^{3} - 8x\)}{\(32x^{3} - 32x^{2} + 8x\)}{}{}}
\LANGsk{
\Question{
Upravte daný výraz \(\left (2x^{2} + 4x\right )^{2} -\left (4x - 2x^{2}\right )^{2}\).}
{\(32x^{3}\)}{\(0\)}{\(32x^{3} - 8x\)}{\(32x^{3} - 32x^{2} + 8x\)}{}{}}
\LANGes{
\Question{
Expand the polynomial \(\left (2x^{2} + 4x\right )^{2} -\left (4x - 2x^{2}\right )^{2}\).}
{\(32x^{3}\)}{\(0\)}{\(32x^{3} - 8x\)}{\(32x^{3} - 32x^{2} + 8x\)}{}{}}
\End

\Data\ID{9000146203}
\Updated{2020-07-15 03:07:37}
\Level{A}
\SubArea{Polynomials and fractions}
\LANGen{
\Question{
Expand the following expression. 
\[ 
 \left (x^{5} -\sqrt{2}y\right )^{2}
\]}
{\(x^{10} - 2\sqrt{2}x^{5}y + 2y^{2}\)}{\(x^{10} -\sqrt{2}x^{5}y + 2y^{2}\)}{\(x^{10} - 2\sqrt{2}x^{5}y - 2y^{2}\)}{\(x^{10} -\sqrt{2}x^{5}y - 2y^{2}\)}{}{}}
\LANGpl{
\Question{
Rozwiń podane wyrażenie:
\[ 
 \left (x^{5} -\sqrt{2}y\right )^{2}
\]}
{\(x^{10} - 2\sqrt{2}x^{5}y + 2y^{2}\)}{\(x^{10} -\sqrt{2}x^{5}y + 2y^{2}\)}{\(x^{10} - 2\sqrt{2}x^{5}y - 2y^{2}\)}{\(x^{10} -\sqrt{2}x^{5}y - 2y^{2}\)}{}{}}
\LANGcs{
\Question{
Umocněním \(\left (x^{5} -\sqrt{2}y\right )^{2}\) 
získáme výraz:}
{\(x^{10} - 2\sqrt{2}x^{5}y + 2y^{2}\)}{\(x^{10} -\sqrt{2}x^{5}y + 2y^{2}\)}{\(x^{10} - 2\sqrt{2}x^{5}y - 2y^{2}\)}{\(x^{10} -\sqrt{2}x^{5}y - 2y^{2}\)}{}{}}
\LANGsk{
\Question{
Umocnite daný výraz.
\[ 
 \left (x^{5} -\sqrt{2}y\right )^{2}
\]}
{\(x^{10} - 2\sqrt{2}x^{5}y + 2y^{2}\)}{\(x^{10} -\sqrt{2}x^{5}y + 2y^{2}\)}{\(x^{10} - 2\sqrt{2}x^{5}y - 2y^{2}\)}{\(x^{10} -\sqrt{2}x^{5}y - 2y^{2}\)}{}{}}
\LANGes{
\Question{
Expand the following expression. 
\[ 
 \left (x^{5} -\sqrt{2}y\right )^{2}
\]}
{\(x^{10} - 2\sqrt{2}x^{5}y + 2y^{2}\)}{\(x^{10} -\sqrt{2}x^{5}y + 2y^{2}\)}{\(x^{10} - 2\sqrt{2}x^{5}y - 2y^{2}\)}{\(x^{10} -\sqrt{2}x^{5}y - 2y^{2}\)}{}{}}
\End

\Data\ID{9000146204}
\Updated{2020-07-15 03:07:37}
\Level{A}
\SubArea{Polynomials and fractions}
\LANGen{
\Question{
Expand the following expression. 
\[ 
 \left (\frac{a}
{2} + 4b^{3}\right )^{2}
\]}
{\(\frac{a^{2}} 
 {4} + 4ab^{3} + 16b^{6}\)}{\(\frac{a^{2}} 
 {4} + 2ab^{3} + 16b^{6}\)}{\(\frac{a^{2}} 
 {4} + 4ab^{3} + 16b^{5}\)}{\(\frac{a^{2}} 
 {4} + 2ab^{3} + 16b^{5}\)}{}{}}
\LANGpl{
\Question{
Rozwiń podane wyrażenie:
\[ 
 \left (\frac{a}
{2} + 4b^{3}\right )^{2}
\]}
{\(\frac{a^{2}} 
 {4} + 4ab^{3} + 16b^{6}\)}{\(\frac{a^{2}} 
 {4} + 2ab^{3} + 16b^{6}\)}{\(\frac{a^{2}} 
 {4} + 4ab^{3} + 16b^{5}\)}{\(\frac{a^{2}} 
 {4} + 2ab^{3} + 16b^{5}\)}{}{}}
\LANGcs{
\Question{
Umocněním \(\left (\frac{a}
{2} + 4b^{3}\right )^{2}\) 
získáme výraz:}
{\(\frac{a^{2}} 
 {4} + 4ab^{3} + 16b^{6}\)}{\(\frac{a^{2}} 
 {4} + 2ab^{3} + 16b^{6}\)}{\(\frac{a^{2}} 
 {4} + 4ab^{3} + 16b^{5}\)}{\(\frac{a^{2}} 
 {4} + 2ab^{3} + 16b^{5}\)}{}{}}
\LANGsk{
\Question{
Umocnite daný výraz.
\[ 
 \left (\frac{a}
{2} + 4b^{3}\right )^{2}
\]}
{\(\frac{a^{2}} 
 {4} + 4ab^{3} + 16b^{6}\)}{\(\frac{a^{2}} 
 {4} + 2ab^{3} + 16b^{6}\)}{\(\frac{a^{2}} 
 {4} + 4ab^{3} + 16b^{5}\)}{\(\frac{a^{2}} 
 {4} + 2ab^{3} + 16b^{5}\)}{}{}}
\LANGes{
\Question{
Expand the following expression. 
\[ 
 \left (\frac{a}
{2} + 4b^{3}\right )^{2}
\]}
{\(\frac{a^{2}} 
 {4} + 4ab^{3} + 16b^{6}\)}{\(\frac{a^{2}} 
 {4} + 2ab^{3} + 16b^{6}\)}{\(\frac{a^{2}} 
 {4} + 4ab^{3} + 16b^{5}\)}{\(\frac{a^{2}} 
 {4} + 2ab^{3} + 16b^{5}\)}{}{}}
\End

\Data\ID{9000146205}
\Updated{2020-07-15 03:07:37}
\Level{A}
\SubArea{Polynomials and fractions}
\LANGen{
\Question{
Factor the following expression. 
\[ 
 9a^{6} - 4b^{2}
\]}
{\(\left (3a^{3} - 2b\right )\left (3a^{3} + 2b\right )\)}{\(\left (3a^{3} - 2b\right )\left (3a^{3} - 2b\right )\)}{\(\left (3a^{3} - 2b\right )\left (3a^{2} + 2b\right )\)}{\(\left (3a^{3} - 2b\right )\left (3a^{2} - 2b\right )\)}{}{}}
\LANGpl{
\Question{
Rozłóż na czynniki wyrażenie:
\[ 
 9a^{6} - 4b^{2}
\]}
{\(\left (3a^{3} - 2b\right )\left (3a^{3} + 2b\right )\)}{\(\left (3a^{3} - 2b\right )\left (3a^{3} - 2b\right )\)}{\(\left (3a^{3} - 2b\right )\left (3a^{2} + 2b\right )\)}{\(\left (3a^{3} - 2b\right )\left (3a^{2} - 2b\right )\)}{}{}}
\LANGcs{
\Question{
Rozložením výrazu \(9a^{6} - 4b^{2}\) 
na součin získáme výsledek:}
{\(\left (3a^{3} - 2b\right )\left (3a^{3} + 2b\right )\)}{\(\left (3a^{3} - 2b\right )\left (3a^{3} - 2b\right )\)}{\(\left (3a^{3} - 2b\right )\left (3a^{2} + 2b\right )\)}{\(\left (3a^{3} - 2b\right )\left (3a^{2} - 2b\right )\)}{}{}}
\LANGsk{
\Question{
Upravte na súčin.
\[ 
 9a^{6} - 4b^{2}
\]}
{\(\left (3a^{3} - 2b\right )\left (3a^{3} + 2b\right )\)}{\(\left (3a^{3} - 2b\right )\left (3a^{3} - 2b\right )\)}{\(\left (3a^{3} - 2b\right )\left (3a^{2} + 2b\right )\)}{\(\left (3a^{3} - 2b\right )\left (3a^{2} - 2b\right )\)}{}{}}
\LANGes{
\Question{
Factor the following expression. 
\[ 
 9a^{6} - 4b^{2}
\]}
{\(\left (3a^{3} - 2b\right )\left (3a^{3} + 2b\right )\)}{\(\left (3a^{3} - 2b\right )\left (3a^{3} - 2b\right )\)}{\(\left (3a^{3} - 2b\right )\left (3a^{2} + 2b\right )\)}{\(\left (3a^{3} - 2b\right )\left (3a^{2} - 2b\right )\)}{}{}}
\End

\Data\ID{9000146206}
\Updated{2020-07-15 03:07:37}
\Level{A}
\SubArea{Polynomials and fractions}
\LANGen{
\Question{
Factor the following expression. 
\[ 
 x^{2}y^{10} - 81
\]}
{\(\left (xy^{5} - 9\right )\left (xy^{5} + 9\right )\)}{\(\left (xy^{5} - 9\right )\left (xy^{5} - 9\right )\)}{\(\left (xy^{5} - 9\right )\left (xy^{2} + 9\right )\)}{\(\left (xy^{5} - 9\right )\left (xy^{2} - 9\right )\)}{}{}}
\LANGpl{
\Question{
Rozłóż na czynniki wyrażenie:
\[ 
 x^{2}y^{10} - 81
\]}
{\(\left (xy^{5} - 9\right )\left (xy^{5} + 9\right )\)}{\(\left (xy^{5} - 9\right )\left (xy^{5} - 9\right )\)}{\(\left (xy^{5} - 9\right )\left (xy^{2} + 9\right )\)}{\(\left (xy^{5} - 9\right )\left (xy^{2} - 9\right )\)}{}{}}
\LANGcs{
\Question{
Rozložením výrazu \(x^{2}y^{10} - 81\) 
na součin získáme výsledek:}
{\(\left (xy^{5} - 9\right )\left (xy^{5} + 9\right )\)}{\(\left (xy^{5} - 9\right )\left (xy^{5} - 9\right )\)}{\(\left (xy^{5} - 9\right )\left (xy^{2} + 9\right )\)}{\(\left (xy^{5} - 9\right )\left (xy^{2} - 9\right )\)}{}{}}
\LANGsk{
\Question{
Upravte na súčin.
\[ 
 x^{2}y^{10} - 81
\]}
{\(\left (xy^{5} - 9\right )\left (xy^{5} + 9\right )\)}{\(\left (xy^{5} - 9\right )\left (xy^{5} - 9\right )\)}{\(\left (xy^{5} - 9\right )\left (xy^{2} + 9\right )\)}{\(\left (xy^{5} - 9\right )\left (xy^{2} - 9\right )\)}{}{}}
\LANGes{
\Question{
Factor the following expression. 
\[ 
 x^{2}y^{10} - 81
\]}
{\(\left (xy^{5} - 9\right )\left (xy^{5} + 9\right )\)}{\(\left (xy^{5} - 9\right )\left (xy^{5} - 9\right )\)}{\(\left (xy^{5} - 9\right )\left (xy^{2} + 9\right )\)}{\(\left (xy^{5} - 9\right )\left (xy^{2} - 9\right )\)}{}{}}
\End

\Data\ID{9000146701}
\Updated{2020-07-15 03:07:37}
\Level{A}
\SubArea{Polynomials and fractions}
\LANGen{
\Question{
Expand the following polynomial. 
\[ 
 2 - (2x + 1) + x(5 - 2x) - 3(x - 2)
\]}
{\(- 2x^{2} + 7\)}{\(- 2x^{2} + 9\)}{\(- 2x^{2} - 3\)}{\(- 2x^{2} - 5\)}{}{}}
\LANGpl{
\Question{
Rozwiń podany wielomian:
\[ 
 2 - (2x + 1) + x(5 - 2x) - 3(x - 2)
\]}
{\(- 2x^{2} + 7\)}{\(- 2x^{2} + 9\)}{\(- 2x^{2} - 3\)}{\(- 2x^{2} - 5\)}{}{}}
\LANGcs{
\Question{
Úpravou výrazu \(2 - (2x + 1) + x(5 - 2x) - 3(x - 2)\) 
získáme dvojčlen:}
{\(- 2x^{2} + 7\)}{\(- 2x^{2} + 9\)}{\(- 2x^{2} - 3\)}{\(- 2x^{2} - 5\)}{}{}}
\LANGsk{
\Question{
Upravte daný výraz.
\[ 
 2 - (2x + 1) + x(5 - 2x) - 3(x - 2)
\]}
{\(- 2x^{2} + 7\)}{\(- 2x^{2} + 9\)}{\(- 2x^{2} - 3\)}{\(- 2x^{2} - 5\)}{}{}}
\LANGes{
\Question{
Expand the following polynomial. 
\[ 
 2 - (2x + 1) + x(5 - 2x) - 3(x - 2)
\]}
{\(- 2x^{2} + 7\)}{\(- 2x^{2} + 9\)}{\(- 2x^{2} - 3\)}{\(- 2x^{2} - 5\)}{}{}}
\End

\Data\ID{9000146702}
\Updated{2020-07-15 03:07:37}
\Level{A}
\SubArea{Polynomials and fractions}
\LANGen{
\Question{
Expand the following expression. 
\[ 
 a - 4(2 - a) - a(5a + 1) + 2a(3 - 2a)
\]}
{\(- 9a^{2} + 10a - 8\)}{\(- 9a^{2} + 12a - 8\)}{\(- 9a^{2} + 2a - 8\)}{\(- 9a^{2} + 4a - 8\)}{}{}}
\LANGpl{
\Question{
Rozwiń podane wyrażenie:
\[ 
 a - 4(2 - a) - a(5a + 1) + 2a(3 - 2a)
\]}
{\(- 9a^{2} + 10a - 8\)}{\(- 9a^{2} + 12a - 8\)}{\(- 9a^{2} + 2a - 8\)}{\(- 9a^{2} + 4a - 8\)}{}{}}
\LANGcs{
\Question{
Úpravou výrazu \(a - 4(2 - a) - a(5a + 1) + 2a(3 - 2a)\) 
získáme trojčlen:}
{\(- 9a^{2} + 10a - 8\)}{\(- 9a^{2} + 12a - 8\)}{\(- 9a^{2} + 2a - 8\)}{\(- 9a^{2} + 4a - 8\)}{}{}}
\LANGsk{
\Question{
Upravte daný výraz.
\[ 
 a - 4(2 - a) - a(5a + 1) + 2a(3 - 2a)
\]}
{\(- 9a^{2} + 10a - 8\)}{\(- 9a^{2} + 12a - 8\)}{\(- 9a^{2} + 2a - 8\)}{\(- 9a^{2} + 4a - 8\)}{}{}}
\LANGes{
\Question{
Expand the following expression. 
\[ 
 a - 4(2 - a) - a(5a + 1) + 2a(3 - 2a)
\]}
{\(- 9a^{2} + 10a - 8\)}{\(- 9a^{2} + 12a - 8\)}{\(- 9a^{2} + 2a - 8\)}{\(- 9a^{2} + 4a - 8\)}{}{}}
\End

\Data\ID{9000146703}
\Updated{2020-07-15 03:07:37}
\Level{A}
\SubArea{Polynomials and fractions}
\LANGen{
\Question{
Expand the following expression. 
\[ 
 (a - 2)(5a + 3) - (2a + 1)(3 - a)
\]}
{\(7a^{2} - 12a - 9\)}{\(3a^{2} - 12a - 9\)}{\(7a^{2} - 2a - 9\)}{\(3a^{2} - 2a - 9\)}{}{}}
\LANGpl{
\Question{
Rozwiń podany wielomian:
\[ 
 (a - 2)(5a + 3) - (2a + 1)(3 - a)
\]}
{\(7a^{2} - 12a - 9\)}{\(3a^{2} - 12a - 9\)}{\(7a^{2} - 2a - 9\)}{\(3a^{2} - 2a - 9\)}{}{}}
\LANGcs{
\Question{
Úpravou výrazu \((a - 2)(5a + 3) - (2a + 1)(3 - a)\) 
získáme trojčlen:}
{\(7a^{2} - 12a - 9\)}{\(3a^{2} - 12a - 9\)}{\(7a^{2} - 2a - 9\)}{\(3a^{2} - 2a - 9\)}{}{}}
\LANGsk{
\Question{
Upravte daný výraz.
\[ 
 (a - 2)(5a + 3) - (2a + 1)(3 - a)
\]}
{\(7a^{2} - 12a - 9\)}{\(3a^{2} - 12a - 9\)}{\(7a^{2} - 2a - 9\)}{\(3a^{2} - 2a - 9\)}{}{}}
\LANGes{
\Question{
Expand the following expression. 
\[ 
 (a - 2)(5a + 3) - (2a + 1)(3 - a)
\]}
{\(7a^{2} - 12a - 9\)}{\(3a^{2} - 12a - 9\)}{\(7a^{2} - 2a - 9\)}{\(3a^{2} - 2a - 9\)}{}{}}
\End

\Data\ID{9000146704}
\Updated{2020-07-15 03:07:37}
\Level{A}
\SubArea{Polynomials and fractions}
\LANGen{
\Question{
Expand the following polynomial. 
\[ 
 (3 - x)(x - 2) - (x + 1)(x - 3)
\]}
{\(- 2x^{2} + 7x - 3\)}{\(- 2x^{2} + 3x - 9\)}{\(- 2x^{2} + 3x - 3\)}{\(- 2x^{2} + 7x - 9\)}{}{}}
\LANGpl{
\Question{
Rozwiń podany wielomian:
\[ 
 (3 - x)(x - 2) - (x + 1)(x - 3)
\]}
{\(- 2x^{2} + 7x - 3\)}{\(- 2x^{2} + 3x - 9\)}{\(- 2x^{2} + 3x - 3\)}{\(- 2x^{2} + 7x - 9\)}{}{}}
\LANGcs{
\Question{
Úpravou výrazu \((3 - x)(x - 2) - (x + 1)(x - 3)\) 
získáme trojčlen:}
{\(- 2x^{2} + 7x - 3\)}{\(- 2x^{2} + 3x - 9\)}{\(- 2x^{2} + 3x - 3\)}{\(- 2x^{2} + 7x - 9\)}{}{}}
\LANGsk{
\Question{
Upravte daný výraz.
\[ 
 (3 - x)(x - 2) - (x + 1)(x - 3)
\]}
{\(- 2x^{2} + 7x - 3\)}{\(- 2x^{2} + 3x - 9\)}{\(- 2x^{2} + 3x - 3\)}{\(- 2x^{2} + 7x - 9\)}{}{}}
\LANGes{
\Question{
Expand the following polynomial. 
\[ 
 (3 - x)(x - 2) - (x + 1)(x - 3)
\]}
{\(- 2x^{2} + 7x - 3\)}{\(- 2x^{2} + 3x - 9\)}{\(- 2x^{2} + 3x - 3\)}{\(- 2x^{2} + 7x - 9\)}{}{}}
\End

\Data\ID{9000146705}
\Updated{2020-07-15 03:07:37}
\Level{A}
\SubArea{Polynomials and fractions}
\LANGen{
\Question{
Expand the following expression. 
\[ 
 (3 - 2a)^{2} - (3a - 4)(3a + 4)
\]}
{\(- 5a^{2} - 12a + 25\)}{\(- 5a^{2} - 12a - 7\)}{\(- 5a^{2} + 25\)}{\(- 5a^{2} - 7\)}{}{}}
\LANGpl{
\Question{
Rozwiń podany wielomian:
\[ 
 (3 - 2a)^{2} - (3a - 4)(3a + 4)
\]}
{\(- 5a^{2} - 12a + 25\)}{\(- 5a^{2} - 12a - 7\)}{\(- 5a^{2} + 25\)}{\(- 5a^{2} - 7\)}{}{}}
\LANGcs{
\Question{
Úpravou výrazu \((3 - 2a)^{2} - (3a - 4)(3a + 4)\) 
získáme mnohočlen:}
{\(- 5a^{2} - 12a + 25\)}{\(- 5a^{2} - 12a - 7\)}{\(- 5a^{2} + 25\)}{\(- 5a^{2} - 7\)}{}{}}
\LANGsk{
\Question{
Upravte daný výraz.
\[ 
 (3 - 2a)^{2} - (3a - 4)(3a + 4)
\]}
{\(- 5a^{2} - 12a + 25\)}{\(- 5a^{2} - 12a - 7\)}{\(- 5a^{2} + 25\)}{\(- 5a^{2} - 7\)}{}{}}
\LANGes{
\Question{
Expand the following expression. 
\[ 
 (3 - 2a)^{2} - (3a - 4)(3a + 4)
\]}
{\(- 5a^{2} - 12a + 25\)}{\(- 5a^{2} - 12a - 7\)}{\(- 5a^{2} + 25\)}{\(- 5a^{2} - 7\)}{}{}}
\End

\Data\ID{9000146706}
\Updated{2020-07-15 03:07:37}
\Level{A}
\SubArea{Polynomials and fractions}
\LANGen{
\Question{
Expand the following polynomial. 
\[ 
 (2x + 3)^{2} - (2 - x)^{2}
\]}
{\(3x^{2} + 16x + 5\)}{\(3x^{2} + 8x + 13\)}{\(3x^{2} + 13\)}{\(3x^{2} + 5\)}{}{}}
\LANGpl{
\Question{
Rozwiń podany wielomian:
\[ 
 (2x + 3)^{2} - (2 - x)^{2}
\]}
{\(3x^{2} + 16x + 5\)}{\(3x^{2} + 8x + 13\)}{\(3x^{2} + 13\)}{\(3x^{2} + 5\)}{}{}}
\LANGcs{
\Question{
Úpravou výrazu \((2x + 3)^{2} - (2 - x)^{2}\) 
získáme mnohočlen:}
{\(3x^{2} + 16x + 5\)}{\(3x^{2} + 8x + 13\)}{\(3x^{2} + 13\)}{\(3x^{2} + 5\)}{}{}}
\LANGsk{
\Question{
Upravte daný výraz.
\[ 
 (2x + 3)^{2} - (2 - x)^{2}
\]}
{\(3x^{2} + 16x + 5\)}{\(3x^{2} + 8x + 13\)}{\(3x^{2} + 13\)}{\(3x^{2} + 5\)}{}{}}
\LANGes{
\Question{
Expand the following polynomial. 
\[ 
 (2x + 3)^{2} - (2 - x)^{2}
\]}
{\(3x^{2} + 16x + 5\)}{\(3x^{2} + 8x + 13\)}{\(3x^{2} + 13\)}{\(3x^{2} + 5\)}{}{}}
\End

\Data\ID{1003032303}
\Updated{2020-07-15 03:07:21}
\Level{B}
\SubArea{Polynomials and fractions}
\LANGen{
\Question{
A car travels by \( 20\,\mathrm{km}/\mathrm{h} \) faster than a second car. The first car covers \( 260\,\mathrm{km} \) in the same time the second car covers \( 195\,\mathrm{km} \). What is the average speed of each car?}
{\( 80\,\mathrm{km}/\mathrm{h} \) and \( 60\,\mathrm{km}/\mathrm{h} \)}{\( 100\,\mathrm{km}/\mathrm{h} \) and \( 80\,\mathrm{km}/\mathrm{h} \)}{\( 90\,\mathrm{km}/\mathrm{h} \) and \( 70\,\mathrm{km}/\mathrm{h} \)}{\( 120\,\mathrm{km}/\mathrm{h} \) and \( 100\,\mathrm{km}/\mathrm{h} \)}{}{}}
\LANGpl{
\Question{
Samochód jadący z prędkością o \( 20\,\mathrm{km}/\mathrm{h} \) większą pokonał trasę \( 260\,\mathrm{km} \). Drugi samochód w tym samym czasie pokonał trasę \( 195\,\mathrm{km} \). Oblicz średnie prędkości obu samochodów.}
{\( 80\,\mathrm{km}/\mathrm{h} \) i \( 60\,\mathrm{km}/\mathrm{h} \)}{\( 100\,\mathrm{km}/\mathrm{h} \) i \( 80\,\mathrm{km}/\mathrm{h} \)}{\( 90\,\mathrm{km}/\mathrm{h} \) i \( 70\,\mathrm{km}/\mathrm{h} \)}{\( 120\,\mathrm{km}/\mathrm{h} \) i \( 100\,\mathrm{km}/\mathrm{h} \)}{}{}}
\LANGcs{
\Question{
První auto jede o \( 20\,\mathrm{km}/\mathrm{h} \) rychleji než druhé auto. První auto ujede \( 260\,\mathrm{km} \) za stejnou dobu jako druhé auto \( 195\,\mathrm{km} \). Jakou rychlostí auta jedou?}
{\( 80\,\mathrm{km}/\mathrm{h} \) a \( 60\,\mathrm{km}/\mathrm{h} \)}{\( 100\,\mathrm{km}/\mathrm{h} \) a \( 80\,\mathrm{km}/\mathrm{h} \)}{\( 90\,\mathrm{km}/\mathrm{h} \) a \( 70\,\mathrm{km}/\mathrm{h} \)}{\( 120\,\mathrm{km}/\mathrm{h} \) a \( 100\,\mathrm{km}/\mathrm{h} \)}{}{}}
\LANGsk{
\Question{
Prvé auto ide o \( 20\,\mathrm{km}/\mathrm{h} \) rýchlejšie ako druhé auto. Prvé auto prejde dráhu \( 260\,\mathrm{km} \) v tom istom čase ako druhé auto prejde dráhu \( 195\,\mathrm{km} \). Aká je priemerná rýchlosť obidvoch áut?}
{\( 80\,\mathrm{km}/\mathrm{h} \) a \( 60\,\mathrm{km}/\mathrm{h} \)}{\( 100\,\mathrm{km}/\mathrm{h} \) a \( 80\,\mathrm{km}/\mathrm{h} \)}{\( 90\,\mathrm{km}/\mathrm{h} \) a \( 70\,\mathrm{km}/\mathrm{h} \)}{\( 120\,\mathrm{km}/\mathrm{h} \) a \( 100\,\mathrm{km}/\mathrm{h} \)}{}{}}
\LANGes{
\Question{
A car travels by \( 20\,\mathrm{km}/\mathrm{h} \) faster than a second car. The first car covers \( 260\,\mathrm{km} \) in the same time the second car covers \( 195\,\mathrm{km} \). What is the average speed of each car?}
{\( 80\,\mathrm{km}/\mathrm{h} \) and \( 60\,\mathrm{km}/\mathrm{h} \)}{\( 100\,\mathrm{km}/\mathrm{h} \) and \( 80\,\mathrm{km}/\mathrm{h} \)}{\( 90\,\mathrm{km}/\mathrm{h} \) and \( 70\,\mathrm{km}/\mathrm{h} \)}{\( 120\,\mathrm{km}/\mathrm{h} \) and \( 100\,\mathrm{km}/\mathrm{h} \)}{}{}}
\End

\Data\ID{1003032401}
\Updated{2020-07-15 03:07:21}
\Level{B}
\SubArea{Polynomials and fractions}
\LANGen{
\Question{
Factoring the polynomial \( 625x^4-1 \) we get:}
{\( \left( 25x^2+1\right)(5x-1)(5x+1) \)}{\( (5x-1)(5x+1)^2 \)}{\( (5x-1)^4 \)}{\( \left( 25x^2+1\right)(5x-1)^2 \)}{}{}}
\LANGpl{
\Question{
Wielomian  \( 625x^4-1 \) jest równy:}
{\( \left( 25x^2+1\right)(5x-1)(5x+1) \)}{\( (5x-1)(5x+1)^2 \)}{\( (5x-1)^4 \)}{\( \left( 25x^2+1\right)(5x-1)^2 \)}{}{}}
\LANGcs{
\Question{
Rozkladem mnohočlenu \( 625x^4-1 \) získáme:}
{\( \left( 25x^2+1\right)(5x-1)(5x+1) \)}{\( (5x-1)(5x+1)^2 \)}{\( (5x-1)^4 \)}{\( \left( 25x^2+1\right)(5x-1)^2 \)}{}{}}
\LANGsk{
\Question{
Rozkladom polynómu \( 625x^4-1 \) na súčin dostaneme:}
{\( \left( 25x^2+1\right)(5x-1)(5x+1) \)}{\( (5x-1)(5x+1)^2 \)}{\( (5x-1)^4 \)}{\( \left( 25x^2+1\right)(5x-1)^2 \)}{}{}}
\LANGes{
\Question{
Factoring the polynomial \( 625x^4-1 \) we get:}
{\( \left( 25x^2+1\right)(5x-1)(5x+1) \)}{\( (5x-1)(5x+1)^2 \)}{\( (5x-1)^4 \)}{\( \left( 25x^2+1\right)(5x-1)^2 \)}{}{}}
\End

\Data\ID{1003032501}
\Updated{2020-07-15 03:07:21}
\Level{B}
\SubArea{Polynomials and fractions}
\LANGen{
\Question{
The product \( (x+y)\left(x^2+y^2\right)\left(x^3+y^3\right) \) equals:}
{\( x^6+x^5y+x^4y^2+2x^3y^3+x^2y^4+xy^5+y^6 \)}{\( x^6-x^5y+x^4y^2-2x^3y^3+x^2y^4-xy^5+y^6 \)}{\( (x+y)^6 \)}{\( x^6+y^6 \)}{}{}}
\LANGpl{
\Question{
Wyrażenie  \( (x+y)\left(x^2+y^2\right)\left(x^3+y^3\right) \) jest równe:}
{\( x^6+x^5y+x^4y^2+2x^3y^3+x^2y^4+xy^5+y^6 \)}{\( x^6-x^5y+x^4y^2-2x^3y^3+x^2y^4-xy^5+y^6 \)}{\( (x+y)^6 \)}{\( x^6+y^6 \)}{}{}}
\LANGcs{
\Question{
Součin \( (x+y)\left(x^2+y^2\right)\left(x^3+y^3\right) \) je roven:}
{\( x^6+x^5y+x^4y^2+2x^3y^3+x^2y^4+xy^5+y^6 \)}{\( x^6-x^5y+x^4y^2-2x^3y^3+x^2y^4-xy^5+y^6 \)}{\( (x+y)^6 \)}{\( x^6+y^6 \)}{}{}}
\LANGsk{
\Question{
Súčin polynómov \( (x+y)\left(x^2+y^2\right)\left(x^3+y^3\right) \) je rovný:}
{\( x^6+x^5y+x^4y^2+2x^3y^3+x^2y^4+xy^5+y^6 \)}{\( x^6-x^5y+x^4y^2-2x^3y^3+x^2y^4-xy^5+y^6 \)}{\( (x+y)^6 \)}{\( x^6+y^6 \)}{}{}}
\LANGes{
\Question{
The product \( (x+y)\left(x^2+y^2\right)\left(x^3+y^3\right) \) equals:}
{\( x^6+x^5y+x^4y^2+2x^3y^3+x^2y^4+xy^5+y^6 \)}{\( x^6-x^5y+x^4y^2-2x^3y^3+x^2y^4-xy^5+y^6 \)}{\( (x+y)^6 \)}{\( x^6+y^6 \)}{}{}}
\End

\Data\ID{9000039301}
\Updated{2020-07-15 03:07:34}
\Level{B}
\SubArea{Polynomials and fractions}
\LANGen{
\Question{
Given formula for an acceleration of the motion in the form 
\[ 
 a = \frac{v - v_{0}} 
 {t} ,
\]
find the initial velocity \(v_{0}\).}
{\(v_{0} = v - at\)}{\(v_{0} = vat\)}{\(v_{0} = v + at\)}{\(v_{0} = at - v\)}{}{}}
\LANGpl{
\Question{
Mając wzór na przyspieszenie  
\[ 
 a = \frac{v - v_{0}} 
 {t} ,
\]
znajdź początkową prędkość \(v_{0}\).}
{\(v_{0} = v - at\)}{\(v_{0} = vat\)}{\(v_{0} = v + at\)}{\(v_{0} = at - v\)}{}{}}
\LANGcs{
\Question{
Ze vzorce pro výpočet zrychlení rovnoměrně zrychleného pohybu
\(a = \frac{v-v_{0}} 
 {t} \) vyjádřete
počáteční rychlost \(v_{0}\).}
{\(v_{0} = v - at\)}{\(v_{0} = vat\)}{\(v_{0} = v + at\)}{\(v_{0} = at - v\)}{}{}}
\LANGsk{
\Question{
Zo vzorca zrýchlenie rovnomerne zrýchleného pohybu
\(a = \frac{v-v_{0}} 
 {t} \) vyjadrite
začiatočnú rýchlosť \(v_{0}\).}
{\(v_{0} = v - at\)}{\(v_{0} = vat\)}{\(v_{0} = v + at\)}{\(v_{0} = at - v\)}{}{}}
\LANGes{
\Question{
Given formula for an acceleration of the motion in the form 
\[ 
 a = \frac{v - v_{0}} 
 {t} ,
\]
find the initial velocity \(v_{0}\).}
{\(v_{0} = v - at\)}{\(v_{0} = vat\)}{\(v_{0} = v + at\)}{\(v_{0} = at - v\)}{}{}}
\End

\Data\ID{9000039302}
\Updated{2020-07-15 03:07:34}
\Level{B}
\SubArea{Polynomials and fractions}
\LANGen{
\Question{
Find \(N\),
the number of the turns, as a function of the other variables in the formula for the
magnetic induction of a solenoid. 
\[ 
 B =\mu \frac{NI} 
 {l}
\]}
{\(N = \frac{Bl} 
 {\mu I} \)}{\(N = \frac{Bl\mu } 
 {I} \)}{\(N = B -\mu \frac{I} 
{l} \)}{\(N = \frac{Bl} 
 {\mu } - I\)}{}{}}
\LANGpl{
\Question{
Znajdź \(N\), liczbę obrotów, jako funkcję innych zmiennych we wzorze na indukcję magnetyczną cewki.
\[ 
 B =\mu \frac{NI} 
 {l}
\]}
{\(N = \frac{Bl} 
 {\mu I} \)}{\(N = \frac{Bl\mu } 
 {I} \)}{\(N = B -\mu \frac{I} 
{l} \)}{\(N = \frac{Bl} 
 {\mu } - I\)}{}{}}
\LANGcs{
\Question{
Ze vzorce pro velikost magnetické indukce
\(B =\mu \frac{NI} 
 {l} \) vyjádřete počet
závitů cívky \(N\).}
{\(N = \frac{Bl} 
 {\mu I} \)}{\(N = \frac{Bl\mu } 
 {I} \)}{\(N = B -\mu \frac{I} 
{l} \)}{\(N = \frac{Bl} 
 {\mu } - I\)}{}{}}
\LANGsk{
\Question{
Zo vzorca pre veľkosť magnetickej indukcie
\(B =\mu \frac{NI} 
 {l} \) vyjadrite počet
závitov cievky \(N\).}
{\(N = \frac{Bl} 
 {\mu I} \)}{\(N = \frac{Bl\mu } 
 {I} \)}{\(N = B -\mu \frac{I} 
{l} \)}{\(N = \frac{Bl} 
 {\mu } - I\)}{}{}}
\LANGes{
\Question{
Find \(N\),
the number of the turns, as a function of the other variables in the formula for the
magnetic induction of a solenoid. 
\[ 
 B =\mu \frac{NI} 
 {l}
\]}
{\(N = \frac{Bl} 
 {\mu I} \)}{\(N = \frac{Bl\mu } 
 {I} \)}{\(N = B -\mu \frac{I} 
{l} \)}{\(N = \frac{Bl} 
 {\mu } - I\)}{}{}}
\End

\Data\ID{9000039303}
\Updated{2020-07-15 03:07:34}
\Level{B}
\SubArea{Polynomials and fractions}
\LANGen{
\Question{
Find the time \(t\) 
as a function of the other variables in the formula for the distance of a motion in the
form \(s = v_{0}t + s_{0}\).}
{\(t = \frac{s-s_{0}} 
 {v_{0}} \)}{\(t = \frac{s} 
{t+s_{0}} \)}{\(t = \frac{s+s_{0}} 
 {v_{0}} \)}{\(t = \frac{v_{0}} 
{s-s_{0}} \)}{}{}}
\LANGpl{
\Question{
Znajdź  \(t\)  jako funkcję innych zmiennych we wzorze na odległość  \(s = v_{0}t + s_{0}\).}
{\(t = \frac{s-s_{0}} 
 {v_{0}} \)}{\(t = \frac{s} 
{t+s_{0}} \)}{\(t = \frac{s+s_{0}} 
 {v_{0}} \)}{\(t = \frac{v_{0}} 
{s-s_{0}} \)}{}{}}
\LANGcs{
\Question{
Ze vztahu pro určení dráhy rovnoměrného přímočarého pohybu
\(s = v_{0}t + s_{0}\) vyjádřete
čas \(t\).}
{\(t = \frac{s-s_{0}} 
 {v_{0}} \)}{\(t = \frac{s} 
{t+s_{0}} \)}{\(t = \frac{s+s_{0}} 
 {v_{0}} \)}{\(t = \frac{v_{0}} 
{s-s_{0}} \)}{}{}}
\LANGsk{
\Question{
Zo vzťahu pre určenie dráhy rovnomerného priamočiareho pohybu
\(s = v_{0}t + s_{0}\) vyjadrite
čas \(t\).}
{\(t = \frac{s-s_{0}} 
 {v_{0}} \)}{\(t = \frac{s} 
{t+s_{0}} \)}{\(t = \frac{s+s_{0}} 
 {v_{0}} \)}{\(t = \frac{v_{0}} 
{s-s_{0}} \)}{}{}}
\LANGes{
\Question{
Find the time \(t\) 
as a function of the other variables in the formula for the distance of a motion in the
form \(s = v_{0}t + s_{0}\).}
{\(t = \frac{s-s_{0}} 
 {v_{0}} \)}{\(t = \frac{s} 
{t+s_{0}} \)}{\(t = \frac{s+s_{0}} 
 {v_{0}} \)}{\(t = \frac{v_{0}} 
{s-s_{0}} \)}{}{}}
\End

\Data\ID{9000039304}
\Updated{2020-07-15 03:07:34}
\Level{B}
\SubArea{Polynomials and fractions}
\LANGen{
\Question{
Find the focus length \(f\) 
as a function of the other variables from the following equation relating this distance with object
and image distances \(a\) 
and \(a'\).
\[ 
 \frac{1} 
{f} = \frac{1} 
{a} + \frac{1} 
{a'}
\]}
{\(f = \frac{aa'} 
{a+a'}\)}{\(f = \frac{a-a'} 
{a+a'}\)}{\(f = a + a'\)}{\(f = \frac{a} 
{a'}\)}{}{}}
\LANGpl{
\Question{
Znajdź długość ogniskowej  \(f\)  
jako funkcję innych zmiennych z następującego równania wiążącego tę odległość z odległością przedmiotu  \(a\) 
i obrazu \(a'\).
\[ 
 \frac{1} 
{f} = \frac{1} 
{a} + \frac{1} 
{a'}
\]}
{\(f = \frac{aa'} 
{a+a'}\)}{\(f = \frac{a-a'} 
{a+a'}\)}{\(f = a + a'\)}{\(f = \frac{a} 
{a'}\)}{}{}}
\LANGcs{
\Question{
Ze zobrazovací rovnice kulového zrcadla
\(\frac{1}
{f} = \frac{1} 
{a} + \frac{1} 
{a'}\) vyjádřete ohniskovou
vzdálenost \(f\).}
{\(f = \frac{aa'} 
{a+a'}\)}{\(f = \frac{a-a'} 
{a+a'}\)}{\(f = a + a'\)}{\(f = \frac{a} 
{a'}\)}{}{}}
\LANGsk{
\Question{
Zo zobrazovacej rovnice guľového zrkadla
\(\frac{1}
{f} = \frac{1} 
{a} + \frac{1} 
{a'}\) vyjadrite ohniskovú
vzdialenosť \(f\).}
{\(f = \frac{aa'} 
{a+a'}\)}{\(f = \frac{a-a'} 
{a+a'}\)}{\(f = a + a'\)}{\(f = \frac{a} 
{a'}\)}{}{}}
\LANGes{
\Question{
Find the focus length \(f\) 
as a function of the other variables from the following equation relating this distance with object
and image distances \(a\) 
and \(a'\).
\[ 
 \frac{1} 
{f} = \frac{1} 
{a} + \frac{1} 
{a'}
\]}
{\(f = \frac{aa'} 
{a+a'}\)}{\(f = \frac{a-a'} 
{a+a'}\)}{\(f = a + a'\)}{\(f = \frac{a} 
{a'}\)}{}{}}
\End

\Data\ID{9000039305}
\Updated{2020-07-15 03:07:34}
\Level{B}
\SubArea{Polynomials and fractions}
\LANGen{
\Question{
Find \(m_{1}\) 
as a function of the other variables from the following mixing equation.
\[ 
 w_{1}m_{1} + w_{2}m_{2} = w_{3}m_{3}
\]}
{\(m_{1} = \frac{w_{3}m_{3}-w_{2}m_{2}} 
 {w_{1}} \)}{\(m_{1} = \frac{w_{3}m_{3}w_{2}m_{2}} 
 {w_{1}} \)}{\(m_{1} = \frac{w_{3}m_{3}+w_{2}m_{2}} 
 {w_{1}} \)}{\(m_{1} = \frac{w_{2}m_{2}-w_{3}m_{3}} 
 {w_{1}} \)}{}{}}
\LANGpl{
\Question{
Znajdź  \(m_{1}\) jako funkcję innych zmiennych z następującego równania.
\[ 
 w_{1}m_{1} + w_{2}m_{2} = w_{3}m_{3}
\]}
{\(m_{1} = \frac{w_{3}m_{3}-w_{2}m_{2}} 
 {w_{1}} \)}{\(m_{1} = \frac{w_{3}m_{3}w_{2}m_{2}} 
 {w_{1}} \)}{\(m_{1} = \frac{w_{3}m_{3}+w_{2}m_{2}} 
 {w_{1}} \)}{\(m_{1} = \frac{w_{2}m_{2}-w_{3}m_{3}} 
 {w_{1}} \)}{}{}}
\LANGcs{
\Question{
Ze směšovací rovnice \(w_{1}m_{1} + w_{2}m_{2} = w_{3}m_{3}\) 
vyjádřete hmotnost \(m_{1}\).}
{\(m_{1} = \frac{w_{3}m_{3}-w_{2}m_{2}} 
 {w_{1}} \)}{\(m_{1} = \frac{w_{3}m_{3}w_{2}m_{2}} 
 {w_{1}} \)}{\(m_{1} = \frac{w_{3}m_{3}+w_{2}m_{2}} 
 {w_{1}} \)}{\(m_{1} = \frac{w_{2}m_{2}-w_{3}m_{3}} 
 {w_{1}} \)}{}{}}
\LANGsk{
\Question{
Zo zmiešavacej rovnice \(w_{1}m_{1} + w_{2}m_{2} = w_{3}m_{3}\) 
vyjadrite hmotnosť \(m_{1}\).}
{\(m_{1} = \frac{w_{3}m_{3}-w_{2}m_{2}} 
 {w_{1}} \)}{\(m_{1} = \frac{w_{3}m_{3}w_{2}m_{2}} 
 {w_{1}} \)}{\(m_{1} = \frac{w_{3}m_{3}+w_{2}m_{2}} 
 {w_{1}} \)}{\(m_{1} = \frac{w_{2}m_{2}-w_{3}m_{3}} 
 {w_{1}} \)}{}{}}
\LANGes{
\Question{
Find \(m_{1}\) 
as a function of the other variables from the following mixing equation.
\[ 
 w_{1}m_{1} + w_{2}m_{2} = w_{3}m_{3}
\]}
{\(m_{1} = \frac{w_{3}m_{3}-w_{2}m_{2}} 
 {w_{1}} \)}{\(m_{1} = \frac{w_{3}m_{3}w_{2}m_{2}} 
 {w_{1}} \)}{\(m_{1} = \frac{w_{3}m_{3}+w_{2}m_{2}} 
 {w_{1}} \)}{\(m_{1} = \frac{w_{2}m_{2}-w_{3}m_{3}} 
 {w_{1}} \)}{}{}}
\End

\Data\ID{9000079202}
\Updated{2020-07-15 03:07:37}
\Level{B}
\SubArea{Polynomials and fractions}
\LANGen{
\Question{
Find the set \(M\) 
of all the real \(x\) 
for which the following expression is not a well defined number. 
\[ 
 \frac{x - 4} 
{x^{3} - 16x}
\]}
{\(M = \{ - 4;0;4\}\)}{\(M = \{ - 4;4\}\)}{\(M = \{0;4\}\)}{\(M = \{0\}\)}{}{}}
\LANGpl{
\Question{
Znajdź zbiór  \(M\) wszystkich rzeczywistych wartości \(x\), dla których podane wyrażenie nie jest określone.
\[ 
 \frac{x - 4} 
{x^{3} - 16x}
\]}
{\(M = \{ - 4;0;4\}\)}{\(M = \{ - 4;4\}\)}{\(M = \{0;4\}\)}{\(M = \{0\}\)}{}{}}
\LANGcs{
\Question{
Určete množinu všech hodnot \(x\),
pro které není výraz \(\frac{x-4}
{x^{3}-16x}\) 
definován.}
{\(M = \{ - 4;0;4\}\)}{\(M = \{ - 4;4\}\)}{\(M = \{0;4\}\)}{\(M = \{0\}\)}{}{}}
\LANGsk{
\Question{
Určte množinu všetkých hodnôt \(x\),
pre ktoré nie je výraz \(\frac{x-4}
{x^{3}-16x}\) 
definovaný.}
{\(M = \{ - 4;0;4\}\)}{\(M = \{ - 4;4\}\)}{\(M = \{0;4\}\)}{\(M = \{0\}\)}{}{}}
\LANGes{
\Question{
Find the set \(M\) 
of all the real \(x\) 
for which the following expression is not a well defined number. 
\[ 
 \frac{x - 4} 
{x^{3} - 16x}
\]}
{\(M = \{ - 4;0;4\}\)}{\(M = \{ - 4;4\}\)}{\(M = \{0;4\}\)}{\(M = \{0\}\)}{}{}}
\End

\Data\ID{9000079204}
\Updated{2020-07-15 03:07:37}
\Level{B}
\SubArea{Polynomials and fractions}
\LANGen{
\Question{
Find the domain of the following expression. 
\[ 
 \frac{x^{2} - x} 
 {x + 1} : \frac{x^{2} - 1} 
{x^{2} + 2x + 1}
\]}
{\(\mathbb{R}\setminus \{ - 1;1\}\)}{\(\mathbb{R}\setminus \{ - 1;0;1\}\)}{\(\mathbb{R}\setminus \{ - 1\}\)}{\(\mathbb{R}\setminus \{ - 1;0\}\)}{}{}}
\LANGpl{
\Question{
Znajdź dziedzinę podanego wyrażenia:
\[ 
 \frac{x^{2} - x} 
 {x + 1} : \frac{x^{2} - 1} 
{x^{2} + 2x + 1}
\]}
{\(\mathbb{R}\setminus \{ - 1;1\}\)}{\(\mathbb{R}\setminus \{ - 1;0;1\}\)}{\(\mathbb{R}\setminus \{ - 1\}\)}{\(\mathbb{R}\setminus \{ - 1;0\}\)}{}{}}
\LANGcs{
\Question{
Určete množinu všech hodnot \(x\),
pro které má výraz \(\frac{x^{2}-x}
 {x+1} : \frac{x^{2}-1} 
{x^{2}+2x+1}\) 
smysl.}
{\(\mathbb{R}\setminus \{ - 1;1\}\)}{\(\mathbb{R}\setminus \{ - 1;0;1\}\)}{\(\mathbb{R}\setminus \{ - 1\}\)}{\(\mathbb{R}\setminus \{ - 1;0\}\)}{}{}}
\LANGsk{
\Question{
Určte množinu všetkých hodnôt \(x\),
pre ktoré má výraz \(\frac{x^{2}-x}
 {x+1} : \frac{x^{2}-1} 
{x^{2}+2x+1}\) 
zmysel.}
{\(\mathbb{R}\setminus \{ - 1;1\}\)}{\(\mathbb{R}\setminus \{ - 1;0;1\}\)}{\(\mathbb{R}\setminus \{ - 1\}\)}{\(\mathbb{R}\setminus \{ - 1;0\}\)}{}{}}
\LANGes{
\Question{
Find the domain of the following expression. 
\[ 
 \frac{x^{2} - x} 
 {x + 1} : \frac{x^{2} - 1} 
{x^{2} + 2x + 1}
\]}
{\(\mathbb{R}\setminus \{ - 1;1\}\)}{\(\mathbb{R}\setminus \{ - 1;0;1\}\)}{\(\mathbb{R}\setminus \{ - 1\}\)}{\(\mathbb{R}\setminus \{ - 1;0\}\)}{}{}}
\End

\Data\ID{9000079208}
\Updated{2020-07-15 03:07:37}
\Level{B}
\SubArea{Polynomials and fractions}
\LANGen{
\Question{
Assuming \(x\neq 0\) 
and \(y\neq 0\),
simplify the following expression. 
\[ 
 \left (\frac{x^{-2}y^{2}}
{x^{0}y^{-8}}\right )^{-2} : \frac{x^{2}} 
{x^{-4}y^{7}}
\]}
{\(\frac{1}
{x^{2}y^{13}} \)}{\(\frac{y^{13}} 
 {x^{2}} \)}{\(\frac{y^{15}} 
 {x^{6}} \)}{\(\frac{x^{4}} 
{y^{27}} \)}{}{}}
\LANGpl{
\Question{
Zakładając, że \(x\neq 0\) 
i \(y\neq 0\),
uprość podane wyrażenie:
\[ 
 \left (\frac{x^{-2}y^{2}}
{x^{0}y^{-8}}\right )^{-2} : \frac{x^{2}} 
{x^{-4}y^{7}}
\]}
{\(\frac{1}
{x^{2}y^{13}} \)}{\(\frac{y^{13}} 
 {x^{2}} \)}{\(\frac{y^{15}} 
 {x^{6}} \)}{\(\frac{x^{4}} 
{y^{27}} \)}{}{}}
\LANGcs{
\Question{
Upravte výraz \(\left (\frac{x^{-2}y^{2}} 
{x^{0}y^{-8}} \right )^{-2} : \frac{x^{2}} 
{x^{-4}y^{7}} \) 
za předpokladu, že \(x\neq 0\) a
\(y\neq 0\).}
{\(\frac{1}
{x^{2}y^{13}} \)}{\(\frac{y^{13}} 
 {x^{2}} \)}{\(\frac{y^{15}} 
 {x^{6}} \)}{\(\frac{x^{4}} 
{y^{27}} \)}{}{}}
\LANGsk{
\Question{
Zjednodušte výraz \(\left (\frac{x^{-2}y^{2}} 
{x^{0}y^{-8}} \right )^{-2} : \frac{x^{2}} 
{x^{-4}y^{7}} \) 
za predpokladu, že \(x\neq 0\) a
\(y\neq 0\).}
{\(\frac{1}
{x^{2}y^{13}} \)}{\(\frac{y^{13}} 
 {x^{2}} \)}{\(\frac{y^{15}} 
 {x^{6}} \)}{\(\frac{x^{4}} 
{y^{27}} \)}{}{}}
\LANGes{
\Question{
Assuming \(x\neq 0\) 
and \(y\neq 0\),
simplify the following expression. 
\[ 
 \left (\frac{x^{-2}y^{2}}
{x^{0}y^{-8}}\right )^{-2} : \frac{x^{2}} 
{x^{-4}y^{7}}
\]}
{\(\frac{1}
{x^{2}y^{13}} \)}{\(\frac{y^{13}} 
 {x^{2}} \)}{\(\frac{y^{15}} 
 {x^{6}} \)}{\(\frac{x^{4}} 
{y^{27}} \)}{}{}}
\End

\Data\ID{9000083601}
\Updated{2020-07-15 03:07:37}
\Level{B}
\SubArea{Polynomials and fractions}
\LANGen{
\Question{
Find the conditions on \(x\) 
under which the expression 
\[ 
 \frac{\frac{x-y}
{x+y} -\frac{x+y} 
{x-y}} 
 { \frac{xy}
{x^{2}-y^{2}} }
\]
is well-defined.}
{\(x\neq 0,\; y\neq 0,\; x\neq \pm y\)}{\(x\neq - y\)}{\(x\neq \pm y\)}{\(x\neq 0,\; y\neq 0\)}{}{}}
\LANGpl{
\Question{
Określ warunki dla  \(x\), w których wyrażenie  
\[ 
 \frac{\frac{x-y}
{x+y} -\frac{x+y} 
{x-y}} 
 { \frac{xy}
{x^{2}-y^{2}} }
\]
ma sens.}
{\(x\neq 0,\; y\neq 0,\; x\neq \pm y\)}{\(x\neq - y\)}{\(x\neq \pm y\)}{\(x\neq 0,\; y\neq 0\)}{}{}}
\LANGcs{
\Question{
Určete, za jakých podmínek má výraz
\(\frac{\frac{x-y}
{x+y}-\frac{x+y}
{x-y}}
 { \frac{xy}
{x^{2}-y^{2}} } \) 
smysl.}
{\(x\neq 0,\; y\neq 0,\; x\neq \pm y\)}{\(x\neq - y\)}{\(x\neq \pm y\)}{\(x\neq 0,\; y\neq 0\)}{}{}}
\LANGsk{
\Question{
Určte podmienky daného výrazu
\(\frac{\frac{x-y}
{x+y}-\frac{x+y}
{x-y}}
 { \frac{xy}
{x^{2}-y^{2}} } \).}
{\(x\neq 0,\; y\neq 0,\; x\neq \pm y\)}{\(x\neq - y\)}{\(x\neq \pm y\)}{\(x\neq 0,\; y\neq 0\)}{}{}}
\LANGes{
\Question{
Find the conditions on \(x\) 
under which the expression 
\[ 
 \frac{\frac{x-y}
{x+y} -\frac{x+y} 
{x-y}} 
 { \frac{xy}
{x^{2}-y^{2}} }
\]
is well-defined.}
{\(x\neq 0,\; y\neq 0,\; x\neq \pm y\)}{\(x\neq - y\)}{\(x\neq \pm y\)}{\(x\neq 0,\; y\neq 0\)}{}{}}
\End

\Data\ID{9000083607}
\Updated{2020-07-15 03:07:37}
\Level{B}
\SubArea{Polynomials and fractions}
\LANGen{
\Question{
Assuming \(x\neq 0\),
\(x\neq \pm 1\),
\(y\neq 0\),
simplify the expression. 
\[ 
 \left [\left ( \frac{x}
{x + 1}\right )^{2} : \left (\frac{x - 1}
 {y} \right )^{2}\right ] : \frac{2xy} 
{x^{2} - 1}
\]}
{\(\frac{xy}
{2\left (x^{2}-1\right )}\)}{\(4\)}{\(\frac{x^{2}-1}
 {4} \)}{\(\frac{x-1}
 {4} \)}{}{}}
\LANGpl{
\Question{
Zakładając, że  \(x\neq 0\),
\(x\neq \pm 1\),
\(y\neq 0\),
uprość wyrażenie.
\[ 
 \left [\left ( \frac{x}
{x + 1}\right )^{2} : \left (\frac{x - 1}
 {y} \right )^{2}\right ] : \frac{2xy} 
{x^{2} - 1}
\]}
{\(\frac{xy}
{2\left (x^{2}-1\right )}\)}{\(4\)}{\(\frac{x^{2}-1}
 {4} \)}{\(\frac{x-1}
 {4} \)}{}{}}
\LANGcs{
\Question{
Za předpokladu, že \(x\neq 0\),
\(x\neq \pm 1\),
\(y\neq 0\), zjednodušte výraz. \[\left [\left ( \frac{x}
{x+1}\right )^{2} : \left (\frac{x-1}
 {y} \right )^{2}\right ] : \frac{2xy} 
{x^{2}-1}\]}
{\(\frac{xy}
{2\left (x^{2}-1\right )}\)}{\(4\)}{\(\frac{x^{2}-1}
 {4} \)}{\(\frac{x-1}
 {4} \)}{}{}}
\LANGsk{
\Question{
Za predpokladu, že \(x\neq 0\),
\(x\neq \pm 1\),
\(y\neq 0\), zjednodušte výraz. \[\left [\left ( \frac{x}
{x+1}\right )^{2} : \left (\frac{x-1}
 {y} \right )^{2}\right ] : \frac{2xy} 
{x^{2}-1}\]}
{\(\frac{xy}
{2\left (x^{2}-1\right )}\)}{\(4\)}{\(\frac{x^{2}-1}
 {4} \)}{\(\frac{x-1}
 {4} \)}{}{}}
\LANGes{
\Question{
Assuming \(x\neq 0\),
\(x\neq \pm 1\),
\(y\neq 0\),
simplify the expression. 
\[ 
 \left [\left ( \frac{x}
{x + 1}\right )^{2} : \left (\frac{x - 1}
 {y} \right )^{2}\right ] : \frac{2xy} 
{x^{2} - 1}
\]}
{\(\frac{xy}
{2\left (x^{2}-1\right )}\)}{\(4\)}{\(\frac{x^{2}-1}
 {4} \)}{\(\frac{x-1}
 {4} \)}{}{}}
\End

\Data\ID{9000083608}
\Updated{2021-01-04 12:01:59}
\Level{B}
\SubArea{Polynomials and fractions}
\LANGen{
\Question{
Assuming \(xy\neq - 1\),
simplify the expression: 
\[ 
 \frac{ \frac{x-y}
{1+xy} + y} 
{1 -\frac{y(x-y)} 
 {1+xy} }
\]}
{\(x\)}{\(\frac{x(1+y^{2})}
 {1-y^{2}} \)}{\(x - 1\)}{\(x(1 + y^{2})\)}{}{}}
\LANGpl{
\Question{
Zakładając, że  \(xy\neq - 1\),
uprość wyrażenie:
\[ 
 \frac{ \frac{x-y}
{1+xy} + y} 
{1 -\frac{y(x-y)} 
 {1+xy} }
\]}
{\(x\)}{\(\frac{x(1+y^{2})}
 {1-y^{2}} \)}{\(x - 1\)}{\(x(1 + y^{2})\)}{}{}}
\LANGcs{
\Question{
Za předpokladu, že \(xy\neq - 1\), zjednodušte výraz: \[\frac{ \frac{x-y}
{1+xy}+y}
{1-\frac{y(x-y)} 
 {1+xy} }\]}
{\(x\)}{\(\frac{x(1+y^{2})}
 {1-y^{2}} \)}{\(x - 1\)}{\(x(1 + y^{2})\)}{}{}}
\LANGsk{
\Question{
Za predpokladu, že \(xy\neq - 1\), zjednodušte výraz: \[\frac{ \frac{x-y}
{1+xy}+y}
{1-\frac{y(x-y)} 
 {1+xy} }\]}
{\(x\)}{\(\frac{x(1+y^{2})}
 {1-y^{2}} \)}{\(x - 1\)}{\(x(1 + y^{2})\)}{}{}}
\LANGes{
\Question{
Assuming \(xy\neq - 1\),
simplify the expression: 
\[ 
 \frac{ \frac{x-y}
{1+xy} + y} 
{1 -\frac{y(x-y)} 
 {1+xy} }
\]}
{\(x\)}{\(\frac{x(1+y^{2})}
 {1-y^{2}} \)}{\(x - 1\)}{\(x(1 + y^{2})\)}{}{}}
\End

\Data\ID{9000083609}
\Updated{2020-07-15 03:07:37}
\Level{B}
\SubArea{Polynomials and fractions}
\LANGen{
\Question{
Assuming \(x\neq 0\),
\(x\neq \pm y\),
\(y\neq 0\),
simplify the expression. 
\[ 
 \frac{\frac{x^{2}+y^{2}} 
 {x} - 2y} 
{\left ( \frac{1}
{y^{2}} - \frac{1} 
{x^{2}} \right )\cdot \frac{xy}
{x+y}}
\]}
{\(y(x - y)\)}{\(\frac{x-y}
 {y} \)}{\(x(x - y)\)}{\(\frac{x-y}
 {x} \)}{}{}}
\LANGpl{
\Question{
Zakładając, że  \(x\neq 0\),
\(x\neq \pm y\),
\(y\neq 0\),
uprość wyrażenie.
\[ 
 \frac{\frac{x^{2}+y^{2}} 
 {x} - 2y} 
{\left ( \frac{1}
{y^{2}} - \frac{1} 
{x^{2}} \right )\cdot \frac{xy}
{x+y}}
\]}
{\(y(x - y)\)}{\(\frac{x-y}
 {y} \)}{\(x(x - y)\)}{\(\frac{x-y}
 {x} \)}{}{}}
\LANGcs{
\Question{
Za předpokladu, že \(x\neq 0\),
\(x\neq \pm y\),
\(y\neq 0\), zjednodušte výraz. \[\frac{\frac{x^{2}+y^{2}} 
 {x} -2y}
{\left ( \frac{1} 
{y^{2}} - \frac{1} 
{x^{2}} \right )\cdot \frac{xy}
{x+y}}\]}
{\(y(x - y)\)}{\(\frac{x-y}
 {y} \)}{\(x(x - y)\)}{\(\frac{x-y}
 {x} \)}{}{}}
\LANGsk{
\Question{
Za predpokladu, že \(x\neq 0\),
\(x\neq \pm y\),
\(y\neq 0\), zjednodušte výraz. \[\frac{\frac{x^{2}+y^{2}} 
 {x} -2y}
{\left ( \frac{1} 
{y^{2}} - \frac{1} 
{x^{2}} \right )\cdot \frac{xy}
{x+y}}\]}
{\(y(x - y)\)}{\(\frac{x-y}
 {y} \)}{\(x(x - y)\)}{\(\frac{x-y}
 {x} \)}{}{}}
\LANGes{
\Question{
Assuming \(x\neq 0\),
\(x\neq \pm y\),
\(y\neq 0\),
simplify the expression. 
\[ 
 \frac{\frac{x^{2}+y^{2}} 
 {x} - 2y} 
{\left ( \frac{1}
{y^{2}} - \frac{1} 
{x^{2}} \right )\cdot \frac{xy}
{x+y}}
\]}
{\(y(x - y)\)}{\(\frac{x-y}
 {y} \)}{\(x(x - y)\)}{\(\frac{x-y}
 {x} \)}{}{}}
\End

\Data\ID{9000083610}
\Updated{2021-01-04 12:01:35}
\Level{B}
\SubArea{Polynomials and fractions}
\LANGen{
\Question{
Assuming \(x\neq \pm y\) 
and \(y\neq 2x\),
simplify the expression: 
\[ 
 \left ( \frac{2x}
{x + y} + \frac{y} 
{x - y} - \frac{y^{2}} 
{x^{2} - y^{2}}\right ) : \left ( \frac{1}
{x + y} + \frac{x} 
{x^{2} - y^{2}}\right )
\]}
{\(x\)}{\(2x - y\)}{\(\frac{x}
{2x-y}\)}{\(1\)}{}{}}
\LANGpl{
\Question{
Zakładając, że \(x\neq \pm y\) 
i \(y\neq 2x\),
uprość wyrażenie:
\[ 
 \left ( \frac{2x}
{x + y} + \frac{y} 
{x - y} - \frac{y^{2}} 
{x^{2} - y^{2}}\right ) : \left ( \frac{1}
{x + y} + \frac{x} 
{x^{2} - y^{2}}\right )
\]}
{\(x\)}{\(2x - y\)}{\(\frac{x}
{2x-y}\)}{\(1\)}{}{}}
\LANGcs{
\Question{
Za předpokladu, že \(x\neq \pm y\),
\(y\neq 2x\), zjednodušte výraz: \[\left ( \frac{2x}
{x+y} + \frac{y} 
{x-y} - \frac{y^{2}} 
{x^{2}-y^{2}} \right ) : \left ( \frac{1}
{x+y} + \frac{x} 
{x^{2}-y^{2}} \right )\]}
{\(x\)}{\(2x - y\)}{\(\frac{x}
{2x-y}\)}{\(1\)}{}{}}
\LANGsk{
\Question{
Za predpokladu, že \(x\neq \pm y\),
\(y\neq 2x\), zjednodušte výraz: \[\left ( \frac{2x}
{x+y} + \frac{y} 
{x-y} - \frac{y^{2}} 
{x^{2}-y^{2}} \right ) : \left ( \frac{1}
{x+y} + \frac{x} 
{x^{2}-y^{2}} \right )\]}
{\(x\)}{\(2x - y\)}{\(\frac{x}
{2x-y}\)}{\(1\)}{}{}}
\LANGes{
\Question{
Assuming \(x\neq \pm y\) 
and \(y\neq 2x\),
simplify the expression: 
\[ 
 \left ( \frac{2x}
{x + y} + \frac{y} 
{x - y} - \frac{y^{2}} 
{x^{2} - y^{2}}\right ) : \left ( \frac{1}
{x + y} + \frac{x} 
{x^{2} - y^{2}}\right )
\]}
{\(x\)}{\(2x - y\)}{\(\frac{x}
{2x-y}\)}{\(1\)}{}{}}
\End

\Data\ID{9000088801}
\Updated{2020-07-15 03:07:37}
\Level{B}
\SubArea{Polynomials and fractions}
\LANGen{
\Question{
Find the domain of the following expression. 
\[ 
 \frac{2x + 1} 
{6x^{2} + 3x}
\]}
{\(\mathbb{R}\setminus\left\{0,-\frac{1}{2}\right\}\)}{\(\mathbb{R}\setminus\{0,- 2\}\)}{\(\mathbb{R}\setminus\{0\}\)}{\(\mathbb{R}\setminus\left\{0,\frac{1}{2}\right\}\)}{}{}}
\LANGpl{
\Question{
Znajdź dziedzinę podanego wyrażenia.
\[ 
 \frac{2x + 1} 
{6x^{2} + 3x}
\]}
{\(\mathbb{R}\setminus\left\{0,-\frac{1}{2}\right\} \)}{\(\mathbb{R}\setminus\{ 0, - 2\} \)}{\(\mathbb{R}\setminus\{0\}\)}{\(\mathbb{R}\setminus\left\{0,\frac{1}{2}\right\} \)}{}{}}
\LANGcs{
\Question{
Určete definiční obor daného výrazu.
\[\frac{2x+1}
{6x^{2}+3x}\]}
{\(\mathbb{R}\setminus\left\{0,-\frac{1}{2}\right\} \)}{\(\mathbb{R}\setminus\{0,- 2\}\)}{\(\mathbb{R}\setminus\{0\}\)}{\(\mathbb{R}\setminus\left\{0,\frac{1}{2}\right\}\)}{}{}}
\LANGsk{
\Question{
Určte definičný obor daného výrazu.
\[ 
 \frac{2x + 1} 
{6x^{2} + 3x}
\]}
{\(\mathbb{R}\setminus\left\{0,-\frac{1}{2}\right\} \)}{\(\mathbb{R}\setminus\{0,- 2\}\)}{\(\mathbb{R}\setminus\{0\}\)}{\(\mathbb{R}\setminus\left\{0,\frac{1}{2}\right\}\)}{}{}}
\LANGes{
\Question{
Find the domain of the following expression. 
\[ 
 \frac{2x + 1} 
{6x^{2} + 3x}
\]}
{\(\mathbb{R}\setminus\left\{0,-\frac{1}{2}\right\}\)}{\(\mathbb{R}\setminus\{0,- 2\}\)}{\(\mathbb{R}\setminus\{0\}\)}{\(\mathbb{R}\setminus\left\{0,\frac{1}{2}\right\}\)}{}{}}
\End

\Data\ID{9000088802}
\Updated{2020-07-15 03:07:37}
\Level{B}
\SubArea{Polynomials and fractions}
\LANGen{
\Question{
Find the domain of the following expression. 
\[ 
 \frac{a} 
{a^{2} + 9}\cdot \frac{a^{2} - 9}
{a^{2} + 3a}
\]}
{\(\mathbb{R}\setminus\{0,- 3\}\)}{\(\mathbb{R}\setminus\{3,- 3\}\)}{\(\mathbb{R}\setminus\{0,3\}\)}{\(\mathbb{R}\setminus\{- 3\}\)}{}{}}
\LANGpl{
\Question{
Znajdź dziedzinę podanego wyrażenia.
\[ 
 \frac{a} 
{a^{2} + 9}\cdot \frac{a^{2} - 9}
{a^{2} + 3a}
\]}
{\(\mathbb{R}\setminus\{0,- 3\}\)}{\(\mathbb{R}\setminus\{3,- 3\}\)}{\(\mathbb{R}\setminus\{0,3\}\)}{\(\mathbb{R}\setminus\{- 3\}\)}{}{}}
\LANGcs{
\Question{
Určete definiční obor daného výrazu.
\[\frac{a}
{a^{2}+9}\cdot \frac{a^{2}-9}
{a^{2}+3a}\]}
{\(\mathbb{R}\setminus\{0,- 3\}\)}{\(\mathbb{R}\setminus\{3,- 3\}\)}{\(\mathbb{R}\setminus\{0,3\}\)}{\(\mathbb{R}\setminus\{- 3\}\)}{}{}}
\LANGsk{
\Question{
Určte definičný obor daného výrazu. 
\[ 
 \frac{a} 
{a^{2} + 9}\cdot \frac{a^{2} - 9}
{a^{2} + 3a}
\]}
{\(\mathbb{R}\setminus\{0,- 3\}\)}{\(\mathbb{R}\setminus\{3,- 3\}\)}{\(\mathbb{R}\setminus\{0,3\}\)}{\(\mathbb{R}\setminus\{- 3\}\)}{}{}}
\LANGes{
\Question{
Find the domain of the following expression. 
\[ 
 \frac{a} 
{a^{2} + 9}\cdot \frac{a^{2} - 9}
{a^{2} + 3a}
\]}
{\(\mathbb{R}\setminus\{0,- 3\}\)}{\(\mathbb{R}\setminus\{3,- 3\}\)}{\(\mathbb{R}\setminus\{0,3\}\)}{\(\mathbb{R}\setminus\{- 3\}\)}{}{}}
\End

\Data\ID{9000088806}
\Updated{2020-07-15 03:07:37}
\Level{B}
\SubArea{Polynomials and fractions}
\LANGen{
\Question{
Identify a correct expression which should be placed to the starred position in the
following expression to obtain a valid statement. 
\[ 
 \frac{mn} 
{m^{2} + 2mn + n^{2}} = \frac{*} 
{2m(m + n)^{3}}
\]}
{\(2m^{2}n(m + n)\)}{\(2mn(m + n)\)}{\(2m(m + n)\)}{\(2m(m + n)^{2}\)}{}{}}
\LANGpl{
\Question{
Określ poprawne wyrażenie, które należy wstawić w miejsce gwiazdki i dla którego wyrażenie ma sens.
\[ 
 \frac{mn} 
{m^{2} + 2mn + n^{2}} = \frac{*} 
{2m(m + n)^{3}}
\]}
{\(2m^{2}n(m + n)\)}{\(2mn(m + n)\)}{\(2m(m + n)\)}{\(2m(m + n)^{2}\)}{}{}}
\LANGcs{
\Question{
Na místo označené hvězdičkou doplňte takový výraz, aby v případě
nenulových jmenovatelů platila následující rovnost výrazů.
\[ 
 \frac{mn} 
{m^{2} + 2mn + n^{2}} = \frac{*} 
{2m(m + n)^{3}}
\]}
{\(2m^{2}n(m + n)\)}{\(2mn(m + n)\)}{\(2m(m + n)\)}{\(2m(m + n)^{2}\)}{}{}}
\LANGsk{
\Question{
Na miesto označené hviezdičkou doplňte taký výraz, aby v prípade
nenulových menovateľov platila následujúca rovnosť výrazov.
\[ 
 \frac{mn} 
{m^{2} + 2mn + n^{2}} = \frac{*} 
{2m(m + n)^{3}}
\]}
{\(2m^{2}n(m + n)\)}{\(2mn(m + n)\)}{\(2m(m + n)\)}{\(2m(m + n)^{2}\)}{}{}}
\LANGes{
\Question{
Identify a correct expression which should be placed to the starred position in the
following expression to obtain a valid statement. 
\[ 
 \frac{mn} 
{m^{2} + 2mn + n^{2}} = \frac{*} 
{2m(m + n)^{3}}
\]}
{\(2m^{2}n(m + n)\)}{\(2mn(m + n)\)}{\(2m(m + n)\)}{\(2m(m + n)^{2}\)}{}{}}
\End

\Data\ID{9000088807}
\Updated{2021-01-04 03:01:07}
\Level{B}
\SubArea{Polynomials and fractions}
\LANGen{
\Question{
Suppose we are given the following equality of two fractions with nonzero denominators. From the given expressions, choose the one that by substituting to the starred position makes the equality true.
\[ 
 \frac{3 - 2x} 
 {x - 2} = \frac{3(4x^{2} - 12x + 9)} 
 {*}
\]}
{\((3x - 6)(3 - 2x)\)}{\((x - 2)(2x - 3)\)}{\((x - 2)(9 - 4x)\)}{\((3x - 6)(2x - 3)\)}{}{}}
\LANGpl{
\Question{
Określ poprawne wyrażenie, które należy wstawić w miejsce gwiazdki i dla którego wyrażenie ma sens.
\[ 
 \frac{3 - 2x} 
 {x - 2} = \frac{3(4x^{2} - 12x + 9)} 
 {*}
\]}
{\((3x - 6)(3 - 2x)\)}{\((x - 2)(2x - 3)\)}{\((x - 2)(9 - 4x)\)}{\((3x - 6)(2x - 3)\)}{}{}}
\LANGcs{
\Question{
Na místo označené hvězdičkou doplňte takový výraz, aby v případě
nenulových jmenovatelů platila následující rovnost výrazů.
\[ 
 \frac{3 - 2x} 
 {x - 2} = \frac{3(4x^{2} - 12x + 9)} 
 {*}
\]}
{\((3x - 6)(3 - 2x)\)}{\((x - 2)(2x - 3)\)}{\((x - 2)(9 - 4x)\)}{\((3x - 6)(2x - 3)\)}{}{}}
\LANGsk{
\Question{
Na miesto označené hviezdičkou doplňte taký výraz, aby v prípade nenulových menovateľov platila následujúca rovnosť výrazov.
\[ 
 \frac{3 - 2x} 
 {x - 2} = \frac{3(4x^{2} - 12x + 9)} 
 {*}
\]}
{\((3x - 6)(3 - 2x)\)}{\((x - 2)(2x - 3)\)}{\((x - 2)(9 - 4x)\)}{\((3x - 6)(2x - 3)\)}{}{}}
\LANGes{
\Question{
Suppose we are given the following equality of two fractions with nonzero denominators. From the given expressions, choose the one that by substituting to the starred position makes the equality true.
\[ 
 \frac{3 - 2x} 
 {x - 2} = \frac{3(4x^{2} - 12x + 9)} 
 {*}
\]}
{\((3x - 6)(3 - 2x)\)}{\((x - 2)(2x - 3)\)}{\((x - 2)(9 - 4x)\)}{\((3x - 6)(2x - 3)\)}{}{}}
\End

\Data\ID{9000088808}
\Updated{2020-07-15 03:07:37}
\Level{B}
\SubArea{Polynomials and fractions}
\LANGen{
\Question{
Find the common denominator of the following three fractions. 
\[ 
 \text{$ \frac{a}
{a^{2}-ab}\ ,\qquad \frac{-b} 
{a^{2}-b^{2}} \ ,\qquad \frac{2b} 
{ab+b^{2}} $}
\]}
{\(ab(a^{2} - b^{2})\)}{\(ab(a - b)\)}{\(ab(a + b)\)}{\(ab(a + b)^{2}\)}{}{}}
\LANGpl{
\Question{
Znajdź wspólny mianownik podanych ułamków.
\[ 
 \text{$ \frac{a}
{a^{2}-ab}\ ,\qquad \frac{-b} 
{a^{2}-b^{2}} \ ,\qquad \frac{2b} 
{ab+b^{2}} $}
\]}
{\(ab(a^{2} - b^{2})\)}{\(ab(a - b)\)}{\(ab(a + b)\)}{\(ab(a + b)^{2}\)}{}{}}
\LANGcs{
\Question{
Společný jmenovatel lomených výrazů
\(\frac{a}
{a^{2}-ab}\),
\(\frac{-b}
{a^{2}-b^{2}} \),
\(\frac{2b}
{ab+b^{2}} \) je:}
{\(ab(a^{2} - b^{2})\)}{\(ab(a - b)\)}{\(ab(a + b)\)}{\(ab(a + b)^{2}\)}{}{}}
\LANGsk{
\Question{
Určte spoločný menovateľ lomených výrazov
\(\frac{a}
{a^{2}-ab}\),
\(\frac{-b}
{a^{2}-b^{2}} \),
\(\frac{2b}
{ab+b^{2}} \).}
{\(ab(a^{2} - b^{2})\)}{\(ab(a - b)\)}{\(ab(a + b)\)}{\(ab(a + b)^{2}\)}{}{}}
\LANGes{
\Question{
Find the common denominator of the following three fractions. 
\[ 
 \text{$ \frac{a}
{a^{2}-ab}\ ,\qquad \frac{-b} 
{a^{2}-b^{2}} \ ,\qquad \frac{2b} 
{ab+b^{2}} $}
\]}
{\(ab(a^{2} - b^{2})\)}{\(ab(a - b)\)}{\(ab(a + b)\)}{\(ab(a + b)^{2}\)}{}{}}
\End

\Data\ID{9000101605}
\Updated{2020-07-15 03:07:37}
\Level{B}
\SubArea{Polynomials and fractions}
\LANGen{
\Question{
Expand \(\left (4x^{2}y + 2xy^{2}\right )^{3}\).}
{\(64x^{6}y^{3} + 96x^{5}y^{4} + 48x^{4}y^{5} + 8x^{3}y^{6}\)}{\(16x^{2}y^{3} + 24x^{3}y^{3} + 8x^{3}y^{6}\)}{\(64x^{6}y^{3} + 96x^{3}y^{3} + 96x^{4}y^{5} + 8x^{3}y^{6}\)}{\(64x^{6}y^{3} + 8x^{3}y^{6}\)}{}{}}
\LANGpl{
\Question{
Rozwiń  \(\left (4x^{2}y + 2xy^{2}\right )^{3}\).}
{\(64x^{6}y^{3} + 96x^{5}y^{4} + 48x^{4}y^{5} + 8x^{3}y^{6}\)}{\(16x^{2}y^{3} + 24x^{3}y^{3} + 8x^{3}y^{6}\)}{\(64x^{6}y^{3} + 96x^{3}y^{3} + 96x^{4}y^{5} + 8x^{3}y^{6}\)}{\(64x^{6}y^{3} + 8x^{3}y^{6}\)}{}{}}
\LANGcs{
\Question{
Úpravou výrazu \(\left (4x^{2}y + 2xy^{2}\right )^{3}\) 
získáme:}
{\(64x^{6}y^{3} + 96x^{5}y^{4} + 48x^{4}y^{5} + 8x^{3}y^{6}\)}{\(16x^{2}y^{3} + 24x^{3}y^{3} + 8x^{3}y^{6}\)}{\(64x^{6}y^{3} + 96x^{3}y^{3} + 96x^{4}y^{5} + 8x^{3}y^{6}\)}{\(64x^{6}y^{3} + 8x^{3}y^{6}\)}{}{}}
\LANGsk{
\Question{
Umocnite daný výraz \(\left (4x^{2}y + 2xy^{2}\right )^{3}\).}
{\(64x^{6}y^{3} + 96x^{5}y^{4} + 48x^{4}y^{5} + 8x^{3}y^{6}\)}{\(16x^{2}y^{3} + 24x^{3}y^{3} + 8x^{3}y^{6}\)}{\(64x^{6}y^{3} + 96x^{3}y^{3} + 96x^{4}y^{5} + 8x^{3}y^{6}\)}{\(64x^{6}y^{3} + 8x^{3}y^{6}\)}{}{}}
\LANGes{
\Question{
Expand \(\left (4x^{2}y + 2xy^{2}\right )^{3}\).}
{\(64x^{6}y^{3} + 96x^{5}y^{4} + 48x^{4}y^{5} + 8x^{3}y^{6}\)}{\(16x^{2}y^{3} + 24x^{3}y^{3} + 8x^{3}y^{6}\)}{\(64x^{6}y^{3} + 96x^{3}y^{3} + 96x^{4}y^{5} + 8x^{3}y^{6}\)}{\(64x^{6}y^{3} + 8x^{3}y^{6}\)}{}{}}
\End

\Data\ID{9000101606}
\Updated{2020-07-15 03:07:37}
\Level{B}
\SubArea{Polynomials and fractions}
\LANGen{
\Question{
Expand \(\left (x - y\right )^{3} - x\left (x + y\right )^{2}\).}
{\(- y^{3} - 5x^{2}y + 2xy^{2}\)}{\(y^{3} - 5x^{2}y + 2xy^{2}\)}{\(- y^{3} - 5x^{2}y - 4xy^{2}\)}{\(- y^{3} - 5x^{2}y + 4xy^{2}\)}{}{}}
\LANGpl{
\Question{
Rozwiń  \(\left (x - y\right )^{3} - x\left (x + y\right )^{2}\).}
{\(- y^{3} - 5x^{2}y + 2xy^{2}\)}{\(y^{3} - 5x^{2}y + 2xy^{2}\)}{\(- y^{3} - 5x^{2}y - 4xy^{2}\)}{\(- y^{3} - 5x^{2}y + 4xy^{2}\)}{}{}}
\LANGcs{
\Question{
Úpravou výrazu \(\left (x - y\right )^{3} - x\left (x + y\right )^{2}\) 
získáme:}
{\(- y^{3} - 5x^{2}y + 2xy^{2}\)}{\(y^{3} - 5x^{2}y + 2xy^{2}\)}{\(- y^{3} - 5x^{2}y - 4xy^{2}\)}{\(- y^{3} - 5x^{2}y + 4xy^{2}\)}{}{}}
\LANGsk{
\Question{
Upravte daný výraz \(\left (x - y\right )^{3} - x\left (x + y\right )^{2}\).}
{\(- y^{3} - 5x^{2}y + 2xy^{2}\)}{\(y^{3} - 5x^{2}y + 2xy^{2}\)}{\(- y^{3} - 5x^{2}y - 4xy^{2}\)}{\(- y^{3} - 5x^{2}y + 4xy^{2}\)}{}{}}
\LANGes{
\Question{
Expand \(\left (x - y\right )^{3} - x\left (x + y\right )^{2}\).}
{\(- y^{3} - 5x^{2}y + 2xy^{2}\)}{\(y^{3} - 5x^{2}y + 2xy^{2}\)}{\(- y^{3} - 5x^{2}y - 4xy^{2}\)}{\(- y^{3} - 5x^{2}y + 4xy^{2}\)}{}{}}
\End

\Data\ID{9000101607}
\Updated{2020-07-15 03:07:37}
\Level{B}
\SubArea{Polynomials and fractions}
\LANGen{
\Question{
Expand \(\left (x^{2} - y\right )^{3} -\left (y + x^{2}\right )^{3}\).}
{\(- 6x^{4}y - 2y^{3}\)}{\(- 2y^{3}\)}{\(- 6x^{4}y - 2y^{3} + 6x^{2}y^{2}\)}{\(6x^{2}y - 2y^{3}\)}{}{}}
\LANGpl{
\Question{
Rozwiń  \(\left (x^{2} - y\right )^{3} -\left (y + x^{2}\right )^{3}\).}
{\(- 6x^{4}y - 2y^{3}\)}{\(- 2y^{3}\)}{\(- 6x^{4}y - 2y^{3} + 6x^{2}y^{2}\)}{\(6x^{2}y - 2y^{3}\)}{}{}}
\LANGcs{
\Question{
Úpravou výrazu \(\left (x^{2} - y\right )^{3} -\left (y + x^{2}\right )^{3}\) 
získáme:}
{\(- 6x^{4}y - 2y^{3}\)}{\(- 2y^{3}\)}{\(- 6x^{4}y - 2y^{3} + 6x^{2}y^{2}\)}{\(6x^{2}y - 2y^{3}\)}{}{}}
\LANGsk{
\Question{
Upravte daný výraz \(\left (x^{2} - y\right )^{3} -\left (y + x^{2}\right )^{3}\).}
{\(- 6x^{4}y - 2y^{3}\)}{\(- 2y^{3}\)}{\(- 6x^{4}y - 2y^{3} + 6x^{2}y^{2}\)}{\(6x^{2}y - 2y^{3}\)}{}{}}
\LANGes{
\Question{
Expand \(\left (x^{2} - y\right )^{3} -\left (y + x^{2}\right )^{3}\).}
{\(- 6x^{4}y - 2y^{3}\)}{\(- 2y^{3}\)}{\(- 6x^{4}y - 2y^{3} + 6x^{2}y^{2}\)}{\(6x^{2}y - 2y^{3}\)}{}{}}
\End

\Data\ID{9000101608}
\Updated{2020-07-15 03:07:37}
\Level{B}
\SubArea{Polynomials and fractions}
\LANGen{
\Question{
Expand \(\left (3x + y\right )\left (9x^{2} - 3xy + y^{2}\right )\).}
{\(27x^{3} + y^{3}\)}{\(27x^{3} - y^{3}\)}{\((3x + y)^{3}\)}{\(27x^{3} + 3y^{3}\)}{}{}}
\LANGpl{
\Question{
Rozwiń \(\left (3x + y\right )\left (9x^{2} - 3xy + y^{2}\right )\).}
{\(27x^{3} + y^{3}\)}{\(27x^{3} - y^{3}\)}{\((3x + y)^{3}\)}{\(27x^{3} + 3y^{3}\)}{}{}}
\LANGcs{
\Question{
Úpravou výrazu \(\left (3x + y\right )\left (9x^{2} - 3xy + y^{2}\right )\) 
získáme:}
{\(27x^{3} + y^{3}\)}{\(27x^{3} - y^{3}\)}{\((3x + y)^{3}\)}{\(27x^{3} + 3y^{3}\)}{}{}}
\LANGsk{
\Question{
Upravte daný výraz \(\left (3x + y\right )\left (9x^{2} - 3xy + y^{2}\right )\).}
{\(27x^{3} + y^{3}\)}{\(27x^{3} - y^{3}\)}{\((3x + y)^{3}\)}{\(27x^{3} + 3y^{3}\)}{}{}}
\LANGes{
\Question{
Expand \(\left (3x + y\right )\left (9x^{2} - 3xy + y^{2}\right )\).}
{\(27x^{3} + y^{3}\)}{\(27x^{3} - y^{3}\)}{\((3x + y)^{3}\)}{\(27x^{3} + 3y^{3}\)}{}{}}
\End

\Data\ID{9000101701}
\Updated{2020-07-15 03:07:37}
\Level{B}
\SubArea{Polynomials and fractions}
\LANGen{
\Question{
Factor the following polynomial. 
\[ 
 15xy - 10x - 3y + 2
\]}
{\(\left (5x - 1\right )\left (3y - 2\right )\)}{\(5x\left (3y - 2\right )\)}{\(4x\left (3y - 2\right )\)}{\(- 5x\left (3y - 2\right )\)}{}{}}
\LANGpl{
\Question{
Rozłóż na czynniki podany wielomian.
\[ 
 15xy - 10x - 3y + 2
\]}
{\(\left (5x - 1\right )\left (3y - 2\right )\)}{\(5x\left (3y - 2\right )\)}{\(4x\left (3y - 2\right )\)}{\(- 5x\left (3y - 2\right )\)}{}{}}
\LANGcs{
\Question{
Rozložením výrazu \(15xy - 10x - 3y + 2\) 
na součin získáme výraz:}
{\(\left (5x - 1\right )\left (3y - 2\right )\)}{\(5x\left (3y - 2\right )\)}{\(4x\left (3y - 2\right )\)}{\(- 5x\left (3y - 2\right )\)}{}{}}
\LANGsk{
\Question{
Upravte na súčin. 
\[ 
 15xy - 10x - 3y + 2
\]}
{\(\left (5x - 1\right )\left (3y - 2\right )\)}{\(5x\left (3y - 2\right )\)}{\(4x\left (3y - 2\right )\)}{\(- 5x\left (3y - 2\right )\)}{}{}}
\LANGes{
\Question{
Factor the following polynomial. 
\[ 
 15xy - 10x - 3y + 2
\]}
{\(\left (5x - 1\right )\left (3y - 2\right )\)}{\(5x\left (3y - 2\right )\)}{\(4x\left (3y - 2\right )\)}{\(- 5x\left (3y - 2\right )\)}{}{}}
\End

\Data\ID{9000101702}
\Updated{2020-07-15 03:07:37}
\Level{B}
\SubArea{Polynomials and fractions}
\LANGen{
\Question{
Factor the following polynomial. 
\[ 
 3x^{3} + 3x^{2}y + 4xy + 4y^{2}
\]}
{\(\left (3x^{2} + 4y\right )\left (x + y\right )\)}{\(\left (3x + y\right )\left (x^{2} + y^{2}\right )\)}{\(\left (3x^{2} + 4\right )\left (x + y^{2}\right )\)}{\(\left (3x + y^{2}\right )\left (x + y\right )\)}{}{}}
\LANGpl{
\Question{
Rozłóż na czynniki podany wielomian:
\[ 
 3x^{3} + 3x^{2}y + 4xy + 4y^{2}
\]}
{\(\left (3x^{2} + 4y\right )\left (x + y\right )\)}{\(\left (3x + y\right )\left (x^{2} + y^{2}\right )\)}{\(\left (3x^{2} + 4\right )\left (x + y^{2}\right )\)}{\(\left (3x + y^{2}\right )\left (x + y\right )\)}{}{}}
\LANGcs{
\Question{
Rozložením výrazu \(3x^{3} + 3x^{2}y + 4xy + 4y^{2}\) 
na součin získáme výraz:}
{\(\left (3x^{2} + 4y\right )\left (x + y\right )\)}{\(\left (3x + y\right )\left (x^{2} + y^{2}\right )\)}{\(\left (3x^{2} + 4\right )\left (x + y^{2}\right )\)}{\(\left (3x + y^{2}\right )\left (x + y\right )\)}{}{}}
\LANGsk{
\Question{
Upravte na súčin. 
\[ 
 3x^{3} + 3x^{2}y + 4xy + 4y^{2}
\]}
{\(\left (3x^{2} + 4y\right )\left (x + y\right )\)}{\(\left (3x + y\right )\left (x^{2} + y^{2}\right )\)}{\(\left (3x^{2} + 4\right )\left (x + y^{2}\right )\)}{\(\left (3x + y^{2}\right )\left (x + y\right )\)}{}{}}
\LANGes{
\Question{
Factor the following polynomial. 
\[ 
 3x^{3} + 3x^{2}y + 4xy + 4y^{2}
\]}
{\(\left (3x^{2} + 4y\right )\left (x + y\right )\)}{\(\left (3x + y\right )\left (x^{2} + y^{2}\right )\)}{\(\left (3x^{2} + 4\right )\left (x + y^{2}\right )\)}{\(\left (3x + y^{2}\right )\left (x + y\right )\)}{}{}}
\End

\Data\ID{9000101703}
\Updated{2020-07-15 03:07:37}
\Level{B}
\SubArea{Polynomials and fractions}
\LANGen{
\Question{
Factor the following polynomial. 
\[ 
 \left (5x - y\right )^{2} -\left (x - y\right )^{2}
\]}
{\(4x\left (6x - 2y\right )\)}{\(x\left (5x - y\right )\)}{\(6x\left (6x - 2y\right )\)}{\(- 32x^{2}\)}{}{}}
\LANGpl{
\Question{
Rozłóż na czynniki podany wielomian:
\[ 
 \left (5x - y\right )^{2} -\left (x - y\right )^{2}
\]}
{\(4x\left (6x - 2y\right )\)}{\(x\left (5x - y\right )\)}{\(6x\left (6x - 2y\right )\)}{\(- 32x^{2}\)}{}{}}
\LANGcs{
\Question{
Rozložením výrazu \(\left (5x - y\right )^{2} -\left (x - y\right )^{2}\) 
na součin získáme výraz:}
{\(4x\left (6x - 2y\right )\)}{\(x\left (5x - y\right )\)}{\(6x\left (6x - 2y\right )\)}{\(- 32x^{2}\)}{}{}}
\LANGsk{
\Question{
Upravte na súčin. 
\[ 
 \left (5x - y\right )^{2} -\left (x - y\right )^{2}
\]}
{\(4x\left (6x - 2y\right )\)}{\(x\left (5x - y\right )\)}{\(6x\left (6x - 2y\right )\)}{\(- 32x^{2}\)}{}{}}
\LANGes{
\Question{
Factor the following polynomial. 
\[ 
 \left (5x - y\right )^{2} -\left (x - y\right )^{2}
\]}
{\(4x\left (6x - 2y\right )\)}{\(x\left (5x - y\right )\)}{\(6x\left (6x - 2y\right )\)}{\(- 32x^{2}\)}{}{}}
\End

\Data\ID{9000101704}
\Updated{2020-07-15 03:07:37}
\Level{B}
\SubArea{Polynomials and fractions}
\LANGen{
\Question{
Factor the following polynomial. 
\[ 
 16x^{2}y^{4} - 25x^{4}y^{2}
\]}
{\(\left (4xy^{2} - 5x^{2}y\right )\left (4xy^{2} + 5x^{2}y\right )\)}{\(\left (4xy - 5x^{2}y\right )\left (4xy^{2} + 5xy\right )\)}{\(\left (4x^{2}y^{2} - 5xy\right )\left (4x^{2}y^{2} + 5xy\right )\)}{\(\left (4xy^{2} - 5x^{2}y\right )^{2}\)}{}{}}
\LANGpl{
\Question{
Rozłóż na czynniki podany wielomian:
\[ 
 16x^{2}y^{4} - 25x^{4}y^{2}
\]}
{\(\left (4xy^{2} - 5x^{2}y\right )\left (4xy^{2} + 5x^{2}y\right )\)}{\(\left (4xy - 5x^{2}y\right )\left (4xy^{2} + 5xy\right )\)}{\(\left (4x^{2}y^{2} - 5xy\right )\left (4x^{2}y^{2} + 5xy\right )\)}{\(\left (4xy^{2} - 5x^{2}y\right )^{2}\)}{}{}}
\LANGcs{
\Question{
Rozložením výrazu \(16x^{2}y^{4} - 25x^{4}y^{2}\) 
na součin získáme výraz:}
{\(\left (4xy^{2} - 5x^{2}y\right )\left (4xy^{2} + 5x^{2}y\right )\)}{\(\left (4xy - 5x^{2}y\right )\left (4xy^{2} + 5xy\right )\)}{\(\left (4x^{2}y^{2} - 5xy\right )\left (4x^{2}y^{2} + 5xy\right )\)}{\(\left (4xy^{2} - 5x^{2}y\right )^{2}\)}{}{}}
\LANGsk{
\Question{
Upravte na súčin. 
\[ 
 16x^{2}y^{4} - 25x^{4}y^{2}
\]}
{\(\left (4xy^{2} - 5x^{2}y\right )\left (4xy^{2} + 5x^{2}y\right )\)}{\(\left (4xy - 5x^{2}y\right )\left (4xy^{2} + 5xy\right )\)}{\(\left (4x^{2}y^{2} - 5xy\right )\left (4x^{2}y^{2} + 5xy\right )\)}{\(\left (4xy^{2} - 5x^{2}y\right )^{2}\)}{}{}}
\LANGes{
\Question{
Factor the following polynomial. 
\[ 
 16x^{2}y^{4} - 25x^{4}y^{2}
\]}
{\(\left (4xy^{2} - 5x^{2}y\right )\left (4xy^{2} + 5x^{2}y\right )\)}{\(\left (4xy - 5x^{2}y\right )\left (4xy^{2} + 5xy\right )\)}{\(\left (4x^{2}y^{2} - 5xy\right )\left (4x^{2}y^{2} + 5xy\right )\)}{\(\left (4xy^{2} - 5x^{2}y\right )^{2}\)}{}{}}
\End

\Data\ID{9000101705}
\Updated{2020-07-15 03:07:37}
\Level{B}
\SubArea{Polynomials and fractions}
\LANGen{
\Question{
Factor the following polynomial expression. 
\[ 
 16a^{2}b^{2} - 4a^{2}c^{2} - 16b^{2}d^{2} + 4c^{2}d^{2}
\]}
{\(4\left (a - d\right )\left (a + d\right )\left (2b + c\right )\left (2b - c\right )\)}{\(4\left (a + b\right )^{2}\left (2b + c\right )^{2}\)}{\(4\left (a - b\right )\left (a + b\right )\left (2b + c\right )\left (2b - c\right )\)}{\(4\left (a - c\right )\left (a + c\right )\left (2b + d\right )\left (2b - d\right )\)}{}{}}
\LANGpl{
\Question{
Rozłóż na czynniki podany wielomian:
\[ 
 16a^{2}b^{2} - 4a^{2}c^{2} - 16b^{2}d^{2} + 4c^{2}d^{2}
\]}
{\(4\left (a - d\right )\left (a + d\right )\left (2b + c\right )\left (2b - c\right )\)}{\(4\left (a + b\right )^{2}\left (2b + c\right )^{2}\)}{\(4\left (a - b\right )\left (a + b\right )\left (2b + c\right )\left (2b - c\right )\)}{\(4\left (a - c\right )\left (a + c\right )\left (2b + d\right )\left (2b - d\right )\)}{}{}}
\LANGcs{
\Question{
Rozložením výrazu \(16a^{2}b^{2} - 4a^{2}c^{2} - 16b^{2}d^{2} + 4c^{2}d^{2}\) 
na součin získáme výraz:}
{\(4\left (a - d\right )\left (a + d\right )\left (2b + c\right )\left (2b - c\right )\)}{\(4\left (a + b\right )^{2}\left (2b + c\right )^{2}\)}{\(4\left (a - b\right )\left (a + b\right )\left (2b + c\right )\left (2b - c\right )\)}{\(4\left (a - c\right )\left (a + c\right )\left (2b + d\right )\left (2b - d\right )\)}{}{}}
\LANGsk{
\Question{
Upravte na súčin. 
\[ 
 16a^{2}b^{2} - 4a^{2}c^{2} - 16b^{2}d^{2} + 4c^{2}d^{2}
\]}
{\(4\left (a - d\right )\left (a + d\right )\left (2b + c\right )\left (2b - c\right )\)}{\(4\left (a + b\right )^{2}\left (2b + c\right )^{2}\)}{\(4\left (a - b\right )\left (a + b\right )\left (2b + c\right )\left (2b - c\right )\)}{\(4\left (a - c\right )\left (a + c\right )\left (2b + d\right )\left (2b - d\right )\)}{}{}}
\LANGes{
\Question{
Factor the following polynomial expression. 
\[ 
 16a^{2}b^{2} - 4a^{2}c^{2} - 16b^{2}d^{2} + 4c^{2}d^{2}
\]}
{\(4\left (a - d\right )\left (a + d\right )\left (2b + c\right )\left (2b - c\right )\)}{\(4\left (a + b\right )^{2}\left (2b + c\right )^{2}\)}{\(4\left (a - b\right )\left (a + b\right )\left (2b + c\right )\left (2b - c\right )\)}{\(4\left (a - c\right )\left (a + c\right )\left (2b + d\right )\left (2b - d\right )\)}{}{}}
\End

\Data\ID{9000101706}
\Updated{2020-07-15 03:07:37}
\Level{B}
\SubArea{Polynomials and fractions}
\LANGen{
\Question{
Factor the following polynomial. 
\[ 
 8x^{4} - 48x^{3} + 72x^{2}
\]}
{\(8x^{2}\left (x - 3\right )^{2}\)}{\(- 8x^{2}\left (3 - x\right )^{2}\)}{\(8\left (x^{2} - 3\right )^{2}\)}{\(8x\left (x^{2} - 3\right )^{2}\)}{}{}}
\LANGpl{
\Question{
Rozłóż na czynniki podany wielomian: 
\[ 
 8x^{4} - 48x^{3} + 72x^{2}
\]}
{\(8x^{2}\left (x - 3\right )^{2}\)}{\(- 8x^{2}\left (3 - x\right )^{2}\)}{\(8\left (x^{2} - 3\right )^{2}\)}{\(8x\left (x^{2} - 3\right )^{2}\)}{}{}}
\LANGcs{
\Question{
Rozložením výrazu \(8x^{4} - 48x^{3} + 72x^{2}\) 
získáme výraz:}
{\(8x^{2}\left (x - 3\right )^{2}\)}{\(- 8x^{2}\left (3 - x\right )^{2}\)}{\(8\left (x^{2} - 3\right )^{2}\)}{\(8x\left (x^{2} - 3\right )^{2}\)}{}{}}
\LANGsk{
\Question{
Upravte na súčin. 
\[ 
 8x^{4} - 48x^{3} + 72x^{2}
\]}
{\(8x^{2}\left (x - 3\right )^{2}\)}{\(- 8x^{2}\left (3 - x\right )^{2}\)}{\(8\left (x^{2} - 3\right )^{2}\)}{\(8x\left (x^{2} - 3\right )^{2}\)}{}{}}
\LANGes{
\Question{
Factor the following polynomial. 
\[ 
 8x^{4} - 48x^{3} + 72x^{2}
\]}
{\(8x^{2}\left (x - 3\right )^{2}\)}{\(- 8x^{2}\left (3 - x\right )^{2}\)}{\(8\left (x^{2} - 3\right )^{2}\)}{\(8x\left (x^{2} - 3\right )^{2}\)}{}{}}
\End

\Data\ID{9000101710}
\Updated{2020-07-15 03:07:37}
\Level{B}
\SubArea{Polynomials and fractions}
\LANGen{
\Question{
Factor the following polynomial. 
\[ 
 x^{2}y - x^{2}z - 4xyz + 4xy^{2} + 4y^{3} - 4y^{2}z
\]}
{\(\left (y - z\right )\left (x + 2y\right )^{2}\)}{\(\left (y - z\right )\left (x - 2y\right )^{2}\)}{\(\left (y - z\right )\left (x^{2} + 4y + 4y^{2}\right )\)}{\(\left (y + z\right )\left (x - 2y\right )^{2}\)}{}{}}
\LANGpl{
\Question{
Rozłóż na czynniki podany wielomian:
\[ 
 x^{2}y - x^{2}z - 4xyz + 4xy^{2} + 4y^{3} - 4y^{2}z
\]}
{\(\left (y - z\right )\left (x + 2y\right )^{2}\)}{\(\left (y - z\right )\left (x - 2y\right )^{2}\)}{\(\left (y - z\right )\left (x^{2} + 4y + 4y^{2}\right )\)}{\(\left (y + z\right )\left (x - 2y\right )^{2}\)}{}{}}
\LANGcs{
\Question{
Rozložením výrazu \(x^{2}y - x^{2}z - 4xyz + 4xy^{2} + 4y^{3} - 4y^{2}z\) 
získáme výraz:}
{\(\left (y - z\right )\left (x + 2y\right )^{2}\)}{\(\left (y - z\right )\left (x - 2y\right )^{2}\)}{\(\left (y - z\right )\left (x^{2} + 4y + 4y^{2}\right )\)}{\(\left (y + z\right )\left (x - 2y\right )^{2}\)}{}{}}
\LANGsk{
\Question{
Upravte na súčin. 
\[ 
 x^{2}y - x^{2}z - 4xyz + 4xy^{2} + 4y^{3} - 4y^{2}z
\]}
{\(\left (y - z\right )\left (x + 2y\right )^{2}\)}{\(\left (y - z\right )\left (x - 2y\right )^{2}\)}{\(\left (y - z\right )\left (x^{2} + 4y + 4y^{2}\right )\)}{\(\left (y + z\right )\left (x - 2y\right )^{2}\)}{}{}}
\LANGes{
\Question{
Factor the following polynomial. 
\[ 
 x^{2}y - x^{2}z - 4xyz + 4xy^{2} + 4y^{3} - 4y^{2}z
\]}
{\(\left (y - z\right )\left (x + 2y\right )^{2}\)}{\(\left (y - z\right )\left (x - 2y\right )^{2}\)}{\(\left (y - z\right )\left (x^{2} + 4y + 4y^{2}\right )\)}{\(\left (y + z\right )\left (x - 2y\right )^{2}\)}{}{}}
\End

\Data\ID{9000146201}
\Updated{2020-07-15 03:07:37}
\Level{B}
\SubArea{Polynomials and fractions}
\LANGen{
\Question{
Expand the following expression. 
\[ 
 \left (2x^{3} - y^{2}\right )^{3}
\]}
{\(8x^{9} - 12x^{6}y^{2} + 6x^{3}y^{4} - y^{6}\)}{\(8x^{9} - 4x^{6}y^{2} + 2x^{3}y^{4} - y^{6}\)}{\(8x^{6} - 12x^{5}y^{2} + 6x^{3}y^{4} - y^{5}\)}{\(8x^{6} - 4x^{5}y^{2} + 2x^{3}y^{4} - y^{5}\)}{}{}}
\LANGpl{
\Question{
Rozwiń podane wyrażenie:
\[ 
 \left (2x^{3} - y^{2}\right )^{3}
\]}
{\(8x^{9} - 12x^{6}y^{2} + 6x^{3}y^{4} - y^{6}\)}{\(8x^{9} - 4x^{6}y^{2} + 2x^{3}y^{4} - y^{6}\)}{\(8x^{6} - 12x^{5}y^{2} + 6x^{3}y^{4} - y^{5}\)}{\(8x^{6} - 4x^{5}y^{2} + 2x^{3}y^{4} - y^{5}\)}{}{}}
\LANGcs{
\Question{
Umocněním \(\left (2x^{3} - y^{2}\right )^{3}\) 
získáme výraz:}
{\(8x^{9} - 12x^{6}y^{2} + 6x^{3}y^{4} - y^{6}\)}{\(8x^{9} - 4x^{6}y^{2} + 2x^{3}y^{4} - y^{6}\)}{\(8x^{6} - 12x^{5}y^{2} + 6x^{3}y^{4} - y^{5}\)}{\(8x^{6} - 4x^{5}y^{2} + 2x^{3}y^{4} - y^{5}\)}{}{}}
\LANGsk{
\Question{
Umocnite daný výraz.
\[ 
 \left (2x^{3} - y^{2}\right )^{3}
\]}
{\(8x^{9} - 12x^{6}y^{2} + 6x^{3}y^{4} - y^{6}\)}{\(8x^{9} - 4x^{6}y^{2} + 2x^{3}y^{4} - y^{6}\)}{\(8x^{6} - 12x^{5}y^{2} + 6x^{3}y^{4} - y^{5}\)}{\(8x^{6} - 4x^{5}y^{2} + 2x^{3}y^{4} - y^{5}\)}{}{}}
\LANGes{
\Question{
Expand the following expression. 
\[ 
 \left (2x^{3} - y^{2}\right )^{3}
\]}
{\(8x^{9} - 12x^{6}y^{2} + 6x^{3}y^{4} - y^{6}\)}{\(8x^{9} - 4x^{6}y^{2} + 2x^{3}y^{4} - y^{6}\)}{\(8x^{6} - 12x^{5}y^{2} + 6x^{3}y^{4} - y^{5}\)}{\(8x^{6} - 4x^{5}y^{2} + 2x^{3}y^{4} - y^{5}\)}{}{}}
\End

\Data\ID{9000146202}
\Updated{2020-07-15 03:07:37}
\Level{B}
\SubArea{Polynomials and fractions}
\LANGen{
\Question{
Expand the following expression. 
\[ 
 \left (a^{2} + \sqrt{3}b\right )^{3}
\]}
{\(a^{6} + 3\sqrt{3}a^{4}b + 9a^{2}b^{2} + 3\sqrt{3}b^{3}\)}{\(a^{6} + \sqrt{3}a^{4}b + 3a^{2}b^{2} + 3\sqrt{3}b^{3}\)}{\(a^{5} + 3\sqrt{3}a^{4}b + 9a^{2}b^{2} + 3\sqrt{3}b^{3}\)}{\(a^{5} + \sqrt{3}a^{4}b + 3a^{2}b^{2} + 3\sqrt{3}b^{3}\)}{}{}}
\LANGpl{
\Question{
Rozwiń podane wyrażenie: 
\[ 
 \left (a^{2} + \sqrt{3}b\right )^{3}
\]}
{\(a^{6} + 3\sqrt{3}a^{4}b + 9a^{2}b^{2} + 3\sqrt{3}b^{3}\)}{\(a^{6} + \sqrt{3}a^{4}b + 3a^{2}b^{2} + 3\sqrt{3}b^{3}\)}{\(a^{5} + 3\sqrt{3}a^{4}b + 9a^{2}b^{2} + 3\sqrt{3}b^{3}\)}{\(a^{5} + \sqrt{3}a^{4}b + 3a^{2}b^{2} + 3\sqrt{3}b^{3}\)}{}{}}
\LANGcs{
\Question{
Umocněním \(\left (a^{2} + \sqrt{3}b\right )^{3}\) 
získáme výraz:}
{\(a^{6} + 3\sqrt{3}a^{4}b + 9a^{2}b^{2} + 3\sqrt{3}b^{3}\)}{\(a^{6} + \sqrt{3}a^{4}b + 3a^{2}b^{2} + 3\sqrt{3}b^{3}\)}{\(a^{5} + 3\sqrt{3}a^{4}b + 9a^{2}b^{2} + 3\sqrt{3}b^{3}\)}{\(a^{5} + \sqrt{3}a^{4}b + 3a^{2}b^{2} + 3\sqrt{3}b^{3}\)}{}{}}
\LANGsk{
\Question{
Umocnite daný výraz.
\[ 
 \left (a^{2} + \sqrt{3}b\right )^{3}
\]}
{\(a^{6} + 3\sqrt{3}a^{4}b + 9a^{2}b^{2} + 3\sqrt{3}b^{3}\)}{\(a^{6} + \sqrt{3}a^{4}b + 3a^{2}b^{2} + 3\sqrt{3}b^{3}\)}{\(a^{5} + 3\sqrt{3}a^{4}b + 9a^{2}b^{2} + 3\sqrt{3}b^{3}\)}{\(a^{5} + \sqrt{3}a^{4}b + 3a^{2}b^{2} + 3\sqrt{3}b^{3}\)}{}{}}
\LANGes{
\Question{
Expand the following expression. 
\[ 
 \left (a^{2} + \sqrt{3}b\right )^{3}
\]}
{\(a^{6} + 3\sqrt{3}a^{4}b + 9a^{2}b^{2} + 3\sqrt{3}b^{3}\)}{\(a^{6} + \sqrt{3}a^{4}b + 3a^{2}b^{2} + 3\sqrt{3}b^{3}\)}{\(a^{5} + 3\sqrt{3}a^{4}b + 9a^{2}b^{2} + 3\sqrt{3}b^{3}\)}{\(a^{5} + \sqrt{3}a^{4}b + 3a^{2}b^{2} + 3\sqrt{3}b^{3}\)}{}{}}
\End

\Data\ID{9000146207}
\Updated{2020-07-15 03:07:37}
\Level{B}
\SubArea{Polynomials and fractions}
\LANGen{
\Question{
Factor the following expression. 
\[ 
 4a^{2} -\left (a - 1\right )^{2}
\]}
{\(\left (a + 1\right )\left (3a - 1\right )\)}{\(\left (a - 1\right )\left (3a - 1\right )\)}{\(\left (a + 1\right )\left (3a + 1\right )\)}{\(\left (a - 1\right )\left (3a + 1\right )\)}{}{}}
\LANGpl{
\Question{
Rozłóż na czynniki wyrażenie:
\[ 
 4a^{2} -\left (a - 1\right )^{2}
\]}
{\(\left (a + 1\right )\left (3a - 1\right )\)}{\(\left (a - 1\right )\left (3a - 1\right )\)}{\(\left (a + 1\right )\left (3a + 1\right )\)}{\(\left (a - 1\right )\left (3a + 1\right )\)}{}{}}
\LANGcs{
\Question{
Rozložením výrazu \(4a^{2} -\left (a - 1\right )^{2}\) 
na součin získáme výsledek:}
{\(\left (a + 1\right )\left (3a - 1\right )\)}{\(\left (a - 1\right )\left (3a - 1\right )\)}{\(\left (a + 1\right )\left (3a + 1\right )\)}{\(\left (a - 1\right )\left (3a + 1\right )\)}{}{}}
\LANGsk{
\Question{
Upravte na súčin.
\[ 
 4a^{2} -\left (a - 1\right )^{2}
\]}
{\(\left (a + 1\right )\left (3a - 1\right )\)}{\(\left (a - 1\right )\left (3a - 1\right )\)}{\(\left (a + 1\right )\left (3a + 1\right )\)}{\(\left (a - 1\right )\left (3a + 1\right )\)}{}{}}
\LANGes{
\Question{
Factor the following expression. 
\[ 
 4a^{2} -\left (a - 1\right )^{2}
\]}
{\(\left (a + 1\right )\left (3a - 1\right )\)}{\(\left (a - 1\right )\left (3a - 1\right )\)}{\(\left (a + 1\right )\left (3a + 1\right )\)}{\(\left (a - 1\right )\left (3a + 1\right )\)}{}{}}
\End

\Data\ID{9000146208}
\Updated{2020-07-15 03:07:37}
\Level{B}
\SubArea{Polynomials and fractions}
\LANGen{
\Question{
Factor the following expression. 
\[ 
 \left (2x - 1\right )^{2} -\left (x + 3\right )^{2}
\]}
{\(\left (x - 4\right )\left (3x + 2\right )\)}{\(\left (x - 4\right )\left (3x - 2\right )\)}{\(\left (x + 4\right )\left (3x + 2\right )\)}{\(\left (x + 4\right )\left (3x - 2\right )\)}{}{}}
\LANGpl{
\Question{
Rozłóż na czynniki wyrażenie:
\[ 
 \left (2x - 1\right )^{2} -\left (x + 3\right )^{2}
\]}
{\(\left (x - 4\right )\left (3x + 2\right )\)}{\(\left (x - 4\right )\left (3x - 2\right )\)}{\(\left (x + 4\right )\left (3x + 2\right )\)}{\(\left (x + 4\right )\left (3x - 2\right )\)}{}{}}
\LANGcs{
\Question{
Rozložením výrazu \(\left (2x - 1\right )^{2} -\left (x + 3\right )^{2}\) 
na součin získáme výsledek:}
{\(\left (x - 4\right )\left (3x + 2\right )\)}{\(\left (x - 4\right )\left (3x - 2\right )\)}{\(\left (x + 4\right )\left (3x + 2\right )\)}{\(\left (x + 4\right )\left (3x - 2\right )\)}{}{}}
\LANGsk{
\Question{
Upravte na súčin.
\[ 
 \left (2x - 1\right )^{2} -\left (x + 3\right )^{2}
\]}
{\(\left (x - 4\right )\left (3x + 2\right )\)}{\(\left (x - 4\right )\left (3x - 2\right )\)}{\(\left (x + 4\right )\left (3x + 2\right )\)}{\(\left (x + 4\right )\left (3x - 2\right )\)}{}{}}
\LANGes{
\Question{
Factor the following expression. 
\[ 
 \left (2x - 1\right )^{2} -\left (x + 3\right )^{2}
\]}
{\(\left (x - 4\right )\left (3x + 2\right )\)}{\(\left (x - 4\right )\left (3x - 2\right )\)}{\(\left (x + 4\right )\left (3x + 2\right )\)}{\(\left (x + 4\right )\left (3x - 2\right )\)}{}{}}
\End

\Data\ID{1003032402}
\Updated{2020-07-15 03:07:21}
\Level{C}
\SubArea{Polynomials and fractions}
\LANGen{
\Question{
The polynomial \( 27+x^3 \) is equal to a product:}
{\( (3+x)\left(9-3x+x^2\right) \)}{\( (3+x)^2(3-x) \)}{\( (3-x)\left(9+3x+x^2\right) \)}{\( (3+x)^3 \)}{}{}}
\LANGpl{
\Question{
Wielomian \( 27+x^3 \) jest równy:}
{\( (3+x)\left(9-3x+x^2\right) \)}{\( (3+x)^2(3-x) \)}{\( (3-x)\left(9+3x+x^2\right) \)}{\( (3+x)^3 \)}{}{}}
\LANGcs{
\Question{
Mnohočlen \( 27+x^3 \) se rovná součinu:}
{\( (3+x)\left(9-3x+x^2\right) \)}{\( (3+x)^2(3-x) \)}{\( (3-x)\left(9+3x+x^2\right) \)}{\( (3+x)^3 \)}{}{}}
\LANGsk{
\Question{
Rozkladom polynómu \( 27+x^3 \) na súčin dostaneme:}
{\( (3+x)\left(9-3x+x^2\right) \)}{\( (3+x)^2(3-x) \)}{\( (3-x)\left(9+3x+x^2\right) \)}{\( (3+x)^3 \)}{}{}}
\LANGes{
\Question{
The polynomial \( 27+x^3 \) is equal to a product:}
{\( (3+x)\left(9-3x+x^2\right) \)}{\( (3+x)^2(3-x) \)}{\( (3-x)\left(9+3x+x^2\right) \)}{\( (3+x)^3 \)}{}{}}
\End

\Data\ID{1003032403}
\Updated{2020-07-15 03:07:21}
\Level{C}
\SubArea{Polynomials and fractions}
\LANGen{
\Question{
Reducing the rational expression \( \frac{4m^2-4mn+n^2}{8m^3-n^3} \) we get:}
{\( \frac{2m-n}{4m^2+2mn+n^2} \)}{\( \frac{m-4mn+1}{2m-n} \)}{\( \frac{2m-n}{4m^2-4mn+n^2} \)}{\( \frac{2m-n}{4m^2+4mn+n^2} \)}{}{}}
\LANGpl{
\Question{
Skracając wyrażenia wymierne  \( \frac{4m^2-4mn+n^2}{8m^3-n^3} \) otrzymamy:}
{\( \frac{2m-n}{4m^2+2mn+n^2} \)}{\( \frac{m-4mn+1}{2m-n} \)}{\( \frac{2m-n}{4m^2-4mn+n^2} \)}{\( \frac{2m-n}{4m^2+4mn+n^2} \)}{}{}}
\LANGcs{
\Question{
Zjednodušením lomeného výrazu \( \frac{4m^2-4mn+n^2}{8m^3-n^3} \) získáme:}
{\( \frac{2m-n}{4m^2+2mn+n^2} \)}{\( \frac{m-4mn+1}{2m-n} \)}{\( \frac{2m-n}{4m^2-4mn+n^2} \)}{\( \frac{2m-n}{4m^2+4mn+n^2} \)}{}{}}
\LANGsk{
\Question{
Zjednodušením výrazu \( \frac{4m^2-4mn+n^2}{8m^3-n^3} \) dostaneme:}
{\( \frac{2m-n}{4m^2+2mn+n^2} \)}{\( \frac{m-4mn+1}{2m-n} \)}{\( \frac{2m-n}{4m^2-4mn+n^2} \)}{\( \frac{2m-n}{4m^2+4mn+n^2} \)}{}{}}
\LANGes{
\Question{
Reducing the rational expression \( \frac{4m^2-4mn+n^2}{8m^3-n^3} \) we get:}
{\( \frac{2m-n}{4m^2+2mn+n^2} \)}{\( \frac{m-4mn+1}{2m-n} \)}{\( \frac{2m-n}{4m^2-4mn+n^2} \)}{\( \frac{2m-n}{4m^2+4mn+n^2} \)}{}{}}
\End

\Data\ID{1003032404}
\Updated{2020-07-15 03:07:21}
\Level{C}
\SubArea{Polynomials and fractions}
\LANGen{
\Question{
Simplifying the rational expression \( \frac{x^6-y^6}{x^2-y^2} \)   we get:}
{\( \left(x^2-xy+y^2\right)\left(x^2+xy+y^2\right) \)}{\( x^4-y^4 \)}{\( x^3 - y^3 \)}{\( \left(x^2-2xy+y^2\right)\left(x^2+2xy+y^2\right) \)}{}{}}
\LANGpl{
\Question{
Skracając wyrażenia wymierne  \( \frac{x^6-y^6}{x^2-y^2} \)   otrzymamy:}
{\( \left(x^2-xy+y^2\right)\left(x^2+xy+y^2\right) \)}{\( x^4-y^4 \)}{\( x^3 - y^3 \)}{\( \left(x^2-2xy+y^2\right)\left(x^2+2xy+y^2\right) \)}{}{}}
\LANGcs{
\Question{
Zjednodušením lomeného výrazu \( \frac{x^6-y^6}{x^2-y^2} \)  dostaneme:}
{\( \left(x^2-xy+y^2\right)\left(x^2+xy+y^2\right) \)}{\( x^4-y^4 \)}{\( x^3 - y^3 \)}{\( \left(x^2-2xy+y^2\right)\left(x^2+2xy+y^2\right) \)}{}{}}
\LANGsk{
\Question{
Zjednodušením výrazu \( \frac{x^6-y^6}{x^2-y^2} \) dostaneme:}
{\( \left(x^2-xy+y^2\right)\left(x^2+xy+y^2\right) \)}{\( x^4-y^4 \)}{\( x^3 - y^3 \)}{\( \left(x^2-2xy+y^2\right)\left(x^2+2xy+y^2\right) \)}{}{}}
\LANGes{
\Question{
Simplifying the rational expression \( \frac{x^6-y^6}{x^2-y^2} \)   we get:}
{\( \left(x^2-xy+y^2\right)\left(x^2+xy+y^2\right) \)}{\( x^4-y^4 \)}{\( x^3 - y^3 \)}{\( \left(x^2-2xy+y^2\right)\left(x^2+2xy+y^2\right) \)}{}{}}
\End

\Data\ID{1003032502}
\Updated{2020-07-15 03:07:21}
\Level{C}
\SubArea{Polynomials and fractions}
\LANGen{
\Question{
Let the polynomial \( (x-2)^5-(x+2)^5 \) be expressed in the form \( a_5x^5+a_4x^4+a_3x^3+a_2x^2+a_1x+a_0 \). Give the sum of \( a_5+a_4+a_3+a_2+a_1 \).}
{\( -180 \)}{\( -244 \)}{\( -242 \)}{\( -212 \)}{}{}}
\LANGpl{
\Question{
Zapisz wielomian  \( (x-2)^5-(x+2)^5 \) w postaci \( a_5x^5+a_4x^4+a_3x^3+a_2x^2+a_1x+a_0 \). Wyznacz sumę \( a_5+a_4+a_3+a_2+a_1 \).}
{\( -180 \)}{\( -244 \)}{\( -242 \)}{\( -212 \)}{}{}}
\LANGcs{
\Question{
Mnohočlen \( (x-2)^5-(x+2)^5 \) je vyjádřen jako \( a_5x^5+a_4x^4+a_3x^3+a_2x^2+a_1x+a_0 \). Uveďte součet \( a_5+a_4+a_3+a_2+a_1 \).}
{\( -180 \)}{\( -244 \)}{\( -242 \)}{\( -212 \)}{}{}}
\LANGsk{
\Question{
Polynóm \( (x-2)^5-(x+2)^5 \) môžme vyjadriť v tvare \( a_5x^5+a_4x^4+a_3x^3+a_2x^2+a_1x+a_0 \). Nájdite súčet koeficientov \( a_5+a_4+a_3+a_2+a_1 \).}
{\( -180 \)}{\( -244 \)}{\( -242 \)}{\( -212 \)}{}{}}
\LANGes{
\Question{
Let the polynomial \( (x-2)^5-(x+2)^5 \) be expressed in the form \( a_5x^5+a_4x^4+a_3x^3+a_2x^2+a_1x+a_0 \). Give the sum of \( a_5+a_4+a_3+a_2+a_1 \).}
{\( -180 \)}{\( -244 \)}{\( -242 \)}{\( -212 \)}{}{}}
\End

\Data\ID{1003032503}
\Updated{2020-07-15 03:07:21}
\Level{C}
\SubArea{Polynomials and fractions}
\LANGen{
\Question{
The polynomial \( 3x^5+px^3-(p-1)x^2+5x-9 \) is divisible by the binomial \( x^2-1 \) if \( p \) is equal to:}
{\( -8 \)}{\( -16 \)}{\( 4 \)}{\( 6 \)}{}{}}
\LANGpl{
\Question{
Wielomian  \( 3x^5+px^3-(p-1)x^2+5x-9 \) jest podzielny przez wielomian \( x^2-1 \), jeśli \( p \) jest równe:}
{\( -8 \)}{\( -16 \)}{\( 4 \)}{\( 6 \)}{}{}}
\LANGcs{
\Question{
Mnohočlen \( 3x^5+px^3-(p-1)x^2+5x-9 \) je dělitelný dvojčlenem \( x^2-1 \), jestliže \( p \) se rovná:}
{\( -8 \)}{\( -16 \)}{\( 4 \)}{\( 6 \)}{}{}}
\LANGsk{
\Question{
Polynóm \( 3x^5+px^3-(p-1)x^2+5x-9 \) je deliteľný dvojčlenom \( x^2-1 \), ak \( p \) je rovné:}
{\( -8 \)}{\( -16 \)}{\( 4 \)}{\( 6 \)}{}{}}
\LANGes{
\Question{
The polynomial \( 3x^5+px^3-(p-1)x^2+5x-9 \) is divisible by the binomial \( x^2-1 \) if \( p \) is equal to:}
{\( -8 \)}{\( -16 \)}{\( 4 \)}{\( 6 \)}{}{}}
\End

\Data\ID{1003032505}
\Updated{2020-07-15 03:07:21}
\Level{C}
\SubArea{Polynomials and fractions}
\LANGen{
\Question{
Factoring the polynomial \( x^4+4 \) you get:}
{\( \left(x^2-2x+2\right)\left(x^2+2x+2\right) \)}{\( \left(x^2+2\right)\left(x^2-(-2)\right) \)}{\( \left( x+\sqrt2 \right)^4 \)}{\( \left(x^2-\sqrt2x+2\right)\left(x^2+\sqrt2x+2\right) \)}{}{}}
\LANGpl{
\Question{
Rozkładając wielomian \( x^4+4 \) na czynniki otrzymamy:}
{\( \left(x^2-2x+2\right)\left(x^2+2x+2\right) \)}{\( \left(x^2+2\right)\left(x^2-(-2)\right) \)}{\( \left( x+\sqrt2 \right)^4 \)}{\( \left(x^2-\sqrt2x+2\right)\left(x^2+\sqrt2x+2\right) \)}{}{}}
\LANGcs{
\Question{
Rozkladem mnohočlenu \( x^4+4 \) dostaneme:}
{\( \left(x^2-2x+2\right)\left(x^2+2x+2\right) \)}{\( \left(x^2+2\right)\left(x^2-(-2)\right) \)}{\( \left( x+\sqrt2 \right)^4 \)}{\( \left(x^2-\sqrt2x+2\right)\left(x^2+\sqrt2x+2\right) \)}{}{}}
\LANGsk{
\Question{
Rozkladom polynómu \( x^4+4 \) na súčin dostaneme:}
{\( \left(x^2-2x+2\right)\left(x^2+2x+2\right) \)}{\( \left(x^2+2\right)\left(x^2-(-2)\right) \)}{\( \left( x+\sqrt2 \right)^4 \)}{\( \left(x^2-\sqrt2x+2\right)\left(x^2+\sqrt2x+2\right) \)}{}{}}
\LANGes{
\Question{
Factoring the polynomial \( x^4+4 \) you get:}
{\( \left(x^2-2x+2\right)\left(x^2+2x+2\right) \)}{\( \left(x^2+2\right)\left(x^2-(-2)\right) \)}{\( \left( x+\sqrt2 \right)^4 \)}{\( \left(x^2-\sqrt2x+2\right)\left(x^2+\sqrt2x+2\right) \)}{}{}}
\End

\Data\ID{1003032506}
\Updated{2020-07-15 03:07:21}
\Level{C}
\SubArea{Polynomials and fractions}
\LANGen{
\Question{
In a parallel circuit the total resistance \( R \) of three components with the resistances \( R_1 \), \( R_2 \), \( R_3 \) is given by the formula \( \frac1R=\frac1{R_1}+\frac1{R_2}+\frac1{R_3} \). Express \( R_1 \) from this formula.}
{\( R_1=\frac{RR_2R_3}{R_2R_3-R(R_2+R_3)} \)}{\( R_1=\frac{R-R_2-R_3}{RR_2R_3} \)}{\( R_1=R-\frac{R_2R_3}{R_2+R_3} \)}{\( R_1=\frac{R_2R_3-R(R_2+R_3)}{RR_2R_3} \)}{}{}}
\LANGpl{
\Question{
W obwodzie równoległym całkowitą rezystancję \( R \) trzech składników z opornością  \( R_1 \), \( R_2 \), \( R_3 \)  określa wzór \( \frac1R=\frac1{R_1}+\frac1{R_2}+\frac1{R_3} \). Wyznacz \( R_1 \) z tego wzoru.}
{\( R_1=\frac{RR_2R_3}{R_2R_3-R(R_2+R_3)} \)}{\( R_1=\frac{R-R_2-R_3}{RR_2R_3} \)}{\( R_1=R-\frac{R_2R_3}{R_2+R_3} \)}{\( R_1=\frac{R_2R_3-R(R_2+R_3)}{RR_2R_3} \)}{}{}}
\LANGcs{
\Question{
V paralelně zapojeném obvodu je celkový odpor \( R \) tří komponentů s odpory \( R_1 \), \( R_2 \), \( R_3 \) dán rovnicí \( \frac1R=\frac1{R_1}+\frac1{R_2}+\frac1{R_3} \). Vyjádřete \( R_1 \) z této rovnice.}
{\( R_1=\frac{RR_2R_3}{R_2R_3-R(R_2+R_3)} \)}{\( R_1=\frac{R-R_2-R_3}{RR_2R_3} \)}{\( R_1=R-\frac{R_2R_3}{R_2+R_3} \)}{\( R_1=\frac{R_2R_3-R(R_2+R_3)}{RR_2R_3} \)}{}{}}
\LANGsk{
\Question{
V paralelnom obvode sú zapojené tri rezistory \( R_1 \), \( R_2 \), \( R_3 \) s celkovým odporom \( R \) dané vzťahom \( \frac1R=\frac1{R_1}+\frac1{R_2}+\frac1{R_3} \). Vyjadrite \( R_1 \) z tohoto vzťahu.}
{\( R_1=\frac{RR_2R_3}{R_2R_3-R(R_2+R_3)} \)}{\( R_1=\frac{R-R_2-R_3}{RR_2R_3} \)}{\( R_1=R-\frac{R_2R_3}{R_2+R_3} \)}{\( R_1=\frac{R_2R_3-R(R_2+R_3)}{RR_2R_3} \)}{}{}}
\LANGes{
\Question{
In a parallel circuit the total resistance \( R \) of three components with the resistances \( R_1 \), \( R_2 \), \( R_3 \) is given by the formula \( \frac1R=\frac1{R_1}+\frac1{R_2}+\frac1{R_3} \). Express \( R_1 \) from this formula.}
{\( R_1=\frac{RR_2R_3}{R_2R_3-R(R_2+R_3)} \)}{\( R_1=\frac{R-R_2-R_3}{RR_2R_3} \)}{\( R_1=R-\frac{R_2R_3}{R_2+R_3} \)}{\( R_1=\frac{R_2R_3-R(R_2+R_3)}{RR_2R_3} \)}{}{}}
\End

\Data\ID{1003032507}
\Updated{2020-11-12 06:11:42}
\Level{C}
\SubArea{Polynomials and fractions}
\LANGen{
\Question{
Express \( v_1 \) from the formula \( v=\frac{v_1v_2(d_1+d_2 )}{d_1v_2+d_2v_1} \).}
{\( v_1=\frac{vd_1 v_2}{v_2d_1+v_2d_2-vd_2} \)}{\( v_1=\frac{vv_2 d_1}{vd_2-v_2d_1-v_2d_2} \)}{\( v_1=\frac{vv_2(d_1+d_2)}{d_1v_2+d_2 v} \)}{\( v_1=\frac{v(d_1v_2+d_2 v_1 )}{v_2 (d_1+d_2 )} \)}{}{}}
\LANGpl{
\Question{
Wyznacz \( v_1 \) ze wzoru \( v=\frac{v_1v_2(d_1+d_2 )}{d_1v_2+d_2v_1} \).}
{\( v_1=\frac{vd_1 v_2}{v_2d_1+v_2d_2-vd_2} \)}{\( v_1=\frac{vv_2 d_1}{vd_2-v_2d_1-v_2d_2} \)}{\( v_1=\frac{vv_2(d_1+d_2)}{d_1v_2+d_2 v} \)}{\( v_1=\frac{v(d_1v_2+d_2 v_1 )}{v_2 (d_1+d_2 )} \)}{}{}}
\LANGcs{
\Question{
Vyjádřete \( v_1 \) z rovnice \( v=\frac{v_1v_2(d_1+d_2 )}{d_1v_2+d_2v_1} \).}
{\( v_1=\frac{vd_1 v_2}{v_2d_1+v_2d_2-vd_2} \)}{\( v_1=\frac{vv_2 d_1}{vd_2-v_2d_1-v_2d_2} \)}{\( v_1=\frac{vv_2(d_1+d_2)}{d_1v_2+d_2 v} \)}{\( v_1=\frac{v(d_1v_2+d_2 v_1 )}{v_2 (d_1+d_2 )} \)}{}{}}
\LANGsk{
\Question{
Vyjadrite neznámu \( v_1 \) zo vzorca \( v=\frac{v_1v_2(d_1+d_2 )}{d_1v_2+d_2v_1} \).}
{\( v_1=\frac{vd_1 v_2}{v_2d_1+v_2d_2-vd_2} \)}{\( v_1=\frac{vv_2 d_1}{vd_2-v_2d_1-v_2d_2} \)}{\( v_1=\frac{vv_2(d_1+d_2)}{d_1v_2+d_2 v} \)}{\( v_1=\frac{v(d_1v_2+d_2 v_1 )}{v_2 (d_1+d_2 )} \)}{}{}}
\LANGes{
\Question{
Express \( v_1 \) from the formula \( v=\frac{v_1v_2(d_1+d_2 )}{d_1v_2+d_2v_1} \).}
{\( v_1=\frac{vd_1 v_2}{v_2d_1+v_2d_2-vd_2} \)}{\( v_1=\frac{vv_2 d_1}{vd_2-v_2d_1-v_2d_2} \)}{\( v_1=\frac{vv_2(d_1+d_2)}{d_1v_2+d_2 v} \)}{\( v_1=\frac{v(d_1v_2+d_2 v_1 )}{v_2 (d_1+d_2 )} \)}{}{}}
\End

\Data\ID{1003032508}
\Updated{2020-07-15 03:07:21}
\Level{C}
\SubArea{Polynomials and fractions}
\LANGen{
\Question{
Find the term that does not contain \( x \) and \( y \) (constant term) in the expansion of \( (x+y)^4(\frac1x+\frac1y)^4 \).}
{\( 70 \)}{\( 2 \)}{\( 36 \)}{\( 0 \)}{}{}}
\LANGpl{
\Question{
Znajdź termin, który nie zawiera \( x \) i \( y \) (stały czas) w ekspansji \( (x+y)^4(\frac1x+\frac1y)^4 \).}
{\( 70 \)}{\( 2 \)}{\( 36 \)}{\( 0 \)}{}{}}
\LANGcs{
\Question{
Určete člen binomického rozvoje výrazu \( (x+y)^4(\frac1x+\frac1y)^4 \), který neobsahuje \( x \) ani \( y \), tj. určete konstantní člen rozvoje.}
{\( 70 \)}{\( 2 \)}{\( 36 \)}{\( 0 \)}{}{}}
\LANGsk{
\Question{
Určte člen binomického rozvoja výrazu \( (x+y)^4(\frac1x+\frac1y)^4 \), ktorý neobsahuje \( x \) ani \( y \), tj. určte konštantný člen rozvoja.}
{\( 70 \)}{\( 2 \)}{\( 36 \)}{\( 0 \)}{}{}}
\LANGes{
\Question{
Find the term that does not contain \( x \) and \( y \) (constant term) in the expansion of \( (x+y)^4(\frac1x+\frac1y)^4 \).}
{\( 70 \)}{\( 2 \)}{\( 36 \)}{\( 0 \)}{}{}}
\End

\Data\ID{9000083606}
\Updated{2021-01-04 12:01:25}
\Level{C}
\SubArea{Polynomials and fractions}
\LANGen{
\Question{
Assuming \(x\neq 2\),
simplify the expression
\(
 \frac{x^{2} + x - 6} 
 {x^{3} - 8}.
\)}
{\(\frac{x+3}
{x^{2}+2x+4}\)}{\(\frac{x+3}
{x^{2}-2x+4}\)}{\(\frac{x+3}
{x^{2}+4x+4}\)}{\(\frac{x+3}
{x^{2}-4}\)}{}{}}
\LANGpl{
\Question{
Zakładając, że \(x\neq 2\),
uprość wyrażenie
\(
 \frac{x^{2} + x - 6} 
 {x^{3} - 8}
\).}
{\(\frac{x+3}
{x^{2}+2x+4}\)}{\(\frac{x+3}
{x^{2}-2x+4}\)}{\(\frac{x+3}
{x^{2}+4x+4}\)}{\(\frac{x+3}
{x^{2}-4}\)}{}{}}
\LANGcs{
\Question{
Za předpokladu, že \(x\neq 2\), zjednodušte výraz \(\frac{x^{2}+x-6}
 {x^{3}-8}\).}
{\(\frac{x+3}
{x^{2}+2x+4}\)}{\(\frac{x+3}
{x^{2}-2x+4}\)}{\(\frac{x+3}
{x^{2}+4x+4}\)}{\(\frac{x+3}
{x^{2}-4}\)}{}{}}
\LANGsk{
\Question{
Za predpokladu, že \(x\neq 2\), zjednodušte výraz \(\frac{x^{2}+x-6}
 {x^{3}-8}\).}
{\(\frac{x+3}
{x^{2}+2x+4}\)}{\(\frac{x+3}
{x^{2}-2x+4}\)}{\(\frac{x+3}
{x^{2}+4x+4}\)}{\(\frac{x+3}
{x^{2}-4}\)}{}{}}
\LANGes{
\Question{
Assuming \(x\neq 2\),
simplify the expression
\(
 \frac{x^{2} + x - 6} 
 {x^{3} - 8}.
\)}
{\(\frac{x+3}
{x^{2}+2x+4}\)}{\(\frac{x+3}
{x^{2}-2x+4}\)}{\(\frac{x+3}
{x^{2}+4x+4}\)}{\(\frac{x+3}
{x^{2}-4}\)}{}{}}
\End

\Data\ID{9000087501}
\Updated{2020-07-15 03:07:37}
\Level{C}
\SubArea{Polynomials and fractions}
\LANGen{
\Question{
Assuming \(x\in \mathbb{R}\setminus \left \{-1\right \}\),
find the quotient of the polynomials. 
\[ 
 (3x^{2} + 2x + 7) : (x + 1)
\]}
{\(3x - 1 + \frac{8} 
{x+1}\)}{\(3x + 2 + \frac{8} 
{x+1}\)}{\(3x - 1 - \frac{5} 
{x+1}\)}{\(3x + 2 - \frac{5} 
{x+1}\)}{}{}}
\LANGpl{
\Question{
Zakładając, że \(x\in \mathbb{R}\setminus \left \{-1\right \}\),
znajdź iloraz wielomianów.
\[ 
 (3x^{2} + 2x + 7) : (x + 1)
\]}
{\(3x - 1 + \frac{8} 
{x+1}\)}{\(3x + 2 + \frac{8} 
{x+1}\)}{\(3x - 1 - \frac{5} 
{x+1}\)}{\(3x + 2 - \frac{5} 
{x+1}\)}{}{}}
\LANGcs{
\Question{
Určete podíl \((3x^{2} + 2x + 7) : (x + 1)\) 
za předpokladu, že \(x\in \mathbb{R}\setminus \left \{-1\right \}\).}
{\(3x - 1 + \frac{8} 
{x+1}\)}{\(3x + 2 + \frac{8} 
{x+1}\)}{\(3x - 1 - \frac{5} 
{x+1}\)}{\(3x + 2 - \frac{5} 
{x+1}\)}{}{}}
\LANGsk{
\Question{
Určte podiel \((3x^{2} + 2x + 7) : (x + 1)\) 
za predpokladu, že \(x\in \mathbb{R}\setminus \left \{-1\right \}\).}
{\(3x - 1 + \frac{8} 
{x+1}\)}{\(3x + 2 + \frac{8} 
{x+1}\)}{\(3x - 1 - \frac{5} 
{x+1}\)}{\(3x + 2 - \frac{5} 
{x+1}\)}{}{}}
\LANGes{
\Question{
Assuming \(x\in \mathbb{R}\setminus \left \{-1\right \}\),
find the quotient of the polynomials. 
\[ 
 (3x^{2} + 2x + 7) : (x + 1)
\]}
{\(3x - 1 + \frac{8} 
{x+1}\)}{\(3x + 2 + \frac{8} 
{x+1}\)}{\(3x - 1 - \frac{5} 
{x+1}\)}{\(3x + 2 - \frac{5} 
{x+1}\)}{}{}}
\End

\Data\ID{9000087502}
\Updated{2021-01-04 12:01:59}
\Level{C}
\SubArea{Polynomials and fractions}
\LANGen{
\Question{
Assuming \(x\in \mathbb{R}\setminus \left \{\pm 1\right \}\),
find the quotient of the polynomials: 
\[ 
 (-2x^{4} - 3x^{2} + 3) : (x^{2} - 1)
\]}
{\(- 2x^{2} - 5 - \frac{2} 
{x^{2}-1}\)}{\(- 2x^{2} - 5 + \frac{2} 
{x^{2}-1}\)}{\(2x^{2} + 5 - \frac{2} 
{x^{2}-1}\)}{\(2x^{2} + 5 + \frac{2} 
{x^{2}-1}\)}{}{}}
\LANGpl{
\Question{
Zakładając, że  \(x\in \mathbb{R}\setminus \left \{\pm 1\right \}\),
znajdź iloraz wielomianów: 
\[ 
 (-2x^{4} - 3x^{2} + 3) : (x^{2} - 1)
\]}
{\(- 2x^{2} - 5 - \frac{2} 
{x^{2}-1}\)}{\(- 2x^{2} - 5 + \frac{2} 
{x^{2}-1}\)}{\(2x^{2} + 5 - \frac{2} 
{x^{2}-1}\)}{\(2x^{2} + 5 + \frac{2} 
{x^{2}-1}\)}{}{}}
\LANGcs{
\Question{
Určete podíl \((-2x^{4} - 3x^{2} + 3) : (x^{2} - 1)\) 
za předpokladu, že \(x\in \mathbb{R}\setminus \left \{\pm 1\right \}\).}
{\(- 2x^{2} - 5 - \frac{2} 
{x^{2}-1}\)}{\(- 2x^{2} - 5 + \frac{2} 
{x^{2}-1}\)}{\(2x^{2} + 5 - \frac{2} 
{x^{2}-1}\)}{\(2x^{2} + 5 + \frac{2} 
{x^{2}-1}\)}{}{}}
\LANGsk{
\Question{
Určte podiel \((-2x^{4} - 3x^{2} + 3) : (x^{2} - 1)\) 
za predpokladu, že \(x\in \mathbb{R}\setminus \left \{\pm 1\right \}\).}
{\(- 2x^{2} - 5 - \frac{2} 
{x^{2}-1}\)}{\(- 2x^{2} - 5 + \frac{2} 
{x^{2}-1}\)}{\(2x^{2} + 5 - \frac{2} 
{x^{2}-1}\)}{\(2x^{2} + 5 + \frac{2} 
{x^{2}-1}\)}{}{}}
\LANGes{
\Question{
Assuming \(x\in \mathbb{R}\setminus \left \{\pm 1\right \}\),
find the quotient of the polynomials: 
\[ 
 (-2x^{4} - 3x^{2} + 3) : (x^{2} - 1)
\]}
{\(- 2x^{2} - 5 - \frac{2} 
{x^{2}-1}\)}{\(- 2x^{2} - 5 + \frac{2} 
{x^{2}-1}\)}{\(2x^{2} + 5 - \frac{2} 
{x^{2}-1}\)}{\(2x^{2} + 5 + \frac{2} 
{x^{2}-1}\)}{}{}}
\End

\Data\ID{9000087503}
\Updated{2021-01-04 01:01:26}
\Level{C}
\SubArea{Polynomials and fractions}
\LANGen{
\Question{
Assuming \(x\in \mathbb{R}\setminus \left \{-\frac{3}
{2}\right \}\),
find the quotient of the polynomials: 
\[ 
 (x^{2} + x + 1) : (2x + 3)
\]}
{\(\frac{1}
{2}x -\frac{1} 
{4} + \frac{\frac{7} 
{4} } 
{2x+3}\)}{\(\frac{1}
{2}x -\frac{1} 
{2} + \frac{\frac{7} 
{4} } 
{2x+3}\)}{\(x + 2 + \frac{7} 
{2x+3}\)}{\(x - 2 + \frac{7} 
{2x+3}\)}{}{}}
\LANGpl{
\Question{
Zakładając, że  \(x\in \mathbb{R}\setminus \left \{-\frac{3}
{2}\right \}\),
znajdź iloraz wielomianów:
\[ 
 (x^{2} + x + 1) : (2x + 3)
\]}
{\(\frac{1}
{2}x -\frac{1} 
{4} + \frac{\frac{7} 
{4} } 
{2x+3}\)}{\(\frac{1}
{2}x -\frac{1} 
{2} + \frac{\frac{7} 
{4} } 
{2x+3}\)}{\(x + 2 + \frac{7} 
{2x+3}\)}{\(x - 2 + \frac{7} 
{2x+3}\)}{}{}}
\LANGcs{
\Question{
Určete podíl \((x^{2} + x + 1) : (2x + 3)\) 
za předpokladu, že \(x\in \mathbb{R}\setminus \left \{-\frac{3}
{2}\right \}\).}
{\(\frac{1}
{2}x -\frac{1} 
{4} + \frac{\frac{7} 
{4} } 
{2x+3}\)}{\(\frac{1}
{2}x -\frac{1} 
{2} + \frac{\frac{7} 
{4} } 
{2x+3}\)}{\(x + 2 + \frac{7} 
{2x+3}\)}{\(x - 2 + \frac{7} 
{2x+3}\)}{}{}}
\LANGsk{
\Question{
Určte podiel \((x^{2} + x + 1) : (2x + 3)\) 
za predpokladu, že \(x\in \mathbb{R}\setminus \left \{-\frac{3}
{2}\right \}\).}
{\(\frac{1}
{2}x -\frac{1} 
{4} + \frac{\frac{7} 
{4} } 
{2x+3}\)}{\(\frac{1}
{2}x -\frac{1} 
{2} + \frac{\frac{7} 
{4} } 
{2x+3}\)}{\(x + 2 + \frac{7} 
{2x+3}\)}{\(x - 2 + \frac{7} 
{2x+3}\)}{}{}}
\LANGes{
\Question{
Assuming \(x\in \mathbb{R}\setminus \left \{-\frac{3}
{2}\right \}\),
find the quotient of the polynomials: 
\[ 
 (x^{2} + x + 1) : (2x + 3)
\]}
{\(\frac{1}
{2}x -\frac{1} 
{4} + \frac{\frac{7} 
{4} } 
{2x+3}\)}{\(\frac{1}
{2}x -\frac{1} 
{2} + \frac{\frac{7} 
{4} } 
{2x+3}\)}{\(x + 2 + \frac{7} 
{2x+3}\)}{\(x - 2 + \frac{7} 
{2x+3}\)}{}{}}
\End

\Data\ID{9000087504}
\Updated{2020-07-15 03:07:37}
\Level{C}
\SubArea{Polynomials and fractions}
\LANGen{
\Question{
Assuming \(x\in \mathbb{R}\setminus \left \{-\frac{3}
{5}\right \}\),
find the quotient of the polynomials. 
\[ 
 (5x^{3} - 2x^{2} + x + 1) : (5x + 3)
\]}
{\(x^{2} - x + \frac{4} 
{5} - \frac{\frac{7} 
{5} } 
{5x+3}\)}{\(x^{2} - x + \frac{4} 
{5} + \frac{\frac{7} 
{5} } 
{5x+3}\)}{\(x^{2} - x + \frac{4} 
{5} - \frac{\frac{9} 
{5} } 
{5x+3}\)}{\(x^{2} - x + \frac{4} 
{5} + \frac{\frac{9} 
{5} } 
{5x+3}\)}{}{}}
\LANGpl{
\Question{
Zakładając, że  \(x\in \mathbb{R}\setminus \left \{-\frac{3}
{5}\right \}\),
znajdź iloraz wielomianów.
\[ 
 (5x^{3} - 2x^{2} + x + 1) : (5x + 3)
\]}
{\(x^{2} - x + \frac{4} 
{5} - \frac{\frac{7} 
{5} } 
{5x+3}\)}{\(x^{2} - x + \frac{4} 
{5} + \frac{\frac{7} 
{5} } 
{5x+3}\)}{\(x^{2} - x + \frac{4} 
{5} - \frac{\frac{9} 
{5} } 
{5x+3}\)}{\(x^{2} - x + \frac{4} 
{5} + \frac{\frac{9} 
{5} } 
{5x+3}\)}{}{}}
\LANGcs{
\Question{
Určete podíl \((5x^{3} - 2x^{2} + x + 1) : (5x + 3)\) 
za předpokladu, že \(x\in \mathbb{R}\setminus \left \{-\frac{3}
{5}\right \}\).}
{\(x^{2} - x + \frac{4} 
{5} - \frac{\frac{7} 
{5} } 
{5x+3}\)}{\(x^{2} - x + \frac{4} 
{5} + \frac{\frac{7} 
{5} } 
{5x+3}\)}{\(x^{2} - x + \frac{4} 
{5} - \frac{\frac{9} 
{5} } 
{5x+3}\)}{\(x^{2} - x + \frac{4} 
{5} + \frac{\frac{9} 
{5} } 
{5x+3}\)}{}{}}
\LANGsk{
\Question{
Určte podiel \((5x^{3} - 2x^{2} + x + 1) : (5x + 3)\) 
za predpokladu, že \(x\in \mathbb{R}\setminus \left \{-\frac{3}
{5}\right \}\).}
{\(x^{2} - x + \frac{4} 
{5} - \frac{\frac{7} 
{5} } 
{5x+3}\)}{\(x^{2} - x + \frac{4} 
{5} + \frac{\frac{7} 
{5} } 
{5x+3}\)}{\(x^{2} - x + \frac{4} 
{5} - \frac{\frac{9} 
{5} } 
{5x+3}\)}{\(x^{2} - x + \frac{4} 
{5} + \frac{\frac{9} 
{5} } 
{5x+3}\)}{}{}}
\LANGes{
\Question{
Assuming \(x\in \mathbb{R}\setminus \left \{-\frac{3}
{5}\right \}\),
find the quotient of the polynomials. 
\[ 
 (5x^{3} - 2x^{2} + x + 1) : (5x + 3)
\]}
{\(x^{2} - x + \frac{4} 
{5} - \frac{\frac{7} 
{5} } 
{5x+3}\)}{\(x^{2} - x + \frac{4} 
{5} + \frac{\frac{7} 
{5} } 
{5x+3}\)}{\(x^{2} - x + \frac{4} 
{5} - \frac{\frac{9} 
{5} } 
{5x+3}\)}{\(x^{2} - x + \frac{4} 
{5} + \frac{\frac{9} 
{5} } 
{5x+3}\)}{}{}}
\End

\Data\ID{9000087505}
\Updated{2020-07-15 03:07:37}
\Level{C}
\SubArea{Polynomials and fractions}
\LANGen{
\Question{
Assuming \(x\in \mathbb{R}\setminus \left \{-\frac{1}
{2}\right \}\),
find the quotient of the polynomials. 
\[ 
 (4x^{3} - 1) : (2x + 1)
\]}
{\(2x^{2} - x + \frac{1} 
{2} - \frac{\frac{3} 
{2} } 
{2x+1}\)}{\(2x^{2} + x + \frac{1} 
{2} - \frac{\frac{3} 
{2} } 
{2x+1}\)}{\(2x^{2} - x + \frac{1} 
{4} - \frac{\frac{3} 
{2} } 
{2x+1}\)}{\(2x^{2} + x + \frac{1} 
{4} - \frac{\frac{3} 
{2} } 
{2x+1}\)}{}{}}
\LANGpl{
\Question{
Zakładając, że  \(x\in \mathbb{R}\setminus \left \{-\frac{1}
{2}\right \}\),
znajdź iloraz wielomianów.
\[ 
 (4x^{3} - 1) : (2x + 1)
\]}
{\(2x^{2} - x + \frac{1} 
{2} - \frac{\frac{3} 
{2} } 
{2x+1}\)}{\(2x^{2} + x + \frac{1} 
{2} - \frac{\frac{3} 
{2} } 
{2x+1}\)}{\(2x^{2} - x + \frac{1} 
{4} - \frac{\frac{3} 
{2} } 
{2x+1}\)}{\(2x^{2} + x + \frac{1} 
{4} - \frac{\frac{3} 
{2} } 
{2x+1}\)}{}{}}
\LANGcs{
\Question{
Určete podíl \((4x^{3} - 1) : (2x + 1)\) 
za předpokladu, že \(x\in \mathbb{R}\setminus \left \{-\frac{1}
{2}\right \}\).}
{\(2x^{2} - x + \frac{1} 
{2} - \frac{\frac{3} 
{2} } 
{2x+1}\)}{\(2x^{2} + x + \frac{1} 
{2} - \frac{\frac{3} 
{2} } 
{2x+1}\)}{\(2x^{2} - x + \frac{1} 
{4} - \frac{\frac{3} 
{2} } 
{2x+1}\)}{\(2x^{2} + x + \frac{1} 
{4} - \frac{\frac{3} 
{2} } 
{2x+1}\)}{}{}}
\LANGsk{
\Question{
Určte podiel \((4x^{3} - 1) : (2x + 1)\) 
za predpokladu, že \(x\in \mathbb{R}\setminus \left \{-\frac{1}
{2}\right \}\).}
{\(2x^{2} - x + \frac{1} 
{2} - \frac{\frac{3} 
{2} } 
{2x+1}\)}{\(2x^{2} + x + \frac{1} 
{2} - \frac{\frac{3} 
{2} } 
{2x+1}\)}{\(2x^{2} - x + \frac{1} 
{4} - \frac{\frac{3} 
{2} } 
{2x+1}\)}{\(2x^{2} + x + \frac{1} 
{4} - \frac{\frac{3} 
{2} } 
{2x+1}\)}{}{}}
\LANGes{
\Question{
Assuming \(x\in \mathbb{R}\setminus \left \{-\frac{1}
{2}\right \}\),
find the quotient of the polynomials. 
\[ 
 (4x^{3} - 1) : (2x + 1)
\]}
{\(2x^{2} - x + \frac{1} 
{2} - \frac{\frac{3} 
{2} } 
{2x+1}\)}{\(2x^{2} + x + \frac{1} 
{2} - \frac{\frac{3} 
{2} } 
{2x+1}\)}{\(2x^{2} - x + \frac{1} 
{4} - \frac{\frac{3} 
{2} } 
{2x+1}\)}{\(2x^{2} + x + \frac{1} 
{4} - \frac{\frac{3} 
{2} } 
{2x+1}\)}{}{}}
\End

\Data\ID{9000087506}
\Updated{2020-07-15 03:07:37}
\Level{C}
\SubArea{Polynomials and fractions}
\LANGen{
\Question{
Assuming \(x\in \mathbb{R}\setminus \left \{1\right \}\),
find the quotient of the polynomials. 
\[ 
 (2x + 2x^{2} - 3) : (x - 1)
\]}
{\(2x + 4 + \frac{1} 
{x-1}\)}{\(2x + 4 + \frac{2} 
{x-1}\)}{\(2x + 2 + \frac{1} 
{x-1}\)}{\(2x + 2 + \frac{2} 
{x-1}\)}{}{}}
\LANGpl{
\Question{
Zakładając, że  \(x\in \mathbb{R}\setminus \left \{1\right \}\),
znajdź iloraz wielomianów.
\[ 
 (2x + 2x^{2} - 3) : (x - 1)
\]}
{\(2x + 4 + \frac{1} 
{x-1}\)}{\(2x + 4 + \frac{2} 
{x-1}\)}{\(2x + 2 + \frac{1} 
{x-1}\)}{\(2x + 2 + \frac{2} 
{x-1}\)}{}{}}
\LANGcs{
\Question{
Určete podíl \((2x + 2x^{2} - 3) : (x - 1)\) 
za předpokladu, že \(x\in \mathbb{R}\setminus \left \{1\right \}\).}
{\(2x + 4 + \frac{1} 
{x-1}\)}{\(2x + 4 + \frac{2} 
{x-1}\)}{\(2x + 2 + \frac{1} 
{x-1}\)}{\(2x + 2 + \frac{2} 
{x-1}\)}{}{}}
\LANGsk{
\Question{
Určte podiel \((2x + 2x^{2} - 3) : (x - 1)\) 
za predpokladu, že \(x\in \mathbb{R}\setminus \left \{1\right \}\).}
{\(2x + 4 + \frac{1} 
{x-1}\)}{\(2x + 4 + \frac{2} 
{x-1}\)}{\(2x + 2 + \frac{1} 
{x-1}\)}{\(2x + 2 + \frac{2} 
{x-1}\)}{}{}}
\LANGes{
\Question{
Assuming \(x\in \mathbb{R}\setminus \left \{1\right \}\),
find the quotient of the polynomials. 
\[ 
 (2x + 2x^{2} - 3) : (x - 1)
\]}
{\(2x + 4 + \frac{1} 
{x-1}\)}{\(2x + 4 + \frac{2} 
{x-1}\)}{\(2x + 2 + \frac{1} 
{x-1}\)}{\(2x + 2 + \frac{2} 
{x-1}\)}{}{}}
\End

\Data\ID{9000087507}
\Updated{2020-07-15 03:07:37}
\Level{C}
\SubArea{Polynomials and fractions}
\LANGen{
\Question{
Assuming \(x\in \mathbb{R}\),
find the quotient of two polynomials. 
\[ 
 (-x^{3} - x^{2} + x - 1) : (x^{2} + 1)
\]}
{\(- x - 1 + \frac{2x} 
{x^{2}+1}\)}{\(- x - 1 + \frac{x} 
{x^{2}+1}\)}{\(x - 1 + \frac{x} 
{x^{2}+1}\)}{\(x - 1 + \frac{2x} 
{x^{2}+1}\)}{}{}}
\LANGpl{
\Question{
Zakładając, że \(x\in \mathbb{R}\),
znajdź iloraz wielomianów.
\[ 
 (-x^{3} - x^{2} + x - 1) : (x^{2} + 1)
\]}
{\(- x - 1 + \frac{2x} 
{x^{2}+1}\)}{\(- x - 1 + \frac{x} 
{x^{2}+1}\)}{\(x - 1 + \frac{x} 
{x^{2}+1}\)}{\(x - 1 + \frac{2x} 
{x^{2}+1}\)}{}{}}
\LANGcs{
\Question{
Určete podíl \((-x^{3} - x^{2} + x - 1) : (x^{2} + 1)\), 
je-li \(x\in \mathbb{R}\).}
{\(- x - 1 + \frac{2x} 
{x^{2}+1}\)}{\(- x - 1 + \frac{x} 
{x^{2}+1}\)}{\(x - 1 + \frac{x} 
{x^{2}+1}\)}{\(x - 1 + \frac{2x} 
{x^{2}+1}\)}{}{}}
\LANGsk{
\Question{
Určte podiel \((-x^{3} - x^{2} + x - 1) : (x^{2} + 1)\), ak je \(x\in \mathbb{R}\).}
{\(- x - 1 + \frac{2x} 
{x^{2}+1}\)}{\(- x - 1 + \frac{x} 
{x^{2}+1}\)}{\(x - 1 + \frac{x} 
{x^{2}+1}\)}{\(x - 1 + \frac{2x} 
{x^{2}+1}\)}{}{}}
\LANGes{
\Question{
Assuming \(x\in \mathbb{R}\),
find the quotient of two polynomials. 
\[ 
 (-x^{3} - x^{2} + x - 1) : (x^{2} + 1)
\]}
{\(- x - 1 + \frac{2x} 
{x^{2}+1}\)}{\(- x - 1 + \frac{x} 
{x^{2}+1}\)}{\(x - 1 + \frac{x} 
{x^{2}+1}\)}{\(x - 1 + \frac{2x} 
{x^{2}+1}\)}{}{}}
\End

\Data\ID{9000087508}
\Updated{2020-07-15 03:07:37}
\Level{C}
\SubArea{Polynomials and fractions}
\LANGen{
\Question{
Assuming \(x\in \mathbb{R}\setminus \left \{0, 1, 3\right \}\),
find the quotient of the polynomials. 
\[ 
 (-5x^{4} + 4x^{2} + 3x - 4) : (x^{3} - 4x^{2} + 3x)
\]}
{\(- 5x - 20 + \frac{-61x^{2}+63x-4} 
 {x^{3}-4x^{2}+3x} \)}{\(- 5x - 20 + \frac{16x^{2}+23x+36} 
 {x^{3}-4x^{2}+3x} \)}{\(- 5x - 10 + \frac{-61x^{2}+63x-4} 
 {x^{3}-4x^{2}+3x} \)}{\(- 5x - 10 + \frac{-16x^{2}+23x-36} 
 {x^{3}-4x^{2}+3x} \)}{}{}}
\LANGpl{
\Question{
Zakładając, że  \(x\in \mathbb{R}\setminus \left \{0, 1, 3\right \}\),
znajdź iloraz wielomianów.
\[ 
 (-5x^{4} + 4x^{2} + 3x - 4) : (x^{3} - 4x^{2} + 3x)
\]}
{\(- 5x - 20 + \frac{-61x^{2}+63x-4} 
 {x^{3}-4x^{2}+3x} \)}{\(- 5x - 20 + \frac{16x^{2}+23x+36} 
 {x^{3}-4x^{2}+3x} \)}{\(- 5x - 10 + \frac{-61x^{2}+63x-4} 
 {x^{3}-4x^{2}+3x} \)}{\(- 5x - 10 + \frac{-16x^{2}+23x-36} 
 {x^{3}-4x^{2}+3x} \)}{}{}}
\LANGcs{
\Question{
Určete podíl \((-5x^{4} + 4x^{2} + 3x - 4) : (x^{3} - 4x^{2} + 3x)\) 
za předpokladu, že \(x\in \mathbb{R}\setminus \left \{0, 1, 3\right \}\).}
{\(- 5x - 20 + \frac{-61x^{2}+63x-4} 
 {x^{3}-4x^{2}+3x} \)}{\(- 5x - 20 + \frac{16x^{2}+23x+36} 
 {x^{3}-4x^{2}+3x} \)}{\(- 5x - 10 + \frac{-61x^{2}+63x-4} 
 {x^{3}-4x^{2}+3x} \)}{\(- 5x - 10 + \frac{-16x^{2}+23x-36} 
 {x^{3}-4x^{2}+3x} \)}{}{}}
\LANGsk{
\Question{
Určte podiel \((-5x^{4} + 4x^{2} + 3x - 4) : (x^{3} - 4x^{2} + 3x)\) 
za predpokladu, že \(x\in \mathbb{R}\setminus \left \{0, 1, 3\right \}\).}
{\(- 5x - 20 + \frac{-61x^{2}+63x-4} 
 {x^{3}-4x^{2}+3x} \)}{\(- 5x - 20 + \frac{16x^{2}+23x+36} 
 {x^{3}-4x^{2}+3x} \)}{\(- 5x - 10 + \frac{-61x^{2}+63x-4} 
 {x^{3}-4x^{2}+3x} \)}{\(- 5x - 10 + \frac{-16x^{2}+23x-36} 
 {x^{3}-4x^{2}+3x} \)}{}{}}
\LANGes{
\Question{
Assuming \(x\in \mathbb{R}\setminus \left \{0, 1, 3\right \}\),
find the quotient of the polynomials. 
\[ 
 (-5x^{4} + 4x^{2} + 3x - 4) : (x^{3} - 4x^{2} + 3x)
\]}
{\(- 5x - 20 + \frac{-61x^{2}+63x-4} 
 {x^{3}-4x^{2}+3x} \)}{\(- 5x - 20 + \frac{16x^{2}+23x+36} 
 {x^{3}-4x^{2}+3x} \)}{\(- 5x - 10 + \frac{-61x^{2}+63x-4} 
 {x^{3}-4x^{2}+3x} \)}{\(- 5x - 10 + \frac{-16x^{2}+23x-36} 
 {x^{3}-4x^{2}+3x} \)}{}{}}
\End

\Data\ID{9000101609}
\Updated{2020-07-15 03:07:37}
\Level{C}
\SubArea{Polynomials and fractions}
\LANGen{
\Question{
Simplify using long division of polynomials. 
\[ 
 \left (3x^{3} + 17x^{2} + 23x + 5\right ) : \left (x^{2} + 4x + 1\right )
\]}
{\(3x + 5\)}{\(3x - 5\)}{\(3x + 1\)}{\(3x - 1\)}{}{}}
\LANGpl{
\Question{
Podziel wielomiany.
\[ 
 \left (3x^{3} + 17x^{2} + 23x + 5\right ) : \left (x^{2} + 4x + 1\right )
\]}
{\(3x + 5\)}{\(3x - 5\)}{\(3x + 1\)}{\(3x - 1\)}{}{}}
\LANGcs{
\Question{
Dělením \(\left (3x^{3} + 17x^{2} + 23x + 5\right ) : \left (x^{2} + 4x + 1\right )\) 
získáme výraz:}
{\(3x + 5\)}{\(3x - 5\)}{\(3x + 1\)}{\(3x - 1\)}{}{}}
\LANGsk{
\Question{
Určte podiel polynómov.
\[ 
 \left (3x^{3} + 17x^{2} + 23x + 5\right ) : \left (x^{2} + 4x + 1\right )
\]}
{\(3x + 5\)}{\(3x - 5\)}{\(3x + 1\)}{\(3x - 1\)}{}{}}
\LANGes{
\Question{
Simplify using long division of polynomials. 
\[ 
 \left (3x^{3} + 17x^{2} + 23x + 5\right ) : \left (x^{2} + 4x + 1\right )
\]}
{\(3x + 5\)}{\(3x - 5\)}{\(3x + 1\)}{\(3x - 1\)}{}{}}
\End

\Data\ID{9000101610}
\Updated{2020-07-15 03:07:37}
\Level{C}
\SubArea{Polynomials and fractions}
\LANGen{
\Question{
Simplify using long division of polynomials. 
\[ 
 \left (2x^{3} - x^{2} - 3x - 1\right ) : \left (2x + 1\right )
\]}
{\(x^{2} - x - 1\)}{\(x^{2} - x + 1\)}{\(x^{2} + x + 1\)}{\(x^{2} - 2x - 1\)}{}{}}
\LANGpl{
\Question{
Podziel wielomiany.
\[ 
 \left (2x^{3} - x^{2} - 3x - 1\right ) : \left (2x + 1\right )
\]}
{\(x^{2} - x - 1\)}{\(x^{2} - x + 1\)}{\(x^{2} + x + 1\)}{\(x^{2} - 2x - 1\)}{}{}}
\LANGcs{
\Question{
Dělením \(\left (2x^{3} - x^{2} - 3x - 1\right ) : \left (2x + 1\right )\) 
získáme výraz:}
{\(x^{2} - x - 1\)}{\(x^{2} - x + 1\)}{\(x^{2} + x + 1\)}{\(x^{2} - 2x - 1\)}{}{}}
\LANGsk{
\Question{
Určte podiel polynómov. 
\[ 
 \left (2x^{3} - x^{2} - 3x - 1\right ) : \left (2x + 1\right )
\]}
{\(x^{2} - x - 1\)}{\(x^{2} - x + 1\)}{\(x^{2} + x + 1\)}{\(x^{2} - 2x - 1\)}{}{}}
\LANGes{
\Question{
Simplify using long division of polynomials. 
\[ 
 \left (2x^{3} - x^{2} - 3x - 1\right ) : \left (2x + 1\right )
\]}
{\(x^{2} - x - 1\)}{\(x^{2} - x + 1\)}{\(x^{2} + x + 1\)}{\(x^{2} - 2x - 1\)}{}{}}
\End

\Data\ID{9000101707}
\Updated{2020-07-15 03:07:37}
\Level{C}
\SubArea{Polynomials and fractions}
\LANGen{
\Question{
Factor the following polynomial. 
\[ 
 x^{6} - 1
\]}
{\(\left (x - 1\right )\left (x + 1\right )\left (x^{2} + x + 1\right )\left (x^{2} - x + 1\right )\)}{\(\left (x - 1\right )\left (x + 1\right )\left (x^{2} + x + 1\right )\left (x^{2} - x - 1\right )\)}{\(\left (x - 1\right )\left (x + 1\right )\left (x^{2} + 2x + 1\right )\left (x^{2} - 2x + 1\right )\)}{\(\left (x - 1\right )\left (x + 1\right )\left (x^{2} + x - 1\right )\left (x^{2} - x + 1\right )\)}{}{}}
\LANGpl{
\Question{
Rozłóż na czynniki podany wielomian:
\[ 
 x^{6} - 1
\]}
{\(\left (x - 1\right )\left (x + 1\right )\left (x^{2} + x + 1\right )\left (x^{2} - x + 1\right )\)}{\(\left (x - 1\right )\left (x + 1\right )\left (x^{2} + x + 1\right )\left (x^{2} - x - 1\right )\)}{\(\left (x - 1\right )\left (x + 1\right )\left (x^{2} + 2x + 1\right )\left (x^{2} - 2x + 1\right )\)}{\(\left (x - 1\right )\left (x + 1\right )\left (x^{2} + x - 1\right )\left (x^{2} - x + 1\right )\)}{}{}}
\LANGcs{
\Question{
Rozložením výrazu \(x^{6} - 1\) 
získáme výraz:}
{\(\left (x - 1\right )\left (x + 1\right )\left (x^{2} + x + 1\right )\left (x^{2} - x + 1\right )\)}{\(\left (x - 1\right )\left (x + 1\right )\left (x^{2} + x + 1\right )\left (x^{2} - x - 1\right )\)}{\(\left (x - 1\right )\left (x + 1\right )\left (x^{2} + 2x + 1\right )\left (x^{2} - 2x + 1\right )\)}{\(\left (x - 1\right )\left (x + 1\right )\left (x^{2} + x - 1\right )\left (x^{2} - x + 1\right )\)}{}{}}
\LANGsk{
\Question{
Upravte na súčin. 
\[ 
 x^{6} - 1
\]}
{\(\left (x - 1\right )\left (x + 1\right )\left (x^{2} + x + 1\right )\left (x^{2} - x + 1\right )\)}{\(\left (x - 1\right )\left (x + 1\right )\left (x^{2} + x + 1\right )\left (x^{2} - x - 1\right )\)}{\(\left (x - 1\right )\left (x + 1\right )\left (x^{2} + 2x + 1\right )\left (x^{2} - 2x + 1\right )\)}{\(\left (x - 1\right )\left (x + 1\right )\left (x^{2} + x - 1\right )\left (x^{2} - x + 1\right )\)}{}{}}
\LANGes{
\Question{
Factor the following polynomial. 
\[ 
 x^{6} - 1
\]}
{\(\left (x - 1\right )\left (x + 1\right )\left (x^{2} + x + 1\right )\left (x^{2} - x + 1\right )\)}{\(\left (x - 1\right )\left (x + 1\right )\left (x^{2} + x + 1\right )\left (x^{2} - x - 1\right )\)}{\(\left (x - 1\right )\left (x + 1\right )\left (x^{2} + 2x + 1\right )\left (x^{2} - 2x + 1\right )\)}{\(\left (x - 1\right )\left (x + 1\right )\left (x^{2} + x - 1\right )\left (x^{2} - x + 1\right )\)}{}{}}
\End

\Data\ID{9000101708}
\Updated{2020-07-15 03:07:37}
\Level{C}
\SubArea{Polynomials and fractions}
\LANGen{
\Question{
Factor the following polynomial. 
\[ 
 8x^{3} - 27
\]}
{\(\left (2x - 3\right )\left (4x^{2} + 6x + 9\right )\)}{\(\left (2x - 3\right )\left (4x^{2} - 6x + 9\right )\)}{\(\left (2x + 9\right )\left (4x^{2} - 6x + 9\right )\)}{\(\left (2x - 3\right )\left (4x^{2} + 6x - 9\right )\)}{}{}}
\LANGpl{
\Question{
Rozłóż na czynniki podany wielomian:
\[ 
 8x^{3} - 27
\]}
{\(\left (2x - 3\right )\left (4x^{2} + 6x + 9\right )\)}{\(\left (2x - 3\right )\left (4x^{2} - 6x + 9\right )\)}{\(\left (2x + 9\right )\left (4x^{2} - 6x + 9\right )\)}{\(\left (2x - 3\right )\left (4x^{2} + 6x - 9\right )\)}{}{}}
\LANGcs{
\Question{
Rozložením výrazu \(8x^{3} - 27\) 
získáme výraz:}
{\(\left (2x - 3\right )\left (4x^{2} + 6x + 9\right )\)}{\(\left (2x - 3\right )\left (4x^{2} - 6x + 9\right )\)}{\(\left (2x + 9\right )\left (4x^{2} - 6x + 9\right )\)}{\(\left (2x - 3\right )\left (4x^{2} + 6x - 9\right )\)}{}{}}
\LANGsk{
\Question{
Upravte na súčin. 
\[ 
 8x^{3} - 27
\]}
{\(\left (2x - 3\right )\left (4x^{2} + 6x + 9\right )\)}{\(\left (2x - 3\right )\left (4x^{2} - 6x + 9\right )\)}{\(\left (2x + 9\right )\left (4x^{2} - 6x + 9\right )\)}{\(\left (2x - 3\right )\left (4x^{2} + 6x - 9\right )\)}{}{}}
\LANGes{
\Question{
Factor the following polynomial. 
\[ 
 8x^{3} - 27
\]}
{\(\left (2x - 3\right )\left (4x^{2} + 6x + 9\right )\)}{\(\left (2x - 3\right )\left (4x^{2} - 6x + 9\right )\)}{\(\left (2x + 9\right )\left (4x^{2} - 6x + 9\right )\)}{\(\left (2x - 3\right )\left (4x^{2} + 6x - 9\right )\)}{}{}}
\End

\Data\ID{9000101709}
\Updated{2020-07-15 03:07:37}
\Level{C}
\SubArea{Polynomials and fractions}
\LANGen{
\Question{
Factor the following polynomial. 
\[ 
 27x^{6}z - 8y^{3}z
\]}
{\(z\left (3x^{2} - 2y\right )\left (9x^{4} + 6x^{2}y + 4y^{2}\right )\)}{\(z\left (3x^{2} + 2y\right )\left (9x^{4} + 6x^{2}y - 4y^{2}\right )\)}{\(z\left (3x^{2} + 2y\right )\left (9x^{4} - 6x^{2}y + 4y^{2}\right )\)}{\(z\left (3x^{2} - 2y\right )\left (9x^{4} + 6x^{2}y^{2} + 4y\right )\)}{}{}}
\LANGpl{
\Question{
Rozłóż na czynniki podany wielomian: 
\[ 
 27x^{6}z - 8y^{3}z
\]}
{\(z\left (3x^{2} - 2y\right )\left (9x^{4} + 6x^{2}y + 4y^{2}\right )\)}{\(z\left (3x^{2} + 2y\right )\left (9x^{4} + 6x^{2}y - 4y^{2}\right )\)}{\(z\left (3x^{2} + 2y\right )\left (9x^{4} - 6x^{2}y + 4y^{2}\right )\)}{\(z\left (3x^{2} - 2y\right )\left (9x^{4} + 6x^{2}y^{2} + 4y\right )\)}{}{}}
\LANGcs{
\Question{
Rozložením výrazu \(27x^{6}z - 8y^{3}z\) 
získáme výraz:}
{\(z\left (3x^{2} - 2y\right )\left (9x^{4} + 6x^{2}y + 4y^{2}\right )\)}{\(z\left (3x^{2} + 2y\right )\left (9x^{4} + 6x^{2}y - 4y^{2}\right )\)}{\(z\left (3x^{2} + 2y\right )\left (9x^{4} - 6x^{2}y + 4y^{2}\right )\)}{\(z\left (3x^{2} - 2y\right )\left (9x^{4} + 6x^{2}y^{2} + 4y\right )\)}{}{}}
\LANGsk{
\Question{
Upravte na súčin. 
\[ 
 27x^{6}z - 8y^{3}z
\]}
{\(z\left (3x^{2} - 2y\right )\left (9x^{4} + 6x^{2}y + 4y^{2}\right )\)}{\(z\left (3x^{2} + 2y\right )\left (9x^{4} + 6x^{2}y - 4y^{2}\right )\)}{\(z\left (3x^{2} + 2y\right )\left (9x^{4} - 6x^{2}y + 4y^{2}\right )\)}{\(z\left (3x^{2} - 2y\right )\left (9x^{4} + 6x^{2}y^{2} + 4y\right )\)}{}{}}
\LANGes{
\Question{
Factor the following polynomial. 
\[ 
 27x^{6}z - 8y^{3}z
\]}
{\(z\left (3x^{2} - 2y\right )\left (9x^{4} + 6x^{2}y + 4y^{2}\right )\)}{\(z\left (3x^{2} + 2y\right )\left (9x^{4} + 6x^{2}y - 4y^{2}\right )\)}{\(z\left (3x^{2} + 2y\right )\left (9x^{4} - 6x^{2}y + 4y^{2}\right )\)}{\(z\left (3x^{2} - 2y\right )\left (9x^{4} + 6x^{2}y^{2} + 4y\right )\)}{}{}}
\End

\Data\ID{9000146209}
\Updated{2020-07-15 03:07:37}
\Level{C}
\SubArea{Polynomials and fractions}
\LANGen{
\Question{
Factor the following expression. 
\[ 
 27a^{3} - 8b^{9}
\]}
{\(\left (3a - 2b^{3}\right )\left (9a^{2} + 6ab^{3} + 4b^{6}\right )\)}{\(\left (3a + 2b^{3}\right )\left (9a^{2} - 6ab^{3} + 4b^{6}\right )\)}{\(\left (3a - 2b^{3}\right )\left (9a^{2} + 12ab^{3} + 4b^{6}\right )\)}{\(\left (3a + 2b^{3}\right )\left (9a^{2} - 12ab^{3} + 4b^{6}\right )\)}{}{}}
\LANGpl{
\Question{
Rozłóż na czynniki wyrażenie:
\[ 
 27a^{3} - 8b^{9}
\]}
{\(\left (3a - 2b^{3}\right )\left (9a^{2} + 6ab^{3} + 4b^{6}\right )\)}{\(\left (3a + 2b^{3}\right )\left (9a^{2} - 6ab^{3} + 4b^{6}\right )\)}{\(\left (3a - 2b^{3}\right )\left (9a^{2} + 12ab^{3} + 4b^{6}\right )\)}{\(\left (3a + 2b^{3}\right )\left (9a^{2} - 12ab^{3} + 4b^{6}\right )\)}{}{}}
\LANGcs{
\Question{
Rozložením výrazu \(27a^{3} - 8b^{9}\) 
na součin získáme výsledek:}
{\(\left (3a - 2b^{3}\right )\left (9a^{2} + 6ab^{3} + 4b^{6}\right )\)}{\(\left (3a + 2b^{3}\right )\left (9a^{2} - 6ab^{3} + 4b^{6}\right )\)}{\(\left (3a - 2b^{3}\right )\left (9a^{2} + 12ab^{3} + 4b^{6}\right )\)}{\(\left (3a + 2b^{3}\right )\left (9a^{2} - 12ab^{3} + 4b^{6}\right )\)}{}{}}
\LANGsk{
\Question{
Upravte na súčin.
\[ 
 27a^{3} - 8b^{9}
\]}
{\(\left (3a - 2b^{3}\right )\left (9a^{2} + 6ab^{3} + 4b^{6}\right )\)}{\(\left (3a + 2b^{3}\right )\left (9a^{2} - 6ab^{3} + 4b^{6}\right )\)}{\(\left (3a - 2b^{3}\right )\left (9a^{2} + 12ab^{3} + 4b^{6}\right )\)}{\(\left (3a + 2b^{3}\right )\left (9a^{2} - 12ab^{3} + 4b^{6}\right )\)}{}{}}
\LANGes{
\Question{
Factor the following expression. 
\[ 
 27a^{3} - 8b^{9}
\]}
{\(\left (3a - 2b^{3}\right )\left (9a^{2} + 6ab^{3} + 4b^{6}\right )\)}{\(\left (3a + 2b^{3}\right )\left (9a^{2} - 6ab^{3} + 4b^{6}\right )\)}{\(\left (3a - 2b^{3}\right )\left (9a^{2} + 12ab^{3} + 4b^{6}\right )\)}{\(\left (3a + 2b^{3}\right )\left (9a^{2} - 12ab^{3} + 4b^{6}\right )\)}{}{}}
\End

\Data\ID{9000146210}
\Updated{2020-07-15 03:07:37}
\Level{C}
\SubArea{Polynomials and fractions}
\LANGen{
\Question{
Factor the following expression. 
\[ 
 64x^{6} + 125
\]}
{\(\left (4x^{2} + 5\right )\left (16x^{4} - 20x^{2} + 25\right )\)}{\(\left (4x^{2} - 5\right )\left (16x^{4} + 20x^{2} + 25\right )\)}{\(\left (4x^{2} + 5\right )\left (16x^{3} - 20x^{2} + 25\right )\)}{\(\left (4x^{2} - 5\right )\left (16x^{3} + 20x^{2} + 25\right )\)}{}{}}
\LANGpl{
\Question{
Rozłóż na czynniki wyrażenie:
\[ 
 64x^{6} + 125
\]}
{\(\left (4x^{2} + 5\right )\left (16x^{4} - 20x^{2} + 25\right )\)}{\(\left (4x^{2} - 5\right )\left (16x^{4} + 20x^{2} + 25\right )\)}{\(\left (4x^{2} + 5\right )\left (16x^{3} - 20x^{2} + 25\right )\)}{\(\left (4x^{2} - 5\right )\left (16x^{3} + 20x^{2} + 25\right )\)}{}{}}
\LANGcs{
\Question{
Rozložením výrazu \(64x^{6} + 125\) 
na součin získáme výsledek:}
{\(\left (4x^{2} + 5\right )\left (16x^{4} - 20x^{2} + 25\right )\)}{\(\left (4x^{2} - 5\right )\left (16x^{4} + 20x^{2} + 25\right )\)}{\(\left (4x^{2} + 5\right )\left (16x^{3} - 20x^{2} + 25\right )\)}{\(\left (4x^{2} - 5\right )\left (16x^{3} + 20x^{2} + 25\right )\)}{}{}}
\LANGsk{
\Question{
Upravte na súčin.
\[ 
 64x^{6} + 125
\]}
{\(\left (4x^{2} + 5\right )\left (16x^{4} - 20x^{2} + 25\right )\)}{\(\left (4x^{2} - 5\right )\left (16x^{4} + 20x^{2} + 25\right )\)}{\(\left (4x^{2} + 5\right )\left (16x^{3} - 20x^{2} + 25\right )\)}{\(\left (4x^{2} - 5\right )\left (16x^{3} + 20x^{2} + 25\right )\)}{}{}}
\LANGes{
\Question{
Factor the following expression. 
\[ 
 64x^{6} + 125
\]}
{\(\left (4x^{2} + 5\right )\left (16x^{4} - 20x^{2} + 25\right )\)}{\(\left (4x^{2} - 5\right )\left (16x^{4} + 20x^{2} + 25\right )\)}{\(\left (4x^{2} + 5\right )\left (16x^{3} - 20x^{2} + 25\right )\)}{\(\left (4x^{2} - 5\right )\left (16x^{3} + 20x^{2} + 25\right )\)}{}{}}
\End

\Data\ID{9000146707}
\Updated{2020-07-15 03:07:37}
\Level{C}
\SubArea{Polynomials and fractions}
\LANGen{
\Question{
Divide the following two polynomials using long division. 
\[ 
 \left (6x^{2} - 5x - 6\right ) : (2x - 3)
\]}
{\(3x + 2\)}{\(3x - 2\)}{\(3x - 7 + \frac{15} 
{2x-3}\)}{\(3x - 7 - \frac{27} 
{2x-3}\)}{}{}}
\LANGpl{
\Question{
Podziel podane wielomiany:
\[ 
 \left (6x^{2} - 5x - 6\right ) : (2x - 3)
\]}
{\(3x + 2\)}{\(3x - 2\)}{\(3x - 7 + \frac{15} 
{2x-3}\)}{\(3x - 7 - \frac{27} 
{2x-3}\)}{}{}}
\LANGcs{
\Question{
Určete podíl mnohočlenů. \[\left (6x^{2} - 5x - 6\right ) : (2x - 3)\]}
{\(3x + 2\)}{\(3x - 2\)}{\(3x - 7 + \frac{15} 
{2x-3}\)}{\(3x - 7 - \frac{27} 
{2x-3}\)}{}{}}
\LANGsk{
\Question{
Určte podiel polynómov.
\[ 
 \left (6x^{2} - 5x - 6\right ) : (2x - 3)
\]}
{\(3x + 2\)}{\(3x - 2\)}{\(3x - 7 + \frac{15} 
{2x-3}\)}{\(3x - 7 - \frac{27} 
{2x-3}\)}{}{}}
\LANGes{
\Question{
Divide the following two polynomials using long division. 
\[ 
 \left (6x^{2} - 5x - 6\right ) : (2x - 3)
\]}
{\(3x + 2\)}{\(3x - 2\)}{\(3x - 7 + \frac{15} 
{2x-3}\)}{\(3x - 7 - \frac{27} 
{2x-3}\)}{}{}}
\End

\Data\ID{9000146708}
\Updated{2020-07-15 03:07:37}
\Level{C}
\SubArea{Polynomials and fractions}
\LANGen{
\Question{
Divide the following two polynomials using long division. 
\[ 
 \left (2x^{3} + x^{2} - 17x + 5\right ) : \left (x^{2} + 3x - 1\right )
\]}
{\(2x - 5\)}{\(2x + 5\)}{\(2x + 7 + \frac{2x-2} 
{x^{2}+3x-1}\)}{\(2x + 7 + \frac{12-40x} 
{x^{2}+3x-1}\)}{}{}}
\LANGpl{
\Question{
Podziel podane wielomiany:
\[ 
 \left (2x^{3} + x^{2} - 17x + 5\right ) : \left (x^{2} + 3x - 1\right )
\]}
{\(2x - 5\)}{\(2x + 5\)}{\(2x + 7 + \frac{2x-2} 
{x^{2}+3x-1}\)}{\(2x + 7 + \frac{12-40x} 
{x^{2}+3x-1}\)}{}{}}
\LANGcs{
\Question{
Určete podíl mnohočlenů. \[\left (2x^{3} + x^{2} - 17x + 5\right ) : \left (x^{2} + 3x - 1\right )\]}
{\(2x - 5\)}{\(2x + 5\)}{\(2x + 7 + \frac{2x-2} 
{x^{2}+3x-1}\)}{\(2x + 7 + \frac{12-40x} 
{x^{2}+3x-1}\)}{}{}}
\LANGsk{
\Question{
Určte podiel polynómov.
\[ 
 \left (2x^{3} + x^{2} - 17x + 5\right ) : \left (x^{2} + 3x - 1\right )
\]}
{\(2x - 5\)}{\(2x + 5\)}{\(2x + 7 + \frac{2x-2} 
{x^{2}+3x-1}\)}{\(2x + 7 + \frac{12-40x} 
{x^{2}+3x-1}\)}{}{}}
\LANGes{
\Question{
Divide the following two polynomials using long division. 
\[ 
 \left (2x^{3} + x^{2} - 17x + 5\right ) : \left (x^{2} + 3x - 1\right )
\]}
{\(2x - 5\)}{\(2x + 5\)}{\(2x + 7 + \frac{2x-2} 
{x^{2}+3x-1}\)}{\(2x + 7 + \frac{12-40x} 
{x^{2}+3x-1}\)}{}{}}
\End

\Data\ID{9000146709}
\Updated{2020-07-15 03:07:37}
\Level{C}
\SubArea{Polynomials and fractions}
\LANGen{
\Question{
Divide the following two polynomials using long division. 
\[ 
 \left (4x^{2} - 10x - 1\right ) : (x - 2)
\]}
{\(4x - 2 - \frac{5} 
{x-2}\)}{\(4x - 2 + \frac{3} 
{x-2}\)}{\(4x + 2 - \frac{5} 
{x-2}\)}{\(4x + 2 + \frac{3} 
{x-2}\)}{}{}}
\LANGpl{
\Question{
Podziel podane wielomiany:
\[ 
 \left (4x^{2} - 10x - 1\right ) : (x - 2)
\]}
{\(4x - 2 - \frac{5} 
{x-2}\)}{\(4x - 2 + \frac{3} 
{x-2}\)}{\(4x + 2 - \frac{5} 
{x-2}\)}{\(4x + 2 + \frac{3} 
{x-2}\)}{}{}}
\LANGcs{
\Question{
Určete podíl mnohočlenů. \[\left (4x^{2} - 10x - 1\right ) : (x - 2)\]}
{\(4x - 2 - \frac{5} 
{x-2}\)}{\(4x - 2 + \frac{3} 
{x-2}\)}{\(4x + 2 - \frac{5} 
{x-2}\)}{\(4x + 2 + \frac{3} 
{x-2}\)}{}{}}
\LANGsk{
\Question{
Určte podiel polynómov.
\[ 
 \left (4x^{2} - 10x - 1\right ) : (x - 2)
\]}
{\(4x - 2 - \frac{5} 
{x-2}\)}{\(4x - 2 + \frac{3} 
{x-2}\)}{\(4x + 2 - \frac{5} 
{x-2}\)}{\(4x + 2 + \frac{3} 
{x-2}\)}{}{}}
\LANGes{
\Question{
Divide the following two polynomials using long division. 
\[ 
 \left (4x^{2} - 10x - 1\right ) : (x - 2)
\]}
{\(4x - 2 - \frac{5} 
{x-2}\)}{\(4x - 2 + \frac{3} 
{x-2}\)}{\(4x + 2 - \frac{5} 
{x-2}\)}{\(4x + 2 + \frac{3} 
{x-2}\)}{}{}}
\End

\Data\ID{9000146710}
\Updated{2020-07-15 03:07:37}
\Level{C}
\SubArea{Polynomials and fractions}
\LANGen{
\Question{
Divide the following two polynomials using long division. 
\[ 
 \left (x^{3} + 3x^{2} - x + 4\right ) : \left (x^{2} - x + 1\right )
\]}
{\(x + 4 + \frac{2x} 
{x^{2}-x+1}\)}{\(x + 4 + \frac{2x+8} 
{x^{2}-x+1}\)}{\(x + 2 + \frac{6-2x} 
{x^{2}-x+1}\)}{\(x + 2 + \frac{2x+2} 
{x^{2}-x+1}\)}{}{}}
\LANGpl{
\Question{
Podziel podane wielomiany:
\[ 
 \left (x^{3} + 3x^{2} - x + 4\right ) : \left (x^{2} - x + 1\right )
\]}
{\(x + 4 + \frac{2x} 
{x^{2}-x+1}\)}{\(x + 4 + \frac{2x+8} 
{x^{2}-x+1}\)}{\(x + 2 + \frac{6-2x} 
{x^{2}-x+1}\)}{\(x + 2 + \frac{2x+2} 
{x^{2}-x+1}\)}{}{}}
\LANGcs{
\Question{
Určete podíl mnohočlenů. \[\left (x^{3} + 3x^{2} - x + 4\right ) : \left (x^{2} - x + 1\right )\]}
{\(x + 4 + \frac{2x} 
{x^{2}-x+1}\)}{\(x + 4 + \frac{2x+8} 
{x^{2}-x+1}\)}{\(x + 2 + \frac{6-2x} 
{x^{2}-x+1}\)}{\(x + 2 + \frac{2x+2} 
{x^{2}-x+1}\)}{}{}}
\LANGsk{
\Question{
Určte podiel polynómov.
\[ 
 \left (x^{3} + 3x^{2} - x + 4\right ) : \left (x^{2} - x + 1\right )
\]}
{\(x + 4 + \frac{2x} 
{x^{2}-x+1}\)}{\(x + 4 + \frac{2x+8} 
{x^{2}-x+1}\)}{\(x + 2 + \frac{6-2x} 
{x^{2}-x+1}\)}{\(x + 2 + \frac{2x+2} 
{x^{2}-x+1}\)}{}{}}
\LANGes{
\Question{
Divide the following two polynomials using long division. 
\[ 
 \left (x^{3} + 3x^{2} - x + 4\right ) : \left (x^{2} - x + 1\right )
\]}
{\(x + 4 + \frac{2x} 
{x^{2}-x+1}\)}{\(x + 4 + \frac{2x+8} 
{x^{2}-x+1}\)}{\(x + 2 + \frac{6-2x} 
{x^{2}-x+1}\)}{\(x + 2 + \frac{2x+2} 
{x^{2}-x+1}\)}{}{}}
\End
